\documentclass[11pt]{article}
\usepackage[T1]{fontenc}
\usepackage{lmodern,amsmath,amssymb,amsthm,booktabs,array}
\usepackage[letterpaper,margin=1in]{geometry}
\usepackage{microtype,xcolor}
\usepackage[colorlinks=true,linkcolor=blue!45!black,urlcolor=blue!45!black]{hyperref}
\newtheorem{theorem}{Theorem}[section]
\newtheorem{lemma}[theorem]{Lemma}
\newtheorem{proposition}[theorem]{Proposition}
\newtheorem{corollary}[theorem]{Corollary}
\newtheorem{remark}[theorem]{Remark}
\setlength{\emergencystretch}{3em}
\title{Further results from the cycle-space machinery\\Extended research report}
\author{Research working draft}
\date{10 November 2026}
\begin{document}
\maketitle

\begin{abstract}
We prove the sharp no-waste rank bound, replacing the additive constant three by one, and describe its equality cases.
Fixed-gap terminal-path stability is additive for arbitrary graphs: every terminal path has length at least $b-C_\varepsilon$, and the reduced terminal core has bounded pathwidth.
This combines Fable's localization and distance-layer argument with the extended forest estimate proved here.
Sharper explicit bounds hold for bounded genus and fixed proper minor-closed classes.
For sparse random hosts at the stated Alon--Krivelevich density, the second-order pancyclicity excess lies between coefficients $1$ and $199/192$ times $\ln n/\ln\ln n$.
An adaptive attachment construction proves the leading sparse-pancyclic bound throughout the full range $np-\ln n-\ln\ln n\to\infty$.
For every fixed proper minor-closed class, the cycle-count error above $2^{r-1}$ is $O(2^r/r)$.
A separate companion paper proves the exact $2^n+n$ cycle bound for four-regular multigraphs, classifies the extremizers, and identifies every graph above $2^n$.
The proofs have undergone independent internal audits and include reproducible finite checks.
They have not been formalised in Lean, and publication priority has not been established.
\end{abstract}

\section*{Scope and notation}
The source manuscript is \emph{The number of cycles in a graph and the minimum number of edges in a pancyclic graph}, dated November 10, 2026.
Its local source is \nolinkurl{paper/cycle-count/pancyclic-paper.tex} at repository revision \texttt{a965788}.
We use its definitions of outside parts and forest probability $p(G,W)$.
Rank means $\beta(G)=|E(G)|-|V(G)|+\operatorname{comp}(G)$.
An even subgraph is an edge set of even degree at every vertex.
A cycle is connected and two-regular on its incident vertices.

This report retains several intermediate estimates because their proofs supply ingredients for the stronger results that follow.
The initial 16-page report remains available as \nolinkurl{research-progress.pdf}.
The existing main manuscripts, formal Lean sources, and public site were not changed.
The shared repository-root handoff is \nolinkurl{RESEARCH-HANDOFF-CODEX-2026-11-10.txt}.

\tableofcontents

\section{The sharp no-waste rank bound}\label{sec:nowaste-sharp}

This section answers the final problem in the source manuscript.
The constant three can be replaced by one for simple graphs.
The proof uses the manuscript's elementary Erd\H{o}s--P\'osa lemma.
It does not use the asymptotic forest estimate.

\begin{theorem}\label{thm:nowaste-one}
Let $G$ be a connected simple graph, and let $W\subseteq G$ contain a cycle.
Suppose every outside part of an even subgraph is the outside part of a cycle.
Then
\[
 \beta(G)-\beta(W)\le |V(W)|+1.
\]
The additive constant one is sharp.
\end{theorem}

We first give the small classification needed in the proof.
A cyclic block means a maximal two-connected subgraph that contains a cycle.
Tree attachments do not change the cycle space.

\begin{lemma}\label{lem:nowaste-small-cores}
Suppose every nonempty even subgraph of a simple graph $F$ is a cycle.
Then $\beta(F)\le4$, and at most one component contains a cycle.
If $\beta(F)=2$, its unique cyclic block consists of three internally disjoint paths between two branch vertices.
If $\beta(F)=3$, that block is a subdivision of $K_4$.
If $\beta(F)=4$, that block is a subdivision of $K_{3,3}$.
All other edges in the cyclic component belong to trees attached to that block.
\end{lemma}

\begin{proof}
Two edge-disjoint cycles would have a union that is not a cycle.
Thus there is at most one cyclic block in the entire graph.
Every other edge is a bridge.
Remove the tree attachments and call the remaining cyclic block $B$.
It is two-connected and has minimum degree at least two.

No vertex of $B$ can have degree at least four.
Suppose $v$ did, and choose four edges incident to $v$.
Since $B-v$ is connected, it has an edge set whose odd-degree vertices are the four chosen neighbors.
Such an edge set exists by taking the appropriate join in a spanning tree.
Its union with the four chosen edges is even, with degree four at $v$.
This contradicts the hypothesis.
Hence the maximum degree of $B$ is three.

If every degree is two, then $B$ is a cycle and has rank one.
Otherwise suppress all degree-two vertices to obtain a connected cubic multigraph $J$.
Suppression preserves rank and the property concerning even subgraphs.
The manuscript's Erd\H{o}s--P\'osa lemma gives
$|E(J)|\le|V(J)|+3$.
As $J$ is cubic, it follows that $|V(J)|\le6$ and $\beta(J)\le4$.

For rank two, the cubic graph has two vertices.
It consists of three parallel edges, since the cyclic block has no bridge.
Its subdivision is the stated theta graph.
For rank three, the cubic graph has four vertices.
For rank four, it has six vertices.
In either case it has no loop or parallel pair.
Deleting such a cycle would leave at least as many edges as vertices, and therefore another edge-disjoint cycle.
The only simple cubic graph on four vertices is $K_4$.

In the six-vertex case, a triangle would also leave six edges after deletion, which again forces another cycle.
Thus $J$ is triangle-free.
Fix a vertex and its three neighbors.
The neighbors are independent and must each join both remaining vertices to have degree three.
This is exactly $K_{3,3}$.
\end{proof}

\begin{lemma}\label{lem:nowaste-path-obstruction}
Let $F=G-E(W)$, on the full vertex set of $G$.
Under the hypotheses of Theorem~\ref{thm:nowaste-one}, the following statements hold.
Every nonempty even subgraph of $F$ is a cycle.
Moreover, if $C$ is a cycle of $F$, then each component of $F-E(C)$ contains at most one vertex from each component of $W$.
\end{lemma}

\begin{proof}
An even edge set of $F$ is an outside part.
If a cycle of $G$ has this nonempty outside part, it contains that even edge set.
A cycle cannot contain a proper nonempty even edge set.
This proves the first assertion.

For the second assertion, suppose distinct $u,v$ lie in the same component of $W$ and the same component of $F-E(C)$.
Choose a $u$--$v$ path $P$ in $F-E(C)$ and a $u$--$v$ path $Q$ in $W$.
Their edge sets are disjoint.
The set $P\cup Q\cup C$ is even, since $P,Q$ have the same odd-degree vertices and $C$ is even.
Its outside part is $P\cup C$.
This contains the entire cycle $C$ and at least one further edge.
No cycle can have such an outside part, a contradiction.
\end{proof}

\begin{proof}[Proof of Theorem~\ref{thm:nowaste-one}]
Put $s=\operatorname{comp}(F)$, $c=\operatorname{comp}(W)$, $m=|V(W)|$, and $k=\beta(F)$.
Counting edges and vertices gives
\begin{equation}\label{eq:nowaste-rank-identity}
 \beta(G)-\beta(W)=m+k-s-c+1.
\end{equation}
It suffices to show $k\le s+c$.
Here $s,c\ge1$, so the assertion is immediate for $k\le2$.
By Lemmas~\ref{lem:nowaste-small-cores} and~\ref{lem:nowaste-path-obstruction}, only $k=3,4$ remain.

Choose a component $W_0$ of $W$ that contains a cycle.
It has at least three vertices, since $G$ is simple.
Each tree component of $F$ contains at most one vertex of $W_0$.
Indeed, the path between two such vertices would avoid any cycle in the cyclic component of $F$.

Suppose first that $k=3$.
The cyclic component of $F$ has a core that is a subdivision of $K_4$.
For each of the four triangles of $K_4$, delete the edges of its subdivided cycle.
One component of the remaining graph contains all four branch vertices.
Its core edges form the complementary spanning star.
Call these four components $T_1,\ldots,T_4$, including their surviving tree attachments.
Each $T_i$ contains at most one vertex of $W_0$, by Lemma~\ref{lem:nowaste-path-obstruction}.

Every vertex of the cyclic component belongs to at least two of these four components.
A branch vertex belongs to all four.
An internal vertex of a subdivided edge belongs to the two complementary stars that contain that edge.
A vertex in an attached tree has the same memberships as its attachment vertex.
Therefore the cyclic component contains at most two vertices of $W_0$.
The $s-1$ tree components each contain at most one further vertex.
Thus $3\le |V(W_0)|\le s+1$, and $s\ge2$.
It follows that $k=3\le s+c$.

Suppose next that $k=4$.
The core is a subdivision of $K_{3,3}$.
Delete the subdivided cycles corresponding to its nine four-cycles, one at a time.
Each deletion leaves a component containing all six branch vertices.
The complementary kernel edges form a spanning tree with two adjacent centers.
Each such component contains at most one vertex of $W_0$.

Every branch vertex belongs to all nine components.
Each kernel edge belongs to five complementary trees, since it lies in four of the nine four-cycles.
Thus an internal vertex of a subdivided edge belongs to five components.
Tree attachments again inherit the membership of their attachment vertex.
Every vertex of the cyclic component therefore belongs to at least five of the nine components.
The cyclic component contains at most $\lfloor9/5\rfloor=1$ vertex of $W_0$.
Consequently $3\le |V(W_0)|\le s$, and $s\ge3$.
Hence $k=4\le s+c$.
Equation~\ref{eq:nowaste-rank-identity} proves the bound.

For sharpness, take $G=K_5$ minus one edge.
Let $W$ be the triangle on its three vertices of degree four.
Here $F=K_{2,3}$, with the three marked vertices in its larger part.
Its outside parts are the empty set, its three four-cycles, and four paths for each of the three marked vertex pairs.
Each path closes with the corresponding edge of $W$.
The empty outside part is represented by $W$ itself.
Here $\beta(G)=5$, $\beta(W)=1$, and $|V(W)|=3$.
Thus the difference is $4=|V(W)|+1$.
\end{proof}

The simple-graph assumption supplies three vertices in $W_0$.
The theorem is not asserted for marked loops or marked parallel pairs.
The new proof has not been formalized in Lean.


\begin{corollary}[Dependence on the marked components]\label{cor:nowaste-component}
Under the same hypotheses, put $m=|V(W)|$, $c=\operatorname{comp}(W)$, and
$M=\max_i |V(W_i)|$, where $W_i$ ranges over the components of $W$.
Then
\[
 \beta(G)-\beta(W)
 \le m-c+1+\max\{0,4-M\}.
\]
In particular, the difference is at most $m$ unless $W$ is a triangle.
If $W$ is connected, the difference is at most $\max\{m,4\}$.
\end{corollary}

\begin{proof}
Keep $F,s,k$ from the preceding proof.
If $k\le1$, then $k-s\le0$, and the rank identity gives the result.
Suppose $k\ge2$.
For ranks three and four, the preceding covering arguments apply to every component $W_i$.
They bound its intersection with the cyclic component of $F$ by two and one, respectively.

For rank two, use the three cycles of the theta core.
The three complementary components contain its two branch vertices and one subdivided branch each.
Each such component contains at most one vertex of $W_i$.
Every vertex of the cyclic component belongs to at least one of them.
Thus at most three vertices of $W_i$ lie in that component.

In each of the three cases, at most $5-k$ vertices of $W_i$ lie in the cyclic component.
Each of the other $s-1$ components contributes at most one.
Consequently $|V(W_i)|\le s+4-k$, and $k-s\le4-M$.
Substitute in the rank identity.
Since $M\ge3$, an excess of one over $m$ requires $c=1$ and $M=3$.
A connected simple graph with three vertices and a cycle is a triangle.
\end{proof}

\paragraph{Equality structure.}
The equality case in Theorem~\ref{thm:nowaste-one} has an explicit description.
For a vertex in the cyclic component of $F$, call its nearest vertex in the cyclic block its \emph{anchor}.
A vertex already in the block anchors itself.
Every tree attached to the block has exactly one attachment vertex, so the anchor is well-defined.

\begin{proposition}\label{prop:nowaste-equality}
Under the no-waste hypothesis,
$\beta(G)-\beta(W)=|V(W)|+1$ if and only if $W$ is a triangle and one of the following configurations holds.
\begin{enumerate}
\item $F$ is connected of rank two.
Its cyclic block is a theta graph.
The three marked vertices have anchors in the interiors of its three different branches.
\item $F$ has two components and rank three.
Its cyclic block is a subdivision of $K_4$.
Two marked vertices lie in this component, with anchors interior to two opposite subdivided edges.
The third marked vertex lies in the tree component.
\item $F$ has three components and rank four.
Its cyclic block is a subdivision of $K_{3,3}$.
Each component of $F$ contains exactly one marked vertex.
The two other components are trees.
\end{enumerate}
Conversely, every connected simple graph with $W$ a triangle and one of these configurations satisfies the no-waste hypothesis and equality.
\end{proposition}

\begin{proof}
Corollary~\ref{cor:nowaste-component} forces $W$ to be a triangle.
The rank identity then gives $k=s+1$.
The three possibilities are $(k,s)=(2,1),(3,2),(4,3)$.

For the theta core, the three marked vertices must use the three complementary-component slots exactly once.
An anchor at a branch vertex would use all three slots by itself.
Hence the anchors are internal and lie on different branches.
For the $K_4$ core, the two marked vertices in the cyclic component must use the four star slots exactly twice each.
Thus their anchors are internal and their kernel edges have disjoint endpoints.
These edges are opposite.
For the $K_{3,3}$ core, the component bound gives at most one marked vertex in each component.
All three components therefore contain one.

For the converse, an outside part is an edge set in $F$ whose odd-degree vertices lie among the three marked vertices.
Its boundary is empty or a marked pair.
Every allowable edge set can be completed inside the triangle.
The empty-boundary sets are the even subgraphs of $F$.
Their nonempty members are cycles by the core descriptions.

In the theta configuration, each marked pair has exactly four simple paths in $F$.
These exhaust its $2^2$ joins.
In the $K_4$ configuration, only the pair in the cyclic component can be a boundary.
The eight paths between opposite-edge interior points exhaust its $2^3$ joins.
Tree tails do not change these counts or create repeated vertices.
In the $K_{3,3}$ configuration, component parity excludes every nonempty boundary.
Every nonempty outside part is therefore either a cycle of $F$ or a marked-pair path.
Close each path with the corresponding triangle edge, and use $W$ for the empty part.
This proves no waste.
The rank identity gives equality in each configuration.
\end{proof}

\paragraph{Finite checks.}
The script \texttt{check\_no\_waste.py} checks every marked edge subset containing a cycle in every connected simple graph through six vertices.
It uses the NetworkX graph atlas and exact cycle-space enumeration.
The refined component bound passes for all $101{,}904$ tested pairs, including $19{,}035$ pairs with no waste.
Every equality case has a marked triangle.
These checks support the proof but do not replace it.

\section{Path stability and the size of a marked subgraph}
\label{sec:research-path}

These results use the forest proof in the November 10 manuscript.
The parameter extension below is a mathematical argument, not a Lean formalization.
It first improves the general path bound to order $b/\log b$.
Theorem~\ref{thm:general-linear-path-stability} later obtains the full linear bound from this enlarged forest range.
Theorem~\ref{thm:general-additive-path-stability} then strengthens it to an additive loss depending only on the density gap.
For two-terminal series-parallel networks, an elementary argument gives an explicit logarithmic loss.

All logarithms in this section have base two, except for $\ln$ and $\exp$.
Write $P_{uv}(N)$ for the number of simple paths between distinct vertices $u,v$.
Write $b=\beta(N)$.

\subsection{From paths to forest outside parts}

\begin{lemma}\label{lem:research-path-injection}
Let $N$ be a connected simple graph, and let $A$ be any $u$--$v$ path of length $\ell$.
Add a new vertex $w$ and edges $uw,wv$ to obtain $H$.
Let $W=A\cup\{uw,wv\}$.
Then
\[
 \beta(H)=b+1,\qquad \beta(W)=1,\qquad |E(W)|=\ell+2,
 \qquad P_{uv}(N)\le p(H,W)2^b.
\]
\end{lemma}

\begin{proof}
Each terminal path $P$ gives the cycle $P\cup\{uw,wv\}$ of $H$.
Its outside part is $P\setminus E(A)$, which is a forest.
Suppose that $P\setminus E(A)=Q\setminus E(A)$.
Then $P\mathbin\triangle Q$ is an even edge set supported on the forest $A$.
It must be empty, so $P=Q$.
There are $2^{\beta(H)-\beta(W)}=2^b$ outside parts.
Their forest proportion is $p(H,W)$.
\end{proof}

The formalized forest theorem gives $\ell>\sqrt{b+1}-2$ for fixed positive path-density gap and sufficiently large $b$.
The manuscript's size extension gives $\ell>(b+1)/\log^2(b+1)-2$.
The next argument removes one logarithmic factor for a fixed gap.

\subsection{A fixed-gap forest estimate at size \texorpdfstring{$r/\log r$}{r/log r}}

\begin{theorem}\label{thm:research-marked-size}
Fix $0<\delta<1/2$, and set
\[
 A_\delta=6\exp(16/\delta),\qquad
 a_\delta=\frac{1}{12A_\delta}.
\]
For sufficiently large $r$, the following statement holds.
Let $G$ be a connected simple graph of rank $r$.
Let $W$ be a nonempty union of cycles satisfying
\[
 |E(W)|\le a_\delta\frac{r}{\log r},\qquad
 \operatorname{comp}(W)\le\frac{\log r}{100}.
\]
Then $p(G,W)\le1/2+\delta$.
\end{theorem}

\begin{proof}
We use the manuscript's reduction, cut lemmas, spectral estimate, and second-moment argument.
The details below check each estimate affected by the new parameters.
Constants in $O_\delta(\cdot)$ may depend on $\delta$.

Put $x=\log r$, $y=\log x$, $D=\lfloor x/10\rfloor$, and $s=\lfloor(D-1)/2\rfloor$.
Set $L=\lfloor A_\delta x\rfloor$.
The reduction gives a subcubic multigraph $\Gamma$ and a marked vertex set $P$.
Every vertex outside $P$ has degree three, $|P|\le2|E(W)|$, and $\Gamma[P]$ has at most $x/100$ components.
Also $\beta(\Gamma)=r$ and $2r-2\le |V(\Gamma)|\le2r+2|E(W)|$.
Let $\mathcal F$ mean that the random even subgraph is a forest on $\Gamma-P$.
It suffices to rule out $\mathbb P(\mathcal F)>1/2+\delta$.

Assume this strict inequality.
Choose a maximal vertex-disjoint family $\mathcal Q$ of cycles of length at most $D$ in $\Gamma-P$.
The small-cut lemma gives
\[
 |\mathcal Q|<\frac{(D+1)2^{D-1}}{\delta}
       =O_\delta(xr^{1/10}).
\]
For their vertex union $S$, this implies
\[
 |S|=O_\delta(x^2r^{1/10})=o(r/x).
\]
Let $J=\Gamma-(P\cup S)$, and put $v=|V(J)|$.
This graph is simple and has girth greater than $D$.
The vertex and deficit estimates become
\[
 v=2r+O_\delta(r/x),\qquad
 d_J:=\sum_{z\in V(J)}(3-\deg_Jz)
 \le 3(|P|+|S|)
 \le\frac{6a_\delta r}{x}+o(r/x).
\]
Consequently,
\[
 \frac{Ld_J}{2v}\le\frac32 A_\delta a_\delta+o(1)
                 =\frac18+o(1).
\]

Let $B$ be the non-backtracking matrix of $J$.
For even $j$ with $6x\le j\le L$, the spectral estimate gives
\[
 2^{-j}\operatorname{tr}B^j
 \ge\left(1-\frac{d_J}{2v}\right)^j-3v2^{-j/2}
 \ge1-\frac{Ld_J}{2v}-3vr^{-3}\ge\frac12
\]
for sufficiently large $r$.
For the cycles $C$ of $J$ of length at most $L$, put
$\mu=\sum_C2^{-|C|}$.
The repeated-vertex walk estimate from the manuscript has error at most
\[
 \frac32 L^2 2^{-s}\mu=o(\mu).
\]
Therefore
\[
 (1+o(1))\mu
 \ge\frac12\sum_{\substack{6x\le j\le L\\j\ \mathrm{even}}}\frac1{2j}
 =\frac18\ln(A_\delta/6)+o(1).
\]
In particular, $\mu\ge\frac1{16}\ln(A_\delta/6)=1/\delta$ for sufficiently large $r$.

It remains to check the pair correlations.
Use $t_0=\lfloor y/2\rfloor$ in the manuscript's link argument.
For disjoint cycles $C,C'$ of length at most $L$, it gives
\[
 \ell(C,C')\le \operatorname{comp}(\Gamma[P])
 +\frac{2L}{t_0}
 +\frac{(t_0+1)2^{t_0-1}}{\delta}
 +\left(\frac{8L}{D}\right)^2
 \le\frac{x}{100}+o(x).
\]
Here the three added terms are respectively $O_\delta(x/y)$, $O_\delta(y\sqrt{x})$, and $O_\delta(1)$.
The exceptional-partner estimate thus remains
\[
 2^{x/100+o(x)}L^3 2^{-s}
 \le2^{x/100-x/20+o(x)}=o(1).
\]
The same weighted variable $Z$ therefore satisfies
\[
 \mathbb EZ=\mu,\qquad
 \mathbb EZ^2\le2\mu^2+(1+o(1))\mu.
\]
Since $\mathcal F$ implies $Z=0$, the second-moment lemma yields
\[
 \mathbb P(\mathcal F)\le\frac12+\frac{1+o(1)}{4\mu}
 \le\frac12+\frac\delta2
\]
for sufficiently large $r$.
This contradicts the assumption.
\end{proof}

\begin{corollary}\label{cor:research-general-path}
Fix $0<\varepsilon\le1/2$.
For sufficiently large $b$, if
\[
 P_{uv}(N)\ge(1/2+\varepsilon)2^b,
\]
then every terminal path has length
\[
 \ell>a_{\varepsilon/2}\frac{b+1}{\log(b+1)}-2.
\]
In particular, $\ell=\Omega_\varepsilon(b/\log b)$.
\end{corollary}

\begin{proof}
Apply Lemma~\ref{lem:research-path-injection}.
If the displayed length inequality failed, Theorem~\ref{thm:research-marked-size} would give
$P_{uv}(N)\le(1/2+\varepsilon/2)2^b$, a contradiction.
The one-component condition holds for sufficiently large $b$.
\end{proof}

The threshold is not explicit because the supporting asymptotic estimates are not explicit.
The additive $-2$ comes from the added two-edge path.
Neither issue affects the asymptotic order.

\begin{corollary}\label{cor:research-small-marked}
Suppose $r\to\infty$ along connected simple graphs $G$ and nonempty unions of cycles $W$.
If $|E(W)|=o(r/\log r)$ and $\operatorname{comp}(W)\le(\log r)/100$, then
$p(G,W)\le1/2+o(1)$.
\end{corollary}

\begin{proof}
For each fixed $0<\delta<1/2$, the size hypothesis eventually meets Theorem~\ref{thm:research-marked-size}.
Thus $\limsup p(G,W)\le1/2$.
\end{proof}

\subsection{An additive theorem for two-terminal series-parallel networks}

A two-terminal series-parallel network starts with one edge and its two endpoints.
Series composition identifies one terminal from each network.
Parallel composition identifies the corresponding terminal pairs.
Apart from those identifications, the two networks are disjoint.
Parallel edges are permitted and are counted separately.
This definition imposes a condition on the terminals, not only on the underlying graph.

\begin{theorem}\label{thm:research-ttsp}
Let $(N,u,v)$ be a two-terminal series-parallel network of rank $b$ and terminal distance $\ell$.
Then
\[
 P_{uv}(N)\le 2^{b-1}+2^{\ell-1}.
\]
Consequently, for $0<\varepsilon\le1/2$,
\[
 P_{uv}(N)\ge(1/2+\varepsilon)2^b
 \quad\Longrightarrow\quad
 \ell\ge b-\log\frac1{2\varepsilon}
 \quad\Longrightarrow\quad
 \ell\ge\frac{b}{\log(1/\varepsilon)}.
\]
\end{theorem}

\begin{proof}
Every terminal path is a join, so $P_{uv}(N)\le2^b$ for every connected network.
Define $h(t)=(1+2^{-t})/2$ for $t\ge0$.
We prove the stronger normalized statement
\[
 \frac{P_{uv}(N)}{2^b}\le h((b-\ell)_+),
 \qquad (z)_+=\max\{z,0\}.
\]
It holds for a single edge.

For series composition, write $b_i,\ell_i,P_i$ for the two input networks.
Then $b=b_1+b_2$, $\ell=\ell_1+\ell_2$, and $P=P_1P_2$.
For $s,t\ge0$,
\[
 h(s+t)-h(s)h(t)=\frac{(1-2^{-s})(1-2^{-t})}{4}\ge0.
\]
Also $(b_1-\ell_1)_++(b_2-\ell_2)_+\ge(b-\ell)_+$, and $h$ is decreasing.
The normalized statement follows by multiplication.

For parallel composition, $b=b_1+b_2+1$, $\ell=\min\{\ell_1,\ell_2\}\ge1$, and $P=P_1+P_2$.
Since $b_1,b_2\ge0$,
\[
 P\le2^{b_1}+2^{b_2}
 \le2^{b_1+b_2}+1
 =2^{b-1}+1
 \le2^{b-1}+2^{\ell-1}.
\]
Together with $P\le2^b$, this is the normalized statement.
Induction proves the theorem's first inequality.

The density hypothesis now gives $\varepsilon2^b\le2^{\ell-1}$.
Thus $b\le\ell+\log(1/(2\varepsilon))$.
Since $\ell\ge1$, the right side is at most $\ell\log(1/\varepsilon)$.
\end{proof}

The coefficient is sharp when $\varepsilon=2^{-b}$.
Take a series of $b-1$ triangles between $u,v$, and add the edge $uv$.
For $b\ge3$, the result is simple, has rank $b$, distance one, and $2^{b-1}+1$ terminal paths.
For $b=2$, a direct edge in parallel with two internally disjoint two-edge paths has the same count.
For $b=1$, use a triangle with adjacent terminals.

The theorem fails for arbitrary choices of terminals in a series-parallel graph.
On vertices $0,1,2,3,4$, take
\[
 E=\{01,02,03,04,12,14,23\},\qquad u=3,\quad v=4.
\]
This graph has rank three and terminal distance two.
Its seven terminal paths are
\[
 3014,\quad30214,\quad304,\quad32014,\quad3204,\quad32104,\quad3214.
\]
Thus $P=7>2^{3-1}+2^{2-1}=6$.
The graph is a triangle with two vertices successively attached to existing edges, so it has no $K_4$ minor.
It does not satisfy the two-terminal composition hypothesis for terminals $3,4$.

\subsection{An elementary structural reduction}

\begin{proposition}\label{prop:research-offroute}
Let $N$ be connected, with distinct terminals $u,v$.
Let $N_0$ be the union of the blocks on the $u$--$v$ route in the block-cut tree.
Count bridges as blocks.
Put $s=\beta(N)-\beta(N_0)$.
Then
\[
 P_{uv}(N)=P_{uv}(N_0)\le2^{\beta(N_0)}=2^{b-s}.
\]
In particular, $P_{uv}(N)>2^{b-1}$ implies $s=0$.
\end{proposition}

\begin{proof}
A simple terminal path cannot enter a block outside the terminal route.
It would have to revisit the articulation vertex where it left that route.
Thus all terminal paths lie in $N_0$.
Apply the join bound to $N_0$.
Ranks add across blocks, so $s$ is a nonnegative integer.
\end{proof}

This removes cyclic blocks outside the terminal route from the unresolved problem.
It does not control the internal structure of a single block.
The fixed-gap argument still loses a factor $\log b$ there.
Its spectral estimate counts walks of length at least order $\log b$.
A marked set of order $b$ creates a constant degree deficit, whose effect accumulates along those walks.
The present retuning does not prove general linear stability.

\paragraph{Finite checks.}
The accompanying script checks 1,728 distinct composition states for networks with at most 12 edges.
It verifies the inequality with exact integers across 27,264 series or parallel transitions.
It also enumerates the seven counterexample paths and the sharp simple examples for ranks three through nine.
These checks do not verify the asymptotic estimates or replace their proofs.

\section{Linear path stability for planar graphs and bounded genus}
\label{sec:path-phase-two}

This section gives quantitative bounds for planar graphs, using the small-cut lemma and a reduction adapted to one path.
The same argument gives a bound with an additive genus term.
A separator argument also proves linear stability for every fixed proper minor-closed class.
These results are not formalized in Lean.

Write $P_{uv}(N)$ for the number of simple paths between distinct vertices $u,v$.
Write $b=\beta(N)$.
The orientable genus of a planar graph is zero.

\begin{theorem}[Linear stability with a genus term]
\label{thm:path-genus}
Let $N$ be a connected simple graph that embeds in an orientable surface of genus at most $g$.
Let $A$ be a $u$--$v$ path of length $\ell$.
Suppose that
\[
 P_{uv}(N)\ge(1/2+\varepsilon)2^b,
 \qquad 0<\varepsilon\le1/2.
\]
For every integer $L\ge6$,
\begin{equation}\label{eq:path-genus-explicit}
 (L-5)b
 \le 2(L+3)\ell+2(L+1)g+(L+1)
 +\frac{3L(L+1)(L-1)2^{L-2}}{\varepsilon}.
\end{equation}
In particular,
\[
 b\le10\ell+8g+4+\frac{16128}{\varepsilon}.
\]
For every fixed $\eta>0$, a constant $C_\eta$ gives the stronger asymptotic form
\[
 b\le(2+\eta)(\ell+g)+\frac{C_\eta}{\varepsilon}.
\]
\end{theorem}

\begin{corollary}[The planar case]
\label{cor:path-planar}
Under the same density hypothesis, if $N$ is planar, then
\[
 \ell\ge\frac{b-4-16128/\varepsilon}{10}.
\]
For all ranks, the positive linear bound
\[
 \ell\ge\frac{\varepsilon}{14\varepsilon+16128}\,b
\]
also holds.
For each fixed $\varepsilon>0$ and each $\eta>0$, one has
$\ell\ge b/(2+\eta)-O_{\eta,\varepsilon}(1)$.
\end{corollary}

\begin{proof}[Derivation of the corollary]
Set $g=0$ and $L=7$ in Theorem~\ref{thm:path-genus}.
Since $\ell\ge1$, its upper bound for $b$ is at most
$(14+16128/\varepsilon)\ell$.
The last assertion uses the theorem's final inequality.
\end{proof}

\subsection{A reduction that preserves a short marked path}

\begin{lemma}\label{lem:path-planar-split}
Let $N,A,u,v$ be as above, without a path-density assumption.
There is a connected multigraph $\Gamma$ on the same surface and a terminal path $A'$ such that
\begin{enumerate}
\item $\beta(\Gamma)=b$ and $P_{uv}(N)\le P_{u'v'}(\Gamma)$;
\item every vertex outside $A'$ has degree three;
\item every internal vertex of $A'$ has degree two or three, and its endpoints have degree one or two;
\item $a:=|V(A')|\le2\ell$ and $|V(\Gamma)|\le2b+a$.
\end{enumerate}
\end{lemma}

\begin{proof}
First delete leaves other than the terminals, until no such leaves remain.
This preserves the rank and the terminal paths.
It preserves $A$ as well.

Replace each vertex by a tree inside a small disk in the fixed surface embedding.
The external edges retain their cyclic order around that disk.
At a vertex outside $A$, use a tree whose vertices have degree three, counting external edges.
A degree-two vertex can instead remain unsplit at this stage.

Consider an internal vertex of $A$.
Its two path edges divide the other incident edges into two intervals in the cyclic order.
If both intervals are nonempty, use two marked vertices joined by one edge.
Attach one path edge to each marked vertex.
Attach the two intervals through two rooted trees, one at each marked vertex.
Each marked vertex has degree three.
If only one interval is nonempty, one marked vertex suffices, with both path edges and one rooted tree.
If both intervals are empty, retain the degree-two vertex.
These trees lie in the disk and respect its boundary order.
Their unmarked vertices have degree three.

At each terminal, use one marked vertex.
Its incident path edge is one branch.
Attach all its remaining incident edges through one rooted tree in the other branch.
This marked vertex has degree two, or degree one when there are no remaining edges.
All unmarked vertices of the rooted tree have degree three.
For a rooted tree with only one external edge, use that edge directly.

Replacing a vertex by a tree does not change the rank.
Each original terminal path has a unique lift through these trees.
That lift is simple, since the original path visits each original vertex at most once.
Distinct original paths have distinct lifts.
There can be additional paths after splitting, so the required comparison is an inequality.

The lifted path uses at most two vertices for each original internal path vertex and one for each endpoint.
Thus $a\le2(\ell-1)+2=2\ell$.
Suppress every degree-two vertex outside the lifted path.
This preserves the rank, the embedding, and the path count.
It gives the required multigraph $\Gamma$ and path $A'$.

Write $n=|V(\Gamma)|$.
Every vertex outside $A'$ has degree three.
The internal marked vertices have degree at least two, and the two endpoints have degree at least one.
Therefore
\[
 2(n+b-1)=2|E(\Gamma)|\ge3(n-a)+2(a-2)+2=3n-a-2.
\]
This gives $n\le2b+a$.
\end{proof}

The use of two marked vertices is important.
For arbitrary cyclic orders, grouping both path edges at one marked cubic vertex need not preserve the embedding.

\subsection{Short faces supply disjoint cyclic regions}

\begin{lemma}\label{lem:path-face-packing}
Let $\Gamma,A',a,b,g$ satisfy Lemma~\ref{lem:path-planar-split}.
For each integer $L\ge6$, there are $k$ pairwise vertex-disjoint cycles outside $V(A')$, each of length at most $L$, with
\begin{equation}\label{eq:path-packing}
 k\ge
 \frac{(L-5)b-(L+3)a-2(L+1)g-(L+1)}{3L(L+1)}.
\end{equation}
A negative right side imposes no requirement.
Each chosen vertex set is a cyclic region whose edge boundary in $\Gamma$ has size at most $L$.
\end{lemma}

\begin{proof}
Let $J=\Gamma-V(A')$.
The marked vertices have degree at most three and contain the $a-1$ edges of $A'$.
Thus at most $3a-(a-1)=2a+1$ edges disappear.
If $J$ is nonempty, this gives
\[
 \beta(J)\ge b-a-1.
\]
The inequality also holds when $J$ is empty.
In that case subcubicity gives $b\le a/2+1\le a+1$.
The total number of face-boundary edge occurrences is at most
\[
 2|E(J)|\le2|E(\Gamma)|
 =2(|V(\Gamma)|+b-1)\le6b+2a-2.
\]

In the planar case there are exactly $\beta(J)$ bounded faces.
For a surface of genus $g$, take a regular neighborhood of each component and cap its boundary curves.
This gives cellular embeddings with total genus at most $g$.
Discard one face from each component.
Euler's formula leaves at least $\beta(J)-2g$ faces.
Each remaining face has a boundary containing a cycle.
Tree components contribute no remaining faces.
These statements permit repeated boundary edges and parallel edges.

Call a retained face short when its total boundary length is at most $L$.
Every other retained face has at least $L+1$ boundary occurrences.
The number $f_L$ of short faces therefore satisfies
\[
 \begin{split}
 f_L&\ge b-a-1-2g-\frac{6b+2a}{L+1}\\
 &=\frac{(L-5)b-(L+3)a-2(L+1)g-(L+1)}{L+1}.
 \end{split}
\]
If this bound is positive, choose a cycle in a short face boundary.
Delete from the list every short face incident to a vertex of that cycle.
Repeat until the list is empty.

Every chosen cycle has at most $L$ vertices.
Each such vertex is incident to at most three faces, since its degree in $J$ is at most three.
Thus each choice removes at most $3L$ faces from the list.
The selected cycles are vertex-disjoint, so $k\ge f_L/(3L)$.
This proves \eqref{eq:path-packing}.

Every selected vertex has degree three in $\Gamma$ and uses two incidences in its cycle.
Its cyclic region therefore has at most $|C|\le L$ boundary edges.
\end{proof}

\subsection{The probability estimate}

We use a slightly sharper constant in the manuscript's small-cut lemma.
For disjoint proper cyclic regions with boundary at most $L$, the cut formula gives $p_i\ge2^{-L}$.
The exceptional-partner proof bounds the partners with at least three links by $L-2$.
Consequently, for $Z=\sum_i\mathbf1_{E_i}$ and $\mu=\mathbb EZ$,
\[
 \mu\ge k2^{-L},\qquad
 \mathbb EZ^2\le2\mu^2+(L-1)\mu.
\]
The same Cauchy--Schwarz argument gives
\begin{equation}\label{eq:path-small-cuts}
 \Pr\bigl(\text{the even subgraph is a forest in every region}\bigr)
 \le\frac12+\frac{(L-1)2^{L-2}}{k}.
\end{equation}
Here $k\ge1$, and every region in our application is proper because it avoids $A'$.

\begin{proof}[Proof of Theorem~\ref{thm:path-genus}]
Apply Lemma~\ref{lem:path-planar-split}.
Add a new two-edge terminal path to $\Gamma$, and let $H$ be the resulting graph.
Let $W$ be its union with $A'$.
The path injection from the first research note gives
\[
 \frac12+\varepsilon
 \le\frac{P_{uv}(N)}{2^b}
 \le\frac{P_{u'v'}(\Gamma)}{2^b}
 \le p(H,W).
\]
The genus of $H$ is irrelevant.
The regions supplied by Lemma~\ref{lem:path-face-packing} lie in the embedded graph $\Gamma-V(A')$.
Their cuts are unchanged when the new path is added.
On the forest event outside $W$, the even subgraph is a forest inside each such region.
Thus \eqref{eq:path-small-cuts} yields
\[
 k\le\frac{(L-1)2^{L-2}}{\varepsilon}.
\]
Combine this with \eqref{eq:path-packing} and $a\le2\ell$.
If the right side of \eqref{eq:path-packing} is nonpositive, the desired inequality already holds.
Otherwise the displayed upper bound on $k$ proves \eqref{eq:path-genus-explicit}.

For $L=7$, division by $L-5=2$ gives the stated numerical bound.
For the last assertion, choose $L$ so large that $16/(L-5)\le\eta$.
Then
\[
 \frac{2(L+3)}{L-5}\le2+\eta,
 \qquad
 \frac{2(L+1)}{L-5}\le2+\eta.
\]
Absorb the remaining constant into $C_\eta/\varepsilon$.
\end{proof}

\subsection{Restrictions on a possible general proof}

If the path family has more than half of its affine join space, every even edge set is the difference of two paths.
Indeed, a set and each of its translates must intersect when both have size greater than half the ambient space.
This is necessary, but it does not supply a quantitative length bound.
The chain-with-a-shortcut examples satisfy this condition at arbitrarily large rank.

There is also a useful local restriction.
In a two-connected block, the restriction of a uniform even subgraph to the edges at a vertex is uniform over all even patterns.
Any such pattern extends by a join in the connected graph obtained after deleting that vertex.
For a terminal join, the terminal patterns are instead uniform over all odd patterns.
Thus path density above one half forces internal block degrees at most five and terminal block degrees at most three.
The relevant local probabilities are
\[
 \frac{1+\binom d2}{2^{d-1}}\quad\text{and}\quad
 \frac{d}{2^{d-1}},
\]
which are at most one half for $d\ge6$ and $d\ge4$, respectively.
Ranks and normalized path counts factor over the blocks on the terminal route.
Therefore this observation applies separately to each such block.
It still permits large cubic blocks.

The tempting general bound $P_{uv}\le2^{b-1}+2^\ell-1$ is false.
Start with $K_{3,3}$ on parts $\{0,1,2\}$ and $\{3,4,5\}$, and delete the edge $03$.
Add a vertex $6$ and edges $63,64$.
For terminals $0,6$, the resulting simple graph has rank four and distance two.
It has eight paths ending through $36$ and five ending through $46$, for thirteen paths in total.
The proposed upper bound would be $8+4-1=11$.
The following family rules out that broader form as well.

\begin{proposition}[Unbounded excess at distance two]
\label{prop:path-distance-two-obstruction}
For each integer $b\ge3$, there is a connected simple planar graph of rank $b$, terminal distance two, and
\[
 P_{uv}=2^{b-1}+b.
\]
Consequently, there is no bound $P_{uv}\le2^{b-1}+f(\ell)$ for all connected simple graphs, for any finite-valued function $f$.
In particular, $2^{C\ell}$ cannot serve as an additive error for any fixed $C$.
\end{proposition}

\begin{proof}
Take the ladder on columns $(x_i,y_i)$, $0\le i\le b-1$.
Its edges are the two rails $x_ix_{i+1},y_iy_{i+1}$ and all rungs $x_iy_i$.
It has rank $b-1$.
Each $x_0$--$x_{b-1}$ path crosses every two-edge cut between consecutive columns exactly once.
The choices of upper or lower rail at the $b-1$ cuts determine the path uniquely.
All $2^{b-1}$ choices give simple paths.

Add the edge $x_0x_{b-1}$, and call the resulting planar graph $G$.
Its rank is $b$.
The old cycles are rectangles between two columns.
Exactly $b-1$ of these contain $x_0$.
Every new cycle uses the added edge and corresponds to one of the $2^{b-1}$ terminal paths.
Hence
\[
 c_{x_0}(G)=2^{b-1}+b-1.
\]

The vertex $x_0$ has degree three.
Replace it by three vertices $p_0,p_1,p_2$ joined as a path.
Assign one of its three former incident edges to each new vertex.
This local tree replacement preserves planarity and rank.
Use $p_0,p_2$ as the terminals.
Their distance is two.

The path $p_0p_1p_2$ is one terminal path.
Every other terminal path has exactly one segment outside these three vertices.
Its endpoints are two of the three ports.
Contracting the replacement tree turns that segment into a cycle through $x_0$.
Conversely, every such cycle gives one terminal path, using zero or one of the replacement edges as needed.
Thus the terminal path count is $1+c_{x_0}(G)=2^{b-1}+b$.
\end{proof}

The normalized gap in this family is $b/2^b$, which tends to zero.
It therefore does not contradict fixed-gap linear stability.
It shows that a successful general error bound must retain some rank dependence.

The planar proof isolates its extra input: many short faces give many bounded-cut cyclic regions.
General high-girth cubic graphs have no corresponding face count.
The argument above therefore does not settle the unrestricted linear problem.
Theorem~\ref{thm:general-linear-path-stability} settles that problem through localization and shortest-path layers.

\section{Linear path stability with a fixed excluded minor}
\label{sec:path-minor-stability}

The following extension uses the Alon--Seymour--Thomas separator theorem.
It does not require vertex splitting, which can change an excluded-minor condition.

\begin{theorem}[Fixed excluded minor]\label{thm:path-excluded-minor}
Fix an integer $h\ge2$ and a real number $\eta>0$.
There is a constant $C_{h,\eta}$ with the following property.
Let $N$ be a connected simple graph with no $K_h$ minor and rank $b$.
If
\[
 P_{uv}(N)\ge(1/2+\varepsilon)2^b,
 \qquad 0<\varepsilon\le1/2,
\]
then every $u$--$v$ path of length $\ell$ satisfies
\begin{equation}\label{eq:path-minor-linear}
 b\le(3+\eta)\ell+\frac{C_{h,\eta}}{\varepsilon}.
\end{equation}
In particular, for fixed $h$ and $\varepsilon$, every such path has length $\Omega_{h,\varepsilon}(b)$.
The conclusion applies to every fixed proper minor-closed class.
\end{theorem}

\subsection{Degree and rank restrictions from path density}

\begin{lemma}\label{lem:path-route-reduction}
Suppose that $P_{uv}(N)>2^{b-1}$, and fix a terminal path $A$ of length $\ell$.
There is a connected loopless multigraph $\Gamma$ with the following properties.
\begin{enumerate}
\item Its underlying simple graph is a minor of $N$.
\item Its rank and terminal path count equal those of $N$, and it retains $A$.
\item Every vertex outside $A$ has degree between three and five.
\item It has at most $2b+\ell+1$ vertices, and $\beta(\Gamma-V(A))\ge b-3\ell$.
\end{enumerate}
Degrees count parallel edges separately.
\end{lemma}

\begin{proof}
The blocks on the terminal route have ranks $b_i$ and terminal path counts $P_i$.
Let $s$ be the sum of the ranks of all other blocks.
The affine join bound gives $P_i\le2^{b_i}$, and
\[
 \frac{P_{uv}(N)}{2^b}=2^{-s}\prod_i\frac{P_i}{2^{b_i}}.
\]
The strict density assumption implies $s=0$ and $P_i/2^{b_i}>1/2$ for every route block.
Delete all edges outside the terminal route.
These edges belong to trees, so their deletion preserves both the rank and the terminal paths.

Consider a two-connected route block $B$, with its two route terminals.
Deleting any vertex leaves a connected graph.
Thus every incident edge pattern with the required parity extends to a terminal join.
The restriction of a uniform join to the incident edges is consequently uniform among those patterns.
At an internal vertex of degree $d$, a simple terminal path uses zero or two incident edges.
At a terminal, it uses exactly one.
The corresponding probability bounds are
\[
 \frac{1+\binom d2}{2^{d-1}},\qquad \frac{d}{2^{d-1}}.
\]
The first is at most one half for $d\ge6$, and the second for $d\ge4$.
Therefore internal block degrees are at most five and terminal block degrees are at most three.
Bridge blocks also satisfy these bounds.

Every articulation on the terminal route lies on $A$.
Consequently, all vertices outside $A$ have degree at most five.
Suppress degree-two vertices outside $A$ until none remain.
Suppression preserves the rank and gives a bijection on terminal paths.
It is an edge contraction, with parallel edges retained.
It therefore preserves the required minor exclusion for the underlying simple graph.
No loop arises, since each cyclic route block retains its two distinct terminals.
The resulting graph is $\Gamma$.

Write $n=|V(\Gamma)|$ and $a=|V(A)|=\ell+1$.
The degree lower bounds give
\[
 2(n+b-1)\ge3(n-a)+2(a-2)+2=3n-a-2.
\]
Hence $n\le2b+\ell+1$.

For the rank bound, consider one route block after suppression.
Let its part of $A$ have length $t\ge1$.
The sum of degrees at its $t+1$ marked vertices is at most $5(t-1)+6=5t+1$.
These vertices contain the $t$ path edges.
At most $4t+1$ edges therefore disappear when the marked vertices are deleted.
If an unmarked vertex remains, the surviving rank is at least
\[
 \beta(B)+(t+1)-(4t+1)=\beta(B)-3t.
\]
If no unmarked vertex remains, then
$\beta(B)\le(5t+1)/2-t\le3t$, giving the same conclusion.
After the marked vertices are removed, the route blocks become vertex-disjoint.
Their ranks add, and their marked path lengths sum to $\ell$.
Thus $\beta(\Gamma-V(A))\ge b-3\ell$.
\end{proof}

\subsection{Bounded components from separators}

We use the following form of the Alon--Seymour--Thomas theorem.
An $m$-vertex graph with no $K_h$ minor has a vertex set of size at most $h^{3/2}\sqrt m$.
After its deletion, every component has at most $m/2$ vertices.
This is statement (1.4) of their original paper.\footnote{N. Alon, P. Seymour, and R. Thomas,
\emph{A separator theorem for nonplanar graphs}, Journal of the American Mathematical Society 3 (1990), 801--808.
The original paper is available at \url{https://thomas.math.gatech.edu/PAP/nonplsep.pdf}.}

\begin{lemma}\label{lem:path-small-components}
Let $J$ have $m$ vertices, and suppose its underlying simple graph has no $K_h$ minor.
For every integer $K\ge1$, there is a set $S$ such that
\[
 |S|\le\frac{4h^{3/2}m}{\sqrt K},
\]
and every component of $J-S$ has at most $K$ vertices.
\end{lemma}

\begin{proof}
Apply the separator theorem repeatedly to components with more than $K$ vertices.
Let $S$ contain all separators used in this recursion.
For a processed component of order $q$, charge $h^{3/2}/\sqrt q$ to each of its vertices.
The total charge covers its separator size.
For any fixed vertex, the orders of successive components containing it decrease by a factor of at least two.
Read that sequence in reverse order.
Its smallest order exceeds $K$, so the vertex's total charge is at most
\[
 \frac{h^{3/2}}{\sqrt K}\sum_{j\ge0}2^{-j/2}
 =\frac{(2+\sqrt2)h^{3/2}}{\sqrt K}
 <\frac{4h^{3/2}}{\sqrt K}.
\]
Sum over the vertices to obtain the stated bound.
\end{proof}

\begin{proof}[Proof of Theorem~\ref{thm:path-excluded-minor}]
Choose
\[
 \delta=\frac{\eta}{28+8\eta},\qquad
 K=\left\lceil\frac{16h^3}{\delta^2}\right\rceil,
 \qquad D=3K.
\]
Then $0<\delta<1/8$.
Apply Lemma~\ref{lem:path-route-reduction}, and put $J=\Gamma-V(A)$.
Lemma~\ref{lem:path-small-components} gives a set $S$ with
\[
 |S|\le\delta|V(J)|\le\delta(2b+\ell+1)
\]
such that every component of $J-S$ has at most $K$ vertices.
Deleting one vertex of degree at most five decreases rank by at most four.
Thus
\begin{equation}\label{eq:path-minor-surviving-rank}
 \beta(J-S)\ge (1-8\delta)b-(3+4\delta)\ell-4\delta.
\end{equation}

Every component of $J-S$ has maximum degree at most five and rank at most $2K$.
Choose one cycle from every cyclic component, obtaining $k$ vertex-disjoint cycles.
Their vertex sets lie outside $A$, and
\[
 2Kk\ge\beta(J-S).
\]
Each selected vertex has degree at most five in $\Gamma$ and uses two incidences in its cycle.
Every chosen cyclic region therefore has boundary at most $3K=D$.

Add a new two-edge terminal path to $\Gamma$, and mark its union with $A$.
The path injection and \eqref{eq:path-small-cuts} give
\[
 k\le\frac{(D-1)2^{D-2}}{\varepsilon}.
\]
For $k=0$, the same inequality holds trivially.
Combining these estimates with \eqref{eq:path-minor-surviving-rank} gives
\[
 (1-8\delta)b
 \le(3+4\delta)\ell+4\delta
 +\frac{2K(D-1)2^{D-2}}{\varepsilon}.
\]
The definition of $\delta$ gives
$(3+4\delta)/(1-8\delta)=3+\eta$.
Since $\varepsilon\le1/2$, one may take
\[
 C_{h,\eta}
 =\frac{2\delta+2K(D-1)2^{D-2}}{1-8\delta}.
\]
This proves \eqref{eq:path-minor-linear}.
Since $\ell\ge1$, it also gives
\[
 \ell\ge\frac{\varepsilon b}{(3+\eta)\varepsilon+C_{h,\eta}}.
\]
Finally, a proper minor-closed class omits some fixed finite graph $H$.
It then omits $K_{|V(H)|}$, because that complete graph contains $H$ as a minor.
\end{proof}

The constants above are deliberately explicit and are not optimized.
The proof establishes a structural consequence of the existing small-cut estimate and the classical separator theorem.
It does not provide a separator argument for arbitrary graphs.

\section{General linear path stability}
\label{sec:general-linear-path-stability}

Fable 5.2 proposed the localization and shortest-path layer argument in its November 10 research notes.
The Codex root and path-stability agents independently audited that argument.
The proof below includes the minor padding correction found in both audits.
It uses Theorem~\ref{thm:research-marked-size} to supply the required forest estimate.
Thus the conclusion is a mathematical theorem, rather than a conditional reduction.
It has not been formalized in Lean.

All logarithms in this section have base two, except $\ln$.
Put $\Phi(x)=x/\log^2x$ for $x\ge8$.
This function is increasing on its domain.

\begin{theorem}[General linear stability]\label{thm:general-linear-path-stability}
For every $0<\varepsilon\le1/2$, there is a constant $c_\varepsilon>0$ with the following property.
Let $N$ be a connected simple graph of rank $b$, with distinct terminals $u,v$.
If
\[
 P_{uv}(N)\ge(1/2+\varepsilon)2^b,
\]
then every $u$--$v$ path has length at least $c_\varepsilon b$.
In particular, $\operatorname{dist}_N(u,v)\ge c_\varepsilon b$.

An explicit expression in terms of the forest threshold is available.
Set $\delta=\varepsilon/2$, and let $R_\delta$ be an integer as in Lemma~\ref{lem:general-fe2} below.
Put
\[
 M=2R_\delta,\qquad
 \Theta=\max\left\{2^{36},\frac{(M+1)M2^{M+1}}{\delta}\right\}.
\]
One may take $c_\varepsilon=(\Theta+1)^{-2}$.
Every terminal path also satisfies the stronger inequality
\begin{equation}\label{eq:general-linear-strong}
 \ell>\frac{2b}{\Theta+1}-(\Theta+1).
\end{equation}
\end{theorem}

The formula determines a positive constant once the forest threshold is fixed.
The present proof does not give a numerical value for that threshold.

\subsection{The forest input and its localization}

\begin{lemma}\label{lem:general-fe2}
Fix $0<\delta<1/2$.
There is an integer $R_\delta\ge2^{100}$ such that the following statement holds.
If $G$ is connected and simple, $\rho=\beta(G)\ge R_\delta$, and $W$ is one cycle with
$|E(W)|\le\Phi(\rho)$, then
\[
 p(G,W)\le1/2+\delta.
\]
\end{lemma}

\begin{proof}
Use $a_\delta$ from Theorem~\ref{thm:research-marked-size}.
Choose $R_\delta$ above that theorem's rank threshold and so that
$\log R_\delta\ge\max\{100,1/a_\delta\}$.
For $\rho\ge R_\delta$,
\[
 \Phi(\rho)=\frac{\rho}{\log^2\rho}
 \le a_\delta\frac{\rho}{\log\rho}.
\]
The single-component condition also holds, since $1\le\log\rho/100$.
Theorem~\ref{thm:research-marked-size} applies.
\end{proof}

We also use the manuscript's small-cut estimate in the following form.
If $q$ disjoint cyclic regions avoid $V(W)$ and each has boundary at most an integer $D\ge1$, then
\begin{equation}\label{eq:general-small-cut-input}
 p(G,W)\le\frac12+\frac{(D+1)2^{D-1}}{q}.
\end{equation}
This statement permits loops and parallel edges.
It follows directly from the cut formula, the exceptional-partner bound, and the second-moment estimate.

\begin{lemma}[Localization]\label{lem:general-localization}
Let $G$ be a connected multigraph, and let $R$ be a nonempty connected proper vertex set.
Suppose that $G-R$ is connected, and write $d=|\delta_G R|\ge1$.
For each boundary edge $f$, with endpoint $z_f\in R$, add a distinct new vertex $t_f$ and an edge $z_ft_f$.
Add a cycle $W_R$ through the vertices $t_f$, padding it with new vertices if $d<3$.
Together with $G[R]$, this gives a graph $G_R$ satisfying
\[
 |E(W_R)|=\max\{d,3\},\qquad
 \beta(G_R)=\beta(G[R])+d.
\]
For uniform even subgraphs,
\begin{equation}\label{eq:general-localization-law}
 \Pr\bigl(\xi_G\cap E(G[R])\text{ is a forest}\bigr)=p(G_R,W_R).
\end{equation}
If $R\cap V(W)=\varnothing$, then $p(G,W)\le p(G_R,W_R)$.
\end{lemma}

\begin{proof}
Restriction outside an edge set $T$ maps the cycle space onto the cycle space after contracting every component of $T$.
For surjectivity, prescribe an even edge set in the contraction.
Its boundary inside each component of $T$ has even size.
A join in a spanning tree of that component completes the prescribed edge set.
This is a surjective linear map and therefore preserves uniformity on its image.

Contract $E(G-R)$ in $G$.
The resulting graph consists of $G[R]$ and one vertex receiving its $d$ boundary edges.
Contract $E(W_R)$ in $G_R$.
The same graph results, so the restrictions to $E(G[R])$ have the same law.
The added spokes have distinct leaf endpoints outside $W_R$.
They cannot create a cycle outside $W_R$, which proves \eqref{eq:general-localization-law}.
The stated inequality follows because $E(G[R])$ avoids $E(W)$.

For $d\ge3$, the construction adds $d$ vertices and $2d$ edges, giving the rank formula.
For $d=1,2$, padding the cycle adds equally many further edges and vertices.
Thus the rank formula remains valid.
\end{proof}

\begin{lemma}[Connected complementary shore]\label{lem:general-shore-closure}
Let $G$ and $W$ be connected, and let $R$ be a connected vertex set avoiding $V(W)$.
Let $C_0$ be the component of $G-R$ containing $W$, and put $R^+=V(G)\setminus V(C_0)$.
Then $R^+$ is connected, contains $R$, avoids $V(W)$, and satisfies
$\delta_G R^+\subseteq\delta_G R$.
Its complement is connected.
\end{lemma}

\begin{proof}
Every component of $G-R$ other than $C_0$ has an edge to $R$, since $G$ is connected.
Their union with $R$ is therefore connected.
Edges between $R^+$ and $C_0$ must have their other endpoint in $R$.
This gives the boundary inclusion and the remaining conclusions.
\end{proof}

\subsection{A shortest path in a cubic expansion}

\begin{lemma}\label{lem:general-shortest-reduction}
Suppose $P_{uv}(N)>2^{b-1}$.
There is a connected loopless multigraph $\Gamma$, with distinct terminals $u',v'$ and a shortest terminal path $A$, such that
\begin{enumerate}
\item $\beta(\Gamma)=b$ and $P_{u'v'}(\Gamma)\ge P_{uv}(N)$;
\item every nonterminal vertex of $\Gamma$ has degree three, and both terminals have degree one or two;
\item the length $\ell_0$ of $A$ is at most the length of every terminal path in $N$;
\item $|V(\Gamma)|\ge2b$, and $U=V(\Gamma)\setminus V(A)$ has size at least $2b-\ell_0-1$.
\end{enumerate}
Every vertex of $A$ has at most one neighbor incidence outside $A$.
\end{lemma}

\begin{proof}
By Proposition~\ref{prop:research-offroute}, no cyclic block lies outside the terminal route.
Delete the trees outside that route.
Suppress every degree-two vertex other than $u,v$.
This preserves rank and path count, and it can only shorten terminal paths.
No loop arises, since a cyclic block on the route retains its two distinct route terminals.
Call the resulting kernel $N_1$.
Every nonterminal vertex of $N_1$ has degree at least three.

Choose a shortest terminal path $A_1$ in $N_1$.
At each internal vertex of $A_1$, retain one vertex incident with its two path edges.
Attach all other incident edges through a rooted trivalent tree in its third branch.
At each terminal, retain one vertex incident with the path edge and attach its other edges through one rooted tree.
A single remaining incident edge needs no internal tree vertex.
Replace every other vertex of degree at least four by a trivalent tree.
No embedding restriction is imposed on these trees.

These replacements preserve rank.
Every original terminal path has a unique simple lift, so path count cannot decrease.
The image $A$ of $A_1$ has exactly the same length $\ell_0$.
Contracting the replacement trees maps any new terminal path to a terminal walk in $N_1$ of no greater length.
Every such walk has length at least the terminal distance in $N_1$, namely $\ell_0$.
Thus $A$ remains shortest.
This argument uses a shortest path after suppression, not an arbitrary path in the original graph.

All nonterminal degrees are three, and both terminal degrees are at most two.
If $n=|V(\Gamma)|$, the degree sum gives
\[
 2(n+b-1)=3(n-2)+\deg(u')+\deg(v').
\]
Hence $n=2b+4-\deg(u')-\deg(v')\ge2b$.
Subtracting $|V(A)|=\ell_0+1$ proves the last assertion.
The construction also gives the asserted bound on incidences from $A$ to $U$.
\end{proof}

\subsection{Wide distance layers give a region with small boundary}

Write $A=w_0w_1\cdots w_{\ell_0}$.
Let $L_j$ be the distance-$j$ layer from $u'=w_0$ in $\Gamma$.
Since $A$ is shortest, $w_j\in L_j$.
Put $U_j=L_j\cap U$ and $m_j=|U_j|$, taking $m_j=0$ for nonexistent layers.
For integers $i\le i'$, put
\[
 R(i,i')=U_i\cup\cdots\cup U_{i'},\qquad
 m(i,i')=\sum_{j=i}^{i'}m_j.
\]
Every edge joins vertices in the same or consecutive distance layers.
At most one edge joins each marked vertex to $U$.
Consequently,
\begin{equation}\label{eq:general-window-boundary}
 |\delta_\Gamma R(i,i')|
 \le3m_{i-1}+3m_{i'+1}+(i'-i+3).
\end{equation}
The last term counts the possible marked vertices in layers $i-1,\ldots,i'+1$.

For a parameter $\Theta\ge2^{36}$, call $R\subseteq U$ good if
\[
 |R|\ge\Theta,\qquad |\delta_\Gamma R|\le\Phi(|R|)/4.
\]
A good set need not be connected.

\begin{lemma}[Growing a window]\label{lem:general-window}
If $m_j\ge\Theta$ for some $j$, then a good set exists.
\end{lemma}

\begin{proof}
Start with the window $[j,j]$.
While a neighboring layer has mass greater than $\Phi(m(i,i'))/48$, include that layer.
Repeat until neither neighboring layer satisfies the condition.
Only finitely many layers are nonempty, so the procedure stops.

Let $\mu$ be the final mass, and set $L=\log\mu\ge36$.
Every step multiplies the current mass by more than $1+1/(48L^2)$.
Thus the number of steps is at most
\[
 \frac{\ln\mu}{\ln(1+1/(48L^2))}
 \le96L^2\ln\mu<67L^3.
\]
At termination, each neighboring layer has mass at most $\Phi(\mu)/48$.
Equation~\eqref{eq:general-window-boundary} therefore gives
\[
 |\delta_\Gamma R(i,i')|
 \le\Phi(\mu)/8+67L^3+3\le\Phi(\mu)/4.
\]
The final inequality follows from
$2^L\ge8L^2(67L^3+3)$ for $L\ge36$.
It holds at $L=36$, and the ratio of its left side to its right side increases thereafter.
Since $\mu\ge\Theta$, the final window is good.
\end{proof}

\begin{lemma}[No good set forces a long path]\label{lem:general-layer-count}
If no good subset of $U$ exists, then
\[
 \ell_0>\frac{2b}{\Theta+1}-(\Theta+1).
\]
\end{lemma}

\begin{proof}
Lemma~\ref{lem:general-window} gives $m_j<\Theta$ in every layer.
Let $R$ be the union of the layers $U_j$ with $j\ge\ell_0+2$.
No vertex of $A$ has a neighbor in $R$.
Therefore $|\delta_\Gamma R|\le3m_{\ell_0+1}<3\Theta$.
If $|R|\ge\Theta^2$, monotonicity of $\Phi$ and $\Phi(\Theta^2)/4\ge3\Theta$ would make $R$ good.
Hence $|R|<\Theta^2$.
Counting the remaining layers gives
\[
 2b-\ell_0-1\le|U|<(\ell_0+2)\Theta+\Theta^2.
\]
Rearranging proves the assertion.
\end{proof}

\subsection{A good set contradicts excess path density}

Add a new vertex $x$ and edges $u'x,xv'$ to $\Gamma$.
Call the resulting graph $H$, and mark the cycle $W=A\cup\{u'x,xv'\}$.
The proof of Lemma~\ref{lem:research-path-injection} also applies to multigraphs, giving
\begin{equation}\label{eq:general-path-forest}
 P_{u'v'}(\Gamma)\le p(H,W)2^b.
\end{equation}

\begin{lemma}[The good-set dichotomy]\label{lem:general-good-set}
Use $\delta,R_\delta,M,\Theta$ from Theorem~\ref{thm:general-linear-path-stability}.
If a good subset of $U$ exists, then $p(H,W)\le1/2+\delta$.
\end{lemma}

\begin{proof}
Let $R$ be good, with $m=|R|$ and $d=|\delta_\Gamma R|\le m/(4\log^2m)$.
Let $R_s$ be the vertex sets of the components of $\Gamma[R]$.
Their boundaries are disjoint and have total size $d$.
Every vertex in $R$ is cubic in both $\Gamma$ and $H$, so
\[
 \beta(\Gamma[R_s])=\frac{|R_s|-|\delta_\Gamma R_s|}{2}+1.
\]
Separate the components into four classes:
\begin{enumerate}
\item $|R_s|\ge M$ and $|\delta_\Gamma R_s|\le\Phi(|R_s|)/2$;
\item $|R_s|\ge M$ and $|\delta_\Gamma R_s|>\Phi(|R_s|)/2$;
\item $|R_s|<M$ and $\Gamma[R_s]$ is a tree;
\item $|R_s|<M$ and $\Gamma[R_s]$ contains a cycle.
\end{enumerate}

Class 2 has total size less than $2d\log^2m\le m/2$.
Indeed, each of its components satisfies
$|R_s|<2|\delta_\Gamma R_s|\log^2|R_s|\le2|\delta_\Gamma R_s|\log^2m$.
A tree component satisfies $|\delta_\Gamma R_s|=|R_s|+2$.
Thus class 3 has total size less than $d\le m/4$.
A cyclic component in class 4 has boundary at most its size, hence less than $M$.

Suppose first that class 1 is empty.
Class 4 then has total size greater than $m/4$, so it contains more than $m/(4M)$ components.
These are disjoint cyclic regions avoiding $W$, with boundaries less than $M$.
By \eqref{eq:general-small-cut-input},
\[
 p(H,W)\le\frac12+
 \frac{(M+1)2^{M-1}\,4M}{m}
 \le\frac12+\delta.
\]
The last inequality uses $m\ge\Theta$.

Suppose instead that class 1 contains $R_s$.
Apply Lemma~\ref{lem:general-shore-closure} in $H$ to obtain $R^+\supseteq R_s$ with connected complement.
It avoids $V(W)$, so it is contained in $U$ and all its vertices are cubic.
Write $d^+=|\delta_H R^+|$ and
\[
 \rho=\beta(H[R^+])+d^+
 =\frac{|R^+|+d^+}{2}+1
 \ge|R_s|/2\ge R_\delta.
\]
Since $\delta_HR_s=\delta_\Gamma R_s$, the boundary comparison gives
\[
 d^+\le|\delta_\Gamma R_s|
 \le\Phi(|R_s|)/2
 \le\Phi(|R_s|/2)
 \le\Phi(\rho).
\]
The middle inequality uses $\log(s/2)=\log s-1$.

Localize at $R^+$ using Lemma~\ref{lem:general-localization}.
If $d^+\ge3$, its marked cycle has at most $\Phi(\rho)$ edges by the last display.
If $d^+=1$ or $2$, its marked cycle has three edges, which still meets that bound since $\rho\ge2^{100}$.
Subdivide any parallel interior edges until the localized graph is simple.
Such subdivisions preserve rank, marked-cycle length, and outside-forest probability.
Lemma~\ref{lem:general-fe2} now yields
\[
 p(H,W)\le p(H_{R^+},W_{R^+})\le1/2+\delta.
\]
This completes both cases.
\end{proof}

\begin{proof}[Proof of Theorem~\ref{thm:general-linear-path-stability}]
Use Lemma~\ref{lem:general-shortest-reduction}.
The density assumption and $\delta=\varepsilon/2$ give
\[
 p(H,W)\ge\frac{P_{u'v'}(\Gamma)}{2^b}
 \ge\frac12+\varepsilon>\frac12+\delta.
\]
Lemma~\ref{lem:general-good-set} excludes every good subset of $U$.
Lemma~\ref{lem:general-layer-count} proves \eqref{eq:general-linear-strong} for $\ell_0$.
Every original terminal path has length at least $\ell_0$, so the same inequality holds for its length $\ell$.

Put $Q=\Theta+1$.
If $b\ge Q^2$, then $\ell>2b/Q-Q\ge b/Q\ge b/Q^2$.
If $b<Q^2$, then $\ell\ge1>b/Q^2$.
Thus $c_\varepsilon=Q^{-2}$ works at every rank.
\end{proof}

\subsection{Structural consequences}

\begin{corollary}\label{cor:general-short-path-density}
For any sequence of connected graphs with ranks $b\to\infty$ and terminal distances $o(b)$,
\[
 P_{uv}(N)\le(1/2+o(1))2^b.
\]
If $P_{uv}(N)\ge(1-\alpha)2^b$ for a fixed $0\le\alpha<1/2$, every terminal path has length
$\Omega_\alpha(b)$.
\end{corollary}

\begin{proof}
A fixed positive density gap in the first assertion would contradict Theorem~\ref{thm:general-linear-path-stability}.
For the second assertion, use $\varepsilon=1/2-\alpha$.
\end{proof}

\begin{corollary}[Bounded width of the terminal kernel]\label{cor:general-bounded-pathwidth}
Under the hypotheses of Theorem~\ref{thm:general-linear-path-stability}, remove all blocks outside the terminal route and suppress nonterminal degree-two vertices.
The resulting terminal kernel has pathwidth at most $2\lceil\Theta\rceil+1$.
The same bound holds for the cubic expansion in Lemma~\ref{lem:general-shortest-reduction}.
\end{corollary}

\begin{proof}
The proof excludes good sets, so Lemma~\ref{lem:general-window} gives $|U_j|<\Theta$ in every distance layer.
At most one marked path vertex lies in each layer.
Therefore each full layer has at most $\lceil\Theta\rceil+1$ vertices.
The bags $L_j\cup L_{j+1}$ form a path decomposition, since every edge joins equal or consecutive layers.
Their sizes give the stated pathwidth bound.
The terminal kernel is obtained from the cubic expansion by contracting its replacement trees.
Pathwidth does not increase under taking minors.
\end{proof}

The restriction to the terminal kernel matters.
Trees attached outside the terminal route can have arbitrarily large pathwidth without changing rank or terminal path count.

\section{General additive path stability}
\label{sec:general-additive-path-stability}

The bounded-layer conclusion of the preceding proof yields a stronger result.
It forces every terminal path to have length at least the rank, up to a constant depending on the density gap.
This deduction was developed by the Codex path-stability agent and independently audited by the root and four-regular agents.
It uses Fable's localization and layer argument through the preceding section.
The result has not been formalized in Lean.

\begin{theorem}[Additive path stability]\label{thm:general-additive-path-stability}
For every $0<\varepsilon\le1/2$, there is a constant $C_\varepsilon$ such that the following holds.
If $N$ is a connected simple graph of rank $b$, with distinct terminals $u,v$, and
\[
 P_{uv}(N)\ge(1/2+\varepsilon)2^b,
\]
then every $u$--$v$ path has length at least
\[
 b-C_\varepsilon.
\]
In particular, $\operatorname{dist}_N(u,v)\ge b-C_\varepsilon$.

Let $\Theta$ be the constant from Theorem~\ref{thm:general-linear-path-stability}, and put
\[
 B=\lceil\Theta\rceil,\qquad L=20B+2.
\]
One may take
\begin{equation}\label{eq:general-additive-constant}
 C_\varepsilon=\frac{2L(L+1)2^{L-1}}{\varepsilon}.
\end{equation}
For the cubic expansion $\Gamma$ and shortest path $A$ from Lemma~\ref{lem:general-shortest-reduction},
one also has $\beta(\Gamma-V(A))\le C_\varepsilon$.
\end{theorem}

The size of this constant is not optimized.
Its finiteness gives the qualitative additive conclusion.

\subsection{A short cycle in any cyclic layer interval}

Retain $\Gamma,A,U_j$ from the preceding proof.
All vertices outside $A$ are cubic, and each vertex of $A$ has at most one incidence outside $A$.
Assume that $|U_j|\le B$ in every distance layer.

\begin{lemma}[Minimum layer span]\label{lem:additive-minimum-span}
If the subgraph induced by $\bigcup_{j\in I}U_j$ contains a cycle for an interval $I$ of layer indices,
then it contains a cycle of length at most $20B+2$.
\end{lemma}

\begin{proof}
Among all cycles in that subgraph, choose one whose layer span is smallest.
Let its minimum and maximum layer indices be $a,b'$, and put $s=b'-a$.
Let $R=U_a\cup\cdots\cup U_{b'}$.
Every cycle in $\Gamma[R]$ has span at least $s$, by the choice of the first cycle.
Since it lies in this same interval, it must meet $U_a$.
Thus $\Gamma[R]-U_a$ is a forest.

Deleting a vertex of degree at most three decreases rank by at most two.
Consequently,
\[
 \beta(\Gamma[R])\le2|U_a|\le2B.
\]
Put $m=|R|$, $d=|\delta_\Gamma R|$, and $c=\operatorname{comp}(\Gamma[R])\ge1$.
Cubicity outside $A$ gives the exact identity
\[
 m=2\beta(\Gamma[R])+d-2c.
\]
The layer-boundary estimate \eqref{eq:general-window-boundary} yields
\[
 d\le6B+s+3.
\]
Combining these bounds gives
\begin{equation}\label{eq:additive-span-size}
 m\le10B+s+1.
\end{equation}

The chosen cycle has length at least $2s$.
Indeed, its two arcs between a minimum-layer vertex and a maximum-layer vertex each have at least $s$ edges.
This also holds when $s=0$.
Hence $2s\le m\le10B+s+1$, so $s\le10B+1$.
Equation~\eqref{eq:additive-span-size} now gives $m\le20B+2$.
The chosen cycle has at most $m$ edges, proving the assertion.
\end{proof}

\begin{lemma}[Greedy deletion of whole windows]\label{lem:additive-layer-deletion}
Put $J=\Gamma-V(A)$ and $L=20B+2$.
There are $k$ vertex-disjoint cycles in $J$, each of length at most $L$, such that
\[
 \beta(J)\le2Lk.
\]
\end{lemma}

\begin{proof}
Start with all unmarked layers.
If one remaining layer interval contains a cycle, choose a cycle of minimum span in that interval.
Delete its entire tight layer window $R$, as defined in Lemma~\ref{lem:additive-minimum-span}.
That lemma's proof gives $|R|\le L$.
The remaining vertices again form whole layer intervals, and no edge joins two different intervals.
Repeat until the remaining graph is a forest.

Suppose that this procedure selects $k$ cycles, and let $Z$ be the union of the deleted windows.
The selected cycles are vertex-disjoint, have length at most $L$, and satisfy $|Z|\le Lk$.
Since $J$ has maximum degree at most three, deleting $Z$ decreases rank by at most $2|Z|$.
Therefore $\beta(J)\le2|Z|\le2Lk$.
\end{proof}

\subsection{Completing the additive estimate}

\begin{proof}[Proof of Theorem~\ref{thm:general-additive-path-stability}]
Use Lemma~\ref{lem:general-shortest-reduction} to obtain $\Gamma$ and its shortest terminal path $A$ of length $\ell_0$.
In the proof of Theorem~\ref{thm:general-linear-path-stability}, the density assumption excludes every good set.
Lemma~\ref{lem:general-window} therefore gives $|U_j|<\Theta\le B$ for every layer.

Use Lemma~\ref{lem:additive-layer-deletion} to choose $k$ vertex-disjoint cycles of length at most $L$ in $J=\Gamma-V(A)$.
Each of their vertex sets is a cyclic region with boundary at most $L$ in $\Gamma$.
Every vertex of such a cycle has degree three and uses two incidences in that cycle.
Add the new two-edge terminal path and mark its union with $A$, giving $H,W$ as before.
The chosen region boundaries are unchanged, and all the regions avoid $V(W)$.

If $k\ge1$, the path injection and \eqref{eq:general-small-cut-input} give
\[
 \frac12+\varepsilon\le p(H,W)
 \le\frac12+\frac{(L+1)2^{L-1}}{k}.
\]
Thus $k\le(L+1)2^{L-1}/\varepsilon$.
The same upper bound holds when $k=0$.
Lemma~\ref{lem:additive-layer-deletion} yields
\[
 \beta(J)\le2Lk\le C_\varepsilon.
\]

It remains to compare this rank with $b-\ell_0$.
The degree sum at the $\ell_0+1$ vertices of $A$ is at most $3(\ell_0-1)+4=3\ell_0+1$.
These vertices contain the $\ell_0$ path edges.
Deleting them therefore removes at most $2\ell_0+1$ edges.
If $J$ is nonempty, its rank satisfies
\[
 \beta(J)\ge b+(\ell_0+1)-(2\ell_0+1)=b-\ell_0.
\]
If $J$ is empty, then $|E(\Gamma)|\le(3\ell_0+1)/2$, so
\[
 b=|E(\Gamma)|-\ell_0\le(\ell_0+1)/2\le\ell_0.
\]
Thus $b-\ell_0\le\beta(J)$ in both cases.
It follows that $\ell_0\ge b-C_\varepsilon$.
Every original terminal path has length at least $\ell_0$, which proves the theorem.
\end{proof}

\begin{corollary}\label{cor:general-rank-deficit}
For any sequence of connected graphs and terminal pairs with
$\beta(N)-\operatorname{dist}_N(u,v)\to\infty$,
\[
 P_{uv}(N)\le(1/2+o(1))2^{\beta(N)}.
\]
\end{corollary}

\begin{proof}
A fixed positive density gap would bound the rank-distance deficit by Theorem~\ref{thm:general-additive-path-stability}.
\end{proof}

The conclusion does not yet give an exponential error in the rank-distance deficit.
Such a rate would require quantitative control beyond the constants supplied here.

The script \texttt{check\_additive\_layer\_structure.py} checks the finite layer argument on 157 explicit cubic expansions.
It verifies 301 selected windows, including examples with 322 vertices and outside rank 81.
The checks cover the minimum-span feedback set, cubic rank identity, boundary bound, and final greedy rank estimate.
They support these finite identities and do not replace the probability argument.

\section{A quantitative inverse for the small-cut estimate}\label{sec:inverse}

This section refines the small-cut argument from the source manuscript.
It gives both a pair-structure conclusion and a concentration conclusion.
It does not give a decomposition of every graph with many cycles.

Let $G$ be a connected multigraph.
Let $A_1,\ldots,A_k$ be disjoint cyclic regions, where $k\ge2$ and $|\delta A_i|\le d$.
Set $\Delta=\max\{d-2,0\}$ and $K=1+\Delta$.
For a uniform even subgraph $\xi$, define
\[
 E_i=\{\xi\cap\delta A_i=\varnothing,\quad
             \xi\cap E(G[A_i])\ne\varnothing\},\qquad
 p_i=\Pr(E_i),\quad Z=\sum_i\mathbf1_{E_i},\quad \mu=\sum_i p_i.
\]
Let $\ell_{ij}$ count the links between $A_i$ and $A_j$.
The source manuscript gives
\[
 \Pr(E_i\cap E_j)=2^{\ell_{ij}-1}p_i p_j.
\]
Every region is proper, since $k\ge2$.
The cut formula therefore improves the elementary lower bound to
\begin{equation}\label{eq:pi-lower}
 p_i=2^{\operatorname{comp}(G-A_i)-|\delta A_i|}
       (1-2^{-\beta(G[A_i])})\ge2^{-d},\qquad \mu\ge k2^{-d}.
\end{equation}

\begin{theorem}\label{thm:inverse}
Suppose that $0\le\eta<1/2$ and
\[
 f:=\Pr\bigl(\xi\cap E(G[A_i])\text{ is a forest for every }i\bigr)
       \ge\tfrac12-\eta.
\]
Define the ordered one-link mass by
\[
 D_1=\frac1{\mu^2}\sum_{i\ne j:\,\ell_{ij}=1}p_i p_j.
\]
Then
\begin{equation}\label{eq:inverse-mass}
 D_1\le \frac{4\eta}{1+2\eta}+\frac K\mu
       \le \frac{4\eta}{1+2\eta}+\frac{K2^d}{k}.
\end{equation}
Choose $I,J$ independently with $\Pr(I=i)=\Pr(J=i)=p_i/\mu$.
Then
\begin{equation}\label{eq:two-link-mass}
 \Pr(I\ne J,\ \ell_{IJ}=2)
   \ge1-\frac{4\eta}{1+2\eta}-\frac{2K}{\mu}.
\end{equation}
Writing $\nu=\Pr(Z>0)$, we also have
\begin{gather}
 \frac1{2+K/\mu}\le\nu\le\tfrac12+\eta,\label{eq:nu}\\
 \frac{\operatorname{Var}(Z\mid Z>0)}{(\mu/\nu)^2}
       \le2\eta+(\tfrac12+\eta)\frac K\mu.\label{eq:occupancy}
\end{gather}
\end{theorem}

\begin{proof}
The exceptional-partner lemma bounds each region's partners with at least three links by $\Delta$.
For these pairs, use $\Pr(E_i\cap E_j)\le p_i$.
For one-link pairs, the product moment is $p_i p_j$.
For two-link pairs, it is $2p_i p_j$.
Together with the diagonal terms, this gives
\begin{equation}\label{eq:inverse-moment}
 \mathbb E Z^2\le2\mu^2-D_1\mu^2+K\mu.
\end{equation}
On the event defining $f$, every $E_i$ fails.
Thus $\Pr(Z=0)\ge f$, and $\nu\le1/2+\eta$.
Cauchy--Schwarz gives $\mathbb E Z^2\ge\mu^2/\nu$.
Substitution in \eqref{eq:inverse-moment} proves \eqref{eq:inverse-mass}.

The weighted diagonal mass is at most $1/\mu$, since $p_i\le1$.
The weighted mass of exceptional ordered pairs is at most $\Delta/\mu$.
Connectedness gives $\ell_{ij}\ge1$ for distinct regions.
Subtract these masses and \eqref{eq:inverse-mass} from one to obtain \eqref{eq:two-link-mass}.

Dropping $D_1$ from \eqref{eq:inverse-moment} gives the lower bound in \eqref{eq:nu}.
Finally, $\mathbb E(Z\mid Z>0)=\mu/\nu$, and
\[
 \frac{\operatorname{Var}(Z\mid Z>0)}{(\mu/\nu)^2}
   =\frac{\nu\mathbb E Z^2}{\mu^2}-1
   \le2\nu-1+\nu K/\mu.
\]
The upper bound on $\nu$ proves \eqref{eq:occupancy}.
\end{proof}

For fixed $d$, let $k\to\infty$ and $\eta\to0$.
The theorem forces almost all weighted pair mass onto distinct two-link pairs.
It also forces $Z/\mu$ to concentrate near two, conditionally on being positive.
That conditioning event has probability tending to one half.
These conclusions describe the near-equality mechanism of this particular second-moment argument.

\paragraph{Transfer from a cycle count.}
The chord construction gives $c(G)/2^r\le p(G^+,W)+2^{-t}$ when $t$ edges are selected.
Hence $c(G)\ge(1/2-\eta)2^r$ implies $p(G^+,W)\ge1/2-\eta-2^{-t}$.
Suppose that the regions above are disjoint from $V(W)$ in $G^+$.
The forest event outside $W$ implies the event defining $f$.
The theorem therefore applies with $\eta+2^{-t}$, when that number is less than one half.
A separate argument must produce many regions with bounded cuts.

\paragraph{A more general second-moment boundary.}
If a proposed improvement uses only $\mathbb E Z=\mu$ and $\mathbb E Z^2\le2\mu^2+K\mu$, it cannot force an error exponentially small in $\mu$.
For $K>0$, take $Z=0$ with probability $(\mu+K)/(2\mu+K)$ and $Z=2\mu+K$ otherwise.
Both moment conditions hold with equality, and
\[
 \Pr(Z=0)-\tfrac12=\frac{K}{2(2\mu+K)}.
\]
This is a limitation of those moment inequalities alone.
It is not a graph construction or an obstruction to stronger graph estimates.

\section{A cycle-count recurrence without witness chords}
\label{sec:cycle-recursion}

Fable 5.2 proposed this use of the Entringer--Slater deletion identity in its November 10 notes.
The recurrence gives a shorter route from a one-cycle forest estimate to the cycle-count asymptotic theorem.
It does not change the existing formal proof.
It also removes the separate witness-avoidance error from the quantitative argument.

Let $\psi(r)$ be the maximum number of cycles among finite graphs of rank $r$.
The same maximum results if one requires simple connected graphs.
Bridges can join different components, and subdividing each edge twice removes loops and parallel edges without changing cycles or rank.
Since cycle counts are integers bounded by $2^r-1$, this maximum is attained.

\begin{lemma}[One marked edge]\label{lem:recursion-marked-edge}
Let $G$ be a connected multigraph of rank $r$, let $W$ be a cycle, and let $e\in E(W)$.
Then
\[
 c_e(G)\le p(G,W)2^{r-1}.
\]
\end{lemma}

\begin{proof}
A cycle $C$ containing $e$ has a forest outside $W$.
If two such cycles $C,C'$ have the same outside part, their symmetric difference is an even edge set in $W-e$.
Since $W-e$ is a forest, that symmetric difference is empty.
Thus these cycles give distinct forest outside parts.
There are $2^{r-1}$ outside parts, since $W$ has rank one.
Their forest proportion is $p(G,W)$, proving the bound.
\end{proof}

\begin{proposition}[Recurrence from one-cycle estimates]\label{prop:cycle-forest-recurrence}
Suppose that for each sufficiently large integer $r$, every connected simple graph $G$ of rank $r$
and every cycle $W$ with at most $\sqrt r$ edges satisfy
\[
 p(G,W)\le1/2+d(r),
\]
where $d(r)\ge0$.
Then, for all sufficiently large $r$,
\begin{equation}\label{eq:cycle-deletion-recurrence}
 \psi(r)\le\psi(r-1)+(1/4+d(r)/2)2^r.
\end{equation}
Consequently,
\begin{equation}\label{eq:cycle-recursion-window}
 \frac{\psi(r)}{2^r}
 \le\frac12+
 \max_{\lceil r/2\rceil\le s\le r}d(s)+2^{-r/2}.
\end{equation}
In particular, $d(r)=o(1)$ implies $\psi(r)\le(1/2+o(1))2^r$.
\end{proposition}

\begin{proof}
Choose a connected graph $G$ of rank $r$ with $c(G)=\psi(r)$.
Delete leaves, and suppress degree-two vertices.
For $r\ge2$, the resulting connected multigraph $K$ has minimum degree at least three.
These operations preserve its rank and cycle count, and $|V(K)|\le2r-2$.

The elementary Moore bound gives a cycle $W$ of length at most $2\log(2r)+1$ in $K$.
Loops and parallel pairs already give cycles of length one or two.
Otherwise, the usual vertex-centered and edge-centered breadth-first trees give the claimed logarithmic bound.
Choose an edge $e$ of $W$.
Subdivide every edge of $K$ twice, obtaining a simple graph $K'$ of the same rank.
The cycle $W$ becomes a cycle $W'$ with
\[
 |E(W')|\le6\log(2r)+3\le\sqrt r
\]
for sufficiently large $r$.
Subdivision preserves outside-forest probability and cycle counts.
Lemma~\ref{lem:recursion-marked-edge} therefore yields
\[
 c_e(K)\le(1/2+d(r))2^{r-1}.
\]
Since $e$ lies on a cycle, $K-e$ is connected and has rank $r-1$.
The deletion identity gives
\[
 \psi(r)=c(K-e)+c_e(K)
 \le\psi(r-1)+(1/2+d(r))2^{r-1},
\]
which proves \eqref{eq:cycle-deletion-recurrence}.

Put $z_r=\psi(r)/2^r-1/2$.
The recurrence becomes $z_r\le z_{r-1}/2+d(r)/2$.
For $s=\lfloor r/2\rfloor$, iteration gives
\[
 z_r\le2^{-(r-s)}z_s+
 \sum_{j=s+1}^{r}2^{j-r-1}d(j).
\]
Use $z_s\le1/2$ and sum the geometric weights, whose sum is less than one.
This proves \eqref{eq:cycle-recursion-window}.
\end{proof}

The manuscript's forest error $20/\log\log\log r$ gives the same order of cycle-count error through this recurrence.
The constant $21$ also follows for sufficiently large $r$.
If a one-cycle forest estimate has error $O(r^{-a})$ for a fixed $a>0$, the cycle-count error has the same order.
No additional term from cycles avoiding a selected witness set is needed.

\subsection{The additive excess as an edge problem}

For $r\ge2$, define
\[
 X_r=\psi(r)-2^{r-1},\qquad
 g(r)=\max_{\beta(G)=r}\ \min_{e\text{ on a cycle}}c_e(G)-2^{r-2}.
\]
The maximum may be taken over connected simple graphs.
It exists because the minimum edge frequency is an integer between one and $2^{r-1}$.

\begin{proposition}\label{prop:cycle-edge-excess-recurrence}
For every $r\ge2$,
\[
 X_r-X_{r-1}\le g(r).
\]
Hence, for $1\le s<r$,
\[
 X_r\le X_s+\sum_{j=s+1}^r g(j).
\]
\end{proposition}

\begin{proof}
Choose a graph attaining $\psi(r)$ and an edge $e$ of minimum frequency among its cyclic edges.
Then
\[
 \psi(r)=c(G-e)+c_e(G)
 \le\psi(r-1)+2^{r-2}+g(r).
\]
Subtract $2^{r-1}$ and iterate.
\end{proof}

For a simple graph and an edge $e=uv$, one has $c_e(G)=P_{uv}(G-e)$.
Thus an upper bound on $g(r)$ is a terminal-path bound for one suitable edge of each graph.
If $g(r)=O(2^r/r^a)$, then $X_r=O(2^r/r^a)$.
If $g(r)=O(2^{(1-\kappa)r})$ for $0<\kappa<1$, then $X_r$ has that same exponential order.
Both conclusions follow by summing the recurrence, splitting the sum at $r/2$ when needed.

The general additive stability theorem gives a fixed-gap bound, with a threshold depending on that gap.
It does not presently control this threshold strongly enough to give a polynomial error in arbitrary graphs.
Accordingly, the recurrence does not by itself prove such an error bound.

\subsection{An alternative proof of the fixed-minor error bound}

The recurrence also gives a proof of Theorem~\ref{thm:minor-cycle-error} without adding witness triangles.
We record it because it separates the geometric input from the counting step.

Fix $h\ge4$, put $K=4096h^3$, and set
\[
 A_h=8K(3K-1)2^{3K-2}.
\]
Let $\psi_h(r)$ be the maximum cycle count among simple $K_h$-minor-free graphs of rank $r$.
Suppose an extremizer has more than $2^{r-1}$ cycles.
The proof of Theorem~\ref{thm:minor-cycle-error} reduces it to a unique cyclic kernel $\Gamma$ with
\[
 3\le\deg(v)\le5,\qquad |V(\Gamma)|\le2r-2.
\]
Its underlying simple graph remains $K_h$-minor-free.

Choose a shortest cycle $W$ in $\Gamma$.
For $r\ge512$, the Moore bound gives
\[
 |V(W)|\le|E(W)|\le2\log(2r)+1\le r/16.
\]
The graph $J=\Gamma-V(W)$ therefore has rank at least $r-4|V(W)|$.
Apply Lemma~\ref{lem:path-small-components}, with the same $K$ as above, to delete at most $|V(J)|/16$ vertices.
Every remaining component has at most $K$ vertices, while the remaining rank is at least
\[
 r-4|V(W)|-4|V(J)|/16\ge r/2-4|V(W)|\ge r/4.
\]
Each component has rank at most $2K$.
Selecting one cycle from every cyclic component gives at least $r/(8K)$ disjoint cyclic regions outside $W$.
Their boundaries in $\Gamma$ are at most $3K$.
The sharper small-cut estimate \eqref{eq:path-small-cuts} gives
\[
 p(\Gamma,W)\le1/2+A_h/r.
\]

For any edge $e$ of $W$, Lemma~\ref{lem:recursion-marked-edge} now yields
\[
 c_e(\Gamma)\le(1/2+A_h/r)2^{r-1}.
\]
Choose an edge of the corresponding subdivided path in the original simple graph.
Deleting that edge removes exactly the cycles counted by $c_e(\Gamma)$ and decreases rank by one.
It preserves the excluded-minor condition.
Hence, whenever $\psi_h(r)>2^{r-1}$ and $r\ge512$,
\[
 \psi_h(r)\le\psi_h(r-1)+(1/2+A_h/r)2^{r-1}.
\]

For completeness, this recurrence gives the same explicit constant as Theorem~\ref{thm:minor-cycle-error}.
Put $C=A_h+1$.
For $r\le2C$, the trivial bound $\psi_h(r)<2^r$ gives
$\psi_h(r)\le(1/2+C/r)2^r$.
For larger $r$, assume the same bound at rank $r-1$.
If $\psi_h(r)\le2^{r-1}$, there is nothing to prove.
Otherwise the recurrence applies, and its normalized excess is at most
\[
 \frac{C}{2(r-1)}+\frac{A_h}{2r}\le\frac Cr.
\]
The last inequality is equivalent to $C/(r-1)\le1$, which holds in this range.
This proves the error bound by induction, without any witness addition.

\section{A quantitative planar cycle-count consequence}
\label{sec:planar-cycle-error}

Aldred and Thomassen proved the planar asymptotic bound
$c(G)\le(1/2+o(1))2^{\beta(G)}$ in 2008.\footnote{R. E. L. Aldred and C. Thomassen,
\emph{On the maximum number of cycles in a planar graph}, Journal of Graph Theory 57 (2008), 255--264,
\url{https://doi.org/10.1002/jgt.20290}. Their institutional publication record is
\url{https://orbit.dtu.dk/en/publications/on-the-maximum-number-of-cycles-in-a-planar-graph/}.}
The argument below gives an explicit error term from the present machinery.
The original article's full proof was not available during this comparison.
We therefore make no priority claim for the quantitative rate.

\begin{proposition}[An explicit planar error]\label{prop:planar-cycle-error}
Every planar graph of rank $r\ge1$ satisfies
\[
 c(G)\le2^{r-1}+32257\frac{2^r}{r}.
\]
More generally, for graphs of orientable genus $g=o(r)$,
\[
 c(G)\le2^{r-1}+O(2^r/r),
\]
with an absolute implied constant for sufficiently large $r$.
\end{proposition}

\begin{lemma}[Embedded marked sets]\label{lem:embedded-marked-forest}
Let $H$ be a connected graph of rank $R$, embedded in an orientable surface of genus $g$.
Let $W$ be a nonempty union of cycles with $m$ edges.
For every integer $L\ge6$ such that
\[
 B=(L-5)R-(4L+12)m-2(L+1)g>0,
\]
one has
\[
 p(H,W)\le\frac12+
 \frac{3L(L+1)(L-1)2^{L-2}}{B}.
\]
\end{lemma}

\begin{proof}
Inside each vertex disk, replace a high-degree vertex by a trivalent tree respecting the external cyclic order.
Mark every new vertex incident with an edge of $W$.
Retain the original vertex when its degree is at most three, marking it if necessary.
The marked set $P$ has size $p\le2m$.
No connectivity condition on the marked set is needed.

Contracting the replacement trees identifies the two cycle spaces.
If the original even subgraph is a forest outside $W$, its lift is a forest after deleting $P$.
Indeed, a forest remains a forest when its vertices are replaced by trees.
Delete leaves and suppress degree-two vertices outside $P$.
These operations preserve the cycle space and the required forest implication.
They also preserve the embedding.
Call the resulting connected multigraph $\Gamma$.
Every unmarked vertex has degree three, while each marked vertex has degree at least one and at most three.
Consequently,
\[
 |V(\Gamma)|\le2R+2p-2,
 \qquad |E(\Gamma)|\le3R+2p-3.
\]

Put $J=\Gamma-P$.
Deleting $P$ removes at most $3p$ edges, giving $\beta(J)\ge R-2p$ when $J$ is nonempty.
When $J$ is empty, subcubicity gives $R\le p/2+1\le2p$, so the same bound holds.
The face argument of Lemma~\ref{lem:path-face-packing} now gives at least
\[
 R-2p-2g-\frac{6R+4p}{L+1}
 =\frac{(L-5)R-(2L+6)p-2(L+1)g}{L+1}
\]
retained faces with boundary length at most $L$.
Greedy selection yields $k$ vertex-disjoint cycles outside $P$, where
\[
 k\ge\frac{(L-5)R-(2L+6)p-2(L+1)g}{3L(L+1)}
 \ge\frac{B}{3L(L+1)}.
\]
Each cycle's vertex set has boundary at most $L$ in $\Gamma$.
Apply \eqref{eq:path-small-cuts} and the forest implication to obtain the stated bound.
\end{proof}

\begin{proof}[Proof of Proposition~\ref{prop:planar-cycle-error}]
Join different components by bridges, preserving the rank, cycles, and genus bound.
For $r\ge2$, put $k=\lceil\log_2r\rceil$.
Choose $k$ edges $D$ outside a spanning tree.
Exactly $2^{r-k}$ even subgraphs avoid $D$.

For each $e=uv\in D$, add a fresh vertex $x_e$ and the edges $ux_e,x_ev$.
Embed this new path beside $e$ in a narrow neighborhood of that edge.
The resulting graph $H$ has genus at most $g$ and rank $R=r+k$.
Let $W$ be the union of the $k$ resulting triangles, so $|E(W)|=3k$.

For each original cycle $C$ meeting $D$ and each subset of the new triangles, take their symmetric difference with $C$.
These $2^k$ even subgraphs have the same forest outside $W$.
The new edges determine the chosen subset and hence determine $C$.
Thus at most $p(H,W)2^r$ original cycles meet $D$.
The cycles avoiding $D$ contribute at most $2^{r-k}$.

Set $L=7$ in Lemma~\ref{lem:embedded-marked-forest}.
When the denominator is positive, it gives
\begin{equation}\label{eq:planar-error-denominator}
 \frac{c(G)}{2^r}
 \le\frac12+
 \frac{32256}{2r-118\lceil\log_2r\rceil-16g}
 +2^{-\lceil\log_2r\rceil}.
\end{equation}
If $g=0$ and $r\ge2048$, then $r\ge118\lceil\log_2r\rceil$.
The last display is therefore at most $1/2+32257/r$.
For $1\le r<2048$, the trivial bound $c(G)<2^r$ proves the same assertion.
If $g=o(r)$, the denominator in \eqref{eq:planar-error-denominator} is $(2-o(1))r$.
This proves the remaining assertion.
\end{proof}

The constants come from a fixed face-length cutoff and are not optimized.
This proof uses only embedding-preserving additions beside selected edges.

\section{A polynomial cycle-count error with a fixed excluded minor}
\label{sec:minor-cycle-error}

The same order of error holds throughout every fixed proper minor-closed class.
Here the separators are applied before the witness paths are added.
Thus the proof does not need to preserve minor exclusion under those additions.

\begin{theorem}[Cycle-count error with an excluded minor]\label{thm:minor-cycle-error}
For each fixed integer $h\ge4$, there is a constant $C_h$ such that every simple graph $G$
with no $K_h$ minor and rank $r\ge1$ satisfies
\[
 c(G)\le2^{r-1}+C_h\frac{2^r}{r}.
\]
One may take
\[
 K=4096h^3,\qquad C_h=1+8K(3K-1)2^{3K-2}.
\]
In particular, this bound holds in every fixed proper minor-closed class, with a class-dependent constant.
\end{theorem}

\begin{proof}
Suppose first that $c(G)>2^{r-1}$.
All cycles then lie in one block.
Indeed, two nonempty groups of cyclic blocks, with ranks $a,b\ge1$, would give
\[
 c(G)\le(2^a-1)+(2^b-1)\le2^{a+b-1}=2^{r-1}.
\]
Discard all other blocks.
This preserves the rank and cycle count, giving a two-connected graph $B$.

For each vertex $v$ of $B$, the graph $B-v$ is connected.
The incident pattern of a uniform even subgraph is therefore uniform over all even subsets of the incident edges.
A cycle uses zero or two edges at $v$.
If $v$ has degree $d$, this gives
\[
 \frac{c(B)}{2^r}\le\frac{1+\binom d2}{2^{d-1}}.
\]
The right side is at most one half when $d\ge6$.
Consequently, every vertex of $B$ has degree at most five.

It suffices to consider $r\ge512$.
Suppress all degree-two vertices of $B$, retaining parallel edges.
Since $r\ge2$, the resulting loopless multigraph $\Gamma$ has minimum degree three and maximum degree five.
Its rank and cycle count remain unchanged.
Its underlying simple graph is a minor of $G$, and
\[
 |V(\Gamma)|\le2r-2.
\]

Put $k=\lceil\log_2r\rceil$, and choose $k$ edges $D$ outside a spanning tree of $\Gamma$.
Let $P_0$ be their endpoints, so $|P_0|\le2k$.
Deleting a vertex of degree at most five decreases rank by at most four.
Hence the graph $J=\Gamma-P_0$ has
\[
 \beta(J)\ge r-8k.
\]
Apply Lemma~\ref{lem:path-small-components} to $J$, with $K=4096h^3$.
There is a set $S$ of size at most $|V(J)|/16$ such that every component of $J-S$ has at most $K$ vertices.
Another rank estimate gives
\[
 \beta(J-S)\ge r-8k-4|S|
 \ge r/2-8k\ge r/4.
\]
The last inequality uses $r\ge32\lceil\log_2r\rceil$, which holds for $r\ge512$.

Every component of $J-S$ has rank at most $2K$.
Choose one cycle from each cyclic component.
This yields $q\ge r/(8K)$ vertex-disjoint cycles, each using at most $K$ vertices.
Each chosen vertex set has edge boundary at most $3K$ in $\Gamma$.

Now add one fresh two-edge path parallel to each selected edge in $D$.
Let $H$ be the resulting graph, and let $W$ be the union of the resulting triangles.
The added edges only meet $\Gamma$ at vertices in $P_0$.
Thus the selected cyclic regions retain their boundary bound in $H$ and avoid $V(W)$.
The small-cut estimate gives
\[
 p(H,W)\le\frac12+\frac{(3K-1)2^{3K-2}}{q}
 \le\frac12+\frac{8K(3K-1)2^{3K-2}}{r}.
\]
The triangle-flip injection from Proposition~\ref{prop:planar-cycle-error} bounds cycles meeting $D$ by $p(H,W)2^r$.
At most $2^{r-k}\le2^r/r$ cycles avoid $D$.
Together these estimates give the stated constant $C_h$.

If $c(G)\le2^{r-1}$, the result holds immediately.
For $1\le r<512$, use $c(G)<2^r$ and $C_h\ge256$.
Finally, every proper minor-closed class excludes some $K_h$ with fixed $h\ge4$.
\end{proof}

This result is a consequence of the present counting and small-cut arguments together with the classical separator theorem.
Its publication priority has not been checked.

% Research fragment. Include in a document with amsmath, amssymb, amsthm,
% hyperref, and theorem environments. This file does not change the main paper.
\section{The second-order scale in a sparse random host}
\label{sec:research-random-hosts}

\paragraph{Status.}
The lower bound below has a complete first-moment proof.
The upper bound is a parameter refinement of the published Alon--Krivelevich construction.
Its proof uses their expansion and Hamilton-path inputs.
Neither result is formalised in Lean.
We have not established publication priority for these refinements.

All logarithms in this section are natural, except $\log_2$.
Put $D=np$.
For a pancyclic graph $G$ on $n$ vertices, define
\[
 \operatorname{Pex}(G)
 =\min\{|E(H)|-n:H\subseteq G\text{ is spanning and pancyclic}\}.
\]
If $G$ is not pancyclic, set this quantity to $+\infty$.

\begin{theorem}[Second-order bounds in the Alon--Krivelevich range]
\label{thm:rhost-scale}
Define
\[
 p_1=p_5=\frac{2\ln\ln n}{n},\qquad
 p_2=p_3=\frac{50\ln n}{n\ln\ln n},\qquad
 p_4=\frac{\ln n+10\sqrt{\ln n}}{n},
 \qquad p_* =1-\prod_{i=1}^{5}(1-p_i).
\]
For $G\sim G(n,p_*)$, with probability tending to one,
\[
 \log_2 n+\left(\frac23-o(1)\right)\frac{\ln n}{\ln\ln n}
 \ \le\ \operatorname{Pex}(G)
 \ \le\ \log_2 n+(2+o(1))\frac{\ln n}{\ln\ln n}.
\]
Consequently,
\[
 \operatorname{Pex}(G)-\log_2 n
 =\Theta\!\left(\frac{\ln n}{\ln\ln n}\right)
 \qquad\text{with high probability.}
\]
The upper bound also holds for every $p\ge p_*$.
The lower bound holds whenever $np\sim\ln n$, including values arbitrarily close above the pancyclicity threshold.
\end{theorem}

The last lower-bound assertion does not supply a finite upper bound below $p_*$.
Thus this theorem does not settle the full-threshold question.

\subsection{A uniform count of small subgraphs}

\begin{lemma}[Kernel count]
\label{lem:rhost-kernel}
Let $G\sim G(n,D/n)$, with $D\ge2$.
Let $S\ge2$ be an integer, and let $R>0$.
The probability that $G$ contains a connected subgraph $F$ satisfying
\[
 \delta(F)\ge2,\qquad
 q:=|E(F)|-|V(F)|\ge1,\qquad
 |E(F)|\le\min\{S,Rq\}
\]
is at most
\[
 S\sum_{q=1}^{S}
 \left(\frac{(2S)^{10}D^R}{n}\right)^q.
\]
\end{lemma}

\begin{proof}
Write $e=|E(F)|$.
Suppress each maximal path whose internal vertices have degree two.
The resulting connected multigraph has $t$ vertices and $t+q$ edges.
Every vertex has degree at least three.
The degree identity gives $1\le t\le2q$, and hence $t+q\le3q$.

Label the kernel vertices and order its edges.
There are at most $(2q)^{6q}$ lists of edge endpoints.
There are at most $(S+1)^{3q}$ lists of positive path lengths.
There are at most $n^{e-q}$ assignments of distinct host vertices to a fixed subdivision.
These bounds can count invalid or repeated descriptions, which only increases the upper bound.

Each described simple subgraph has occurrence probability $(D/n)^e$.
After summing over $t$, its expected count is at most
\[
 2q(2q)^{6q}(S+1)^{3q}n^{-q}D^e
 \le (2S)^{10q}n^{-q}D^{Rq}.
\]
There are at most $S$ possible values of $e$ for each $q$.
Sum these bounds and apply Markov's inequality.
\end{proof}

\begin{lemma}[Independence of a small family of short cycles]
\label{lem:rhost-independence}
Let $G\sim G(n,D/n)$, with $D\ge2$, and put
\[
 L=\left\lfloor
 \frac{2(\ln n-40\ln\ln n)}{3\ln D}
 \right\rfloor.
\]
Assume $L\ge3$.
With probability $1-O((\ln n)^{-17})$, every family of at most $L$ distinct cycles,
each of length at most $L$, is linearly independent over $\mathbb F_2$.
\end{lemma}

\begin{proof}
Set $S=\lceil(\ln n)^2\rceil$ and $R=3L/2$.
For sufficiently large $n$,
\[
 D^R\le\frac{n}{(\ln n)^{40}},\qquad
 (2S)^{10}\le(\ln n)^{21}.
\]
Lemma~\ref{lem:rhost-kernel} shows that the probability of a forbidden subgraph is
\[
 O\!\left((\ln n)^2(\ln n)^{-19}\right)
 =O((\ln n)^{-17}).
\]

Suppose a family of at most $L$ short cycles is dependent.
Choose a minimal dependent subfamily $C_1,\ldots,C_k$, and let $F$ be its union.
Then $3\le k\le L$.
Minimality gives $\dim\langle C_1,\ldots,C_k\rangle=k-1$.
Each edge occurs an even positive number of times in the family.
Thus, with $e=|E(F)|$,
\[
 e\le\frac12\sum_{i=1}^k|C_i|\le\frac{kL}{2}\le S.
\]
The union is connected.
Otherwise, each component would give a proper dependent subfamily.
Also, $\delta(F)\ge2$ and
\[
 q=e-|V(F)|=\beta(F)-1\ge k-2\ge1.
\]
Since $k\le3(k-2)$ for $k\ge3$,
\[
 e\le\frac{kL}{2}\le\frac{3L}{2}(k-2)\le Rq.
\]
This contradicts the absence of a forbidden subgraph.
\end{proof}

The lemma concerns families of at most $L$ cycles.
It does not assert independence of the entire family of short cycles.
This distinction avoids an unjustified bound on a minimal dependency's size.

\subsection{The improved lower bound}

\begin{proposition}[Finite lower bound]
\label{prop:rhost-lower}
Under the hypotheses of Lemma~\ref{lem:rhost-independence}, with the same probability,
every spanning pancyclic subgraph $H$ of $G$ satisfies
\[
 |E(H)|-n
 \ge \log_2(n-2)+L-3
 -\log_2\!\left(\frac{L(L+1)}2-3\right).
\]
In particular, if $D\sim\ln n$, then with high probability
\[
 \operatorname{Pex}(G)
 \ge\log_2 n+\frac{2\ln n}{3\ln D}-O(\ln\ln n).
\]
\end{proposition}

\begin{proof}
Choose in $H$ one cycle of each length $3,\ldots,L$.
Let $W$ be their union, and write $r=\beta(H)$, $t=\beta(W)$, and $M=|E(W)|$.
Lemma~\ref{lem:rhost-independence} gives
\[
 t\ge L-2,\qquad M\le\frac{L(L+1)}2-3.
\]
The outside-part length count gives
\[
 n-2\le M2^{r-t}.
\]
The Hamilton cycle makes $H$ connected, so $r=|E(H)|-n+1$.
Taking binary logarithms proves the finite bound.
For $D\sim\ln n$, the rounding and the $40\ln\ln n$ margin change $L$ by only $O(1)$.
Also, $\log_2 M=O(\ln\ln n)$.
\end{proof}

The denominator $\ln D$ matters for a quantitative error term.
The assumption $D\sim\ln n$ permits replacement by $\ln\ln n$ with a relative $o(1)$ error.
It does not always permit replacement with only an additive $O(\ln\ln n)$ error.

The coefficient $2/3$ improves the disjoint-cycle coefficient $1/2$ in the proposed research list.
The proof requires independence, rather than vertex-disjointness.
A theta graph explains the first obstruction to this argument.
Its three cycles are dependent, and their total length is twice its number of edges.
At lengths near $\frac23\log_D n$, its occurrence threshold becomes relevant.
This observation limits the independence method, not the true pancyclicity excess.

\subsection{A two-digit refinement of the upper construction}

The primary source is Alon and Krivelevich,
\emph{Sparse Pancyclic Subgraphs of Random Graphs},
SIAM Journal on Discrete Mathematics 39 (2025), 562--574,
\href{https://doi.org/10.1137/23M1598969}{doi:10.1137/23M1598969}.
We checked the authors' \href{https://www.math.tau.ac.il/~krivelev/pancyclic-rg.pdf}{29 November 2024 manuscript}.
Their Theorem~1 gives the threshold $p_*$ used above.
The following proof changes the short-length encoding parameters in Sections~3--5.

\begin{lemma}[Corrected robust expansion statement]
\label{lem:rhost-expansion}
Set
\[
 \beta=\frac{2(\ln\ln n)^2}{\ln n},\qquad
 q_0=\frac{50\ln n}{n\ln\ln n},\qquad
 L_0=\Theta(n/\sqrt{\ln n}).
\]
With high probability, every vertex set $U$ of size at least $n-2L_0$ in $J\sim G(n,q_0)$ contains a set $U'$ satisfying
\[
 |U'|>|U|-\beta n,\qquad
 |N_{J[U']}(A)|\ge\frac{|A|}{3\beta}
 \quad(0<|A|\le\beta n).
\]
\end{lemma}

\begin{proof}
Let $s=\lceil\beta n\rceil$.
With high probability, every two disjoint $s$-vertex sets have an edge between them.
Indeed, the failure probability is at most
\[
 \binom ns^2e^{-q_0s^2}
 \le\exp\!\left(s\left(2\ln(en/s)-q_0s\right)\right)=o(1),
\]
because $q_0s\sim100\ln\ln n$ and $\ln(en/s)\sim\ln\ln n$.
Assume this property holds.

Start with $U$ and repeatedly delete a nonempty set $A$ that violates the required expansion.
Each deleted set has at most $\beta n$ vertices.
Suppose the accumulated deleted set $B$ first reaches size at least $\beta n$.
Then $\beta n\le|B|<2\beta n$.
In the original induced graph $J[U]$, its external neighbourhood has size less than
\[
 \frac{|B|}{3\beta}<\frac{2n}{3}.
\]
Thus $U\setminus(B\cup N_{J[U]}(B))$ has at least
$n/3-2L_0-2\beta n$ vertices.
For large $n$, this number exceeds $s$.
Choosing $s$ vertices there and $s$ vertices in $B$ contradicts the assumed edge property.
Therefore fewer than $\beta n$ vertices are deleted before the procedure stops.
The remaining set is the required $U'$.
\end{proof}

This corrects the size assertion in the source's Claim~4.3.
Its printed bound $|U'|\ge(1-\beta)n$ can exceed $|U|$ when $|U|=n-2L_0$.
The corrected bound gives $|U'|=n-o(n)$, which suffices for every exposure estimate below.
It also supplies the same rooted-tree embedding conditions.
For $k=\lfloor\beta n/2\rfloor$, take maximum degree $\lfloor1/(4\beta)\rfloor$ and tree size at most $\lfloor n/50\rfloor$.
For $|A|\le k$, the expansion surplus is at least $|A|/(12\beta)>1$.
For $k<|A|\le2k$, that surplus exceeds $n/50$ for sufficiently large $n$.
These are the two hypotheses of the source's Theorem~2.2.

\begin{proposition}[Refined upper bound]
\label{prop:rhost-upper}
Let $p\ge p_*$, with $p_*$ as in Theorem~\ref{thm:rhost-scale}.
Then, with high probability,
\[
 \operatorname{Pex}(G(n,p))
 \le \log_2 n+(2+o(1))\frac{\ln n}{\ln\ln n}.
\]
\end{proposition}

\begin{proof}
Use the five independent exposures from the cited construction.
In $G_2$, first expose only edges with both endpoints outside the initial graph $H_1$.
In $G_3$, first expose only edges with both endpoints outside $H_2$.
The two-set estimate in Lemma~\ref{lem:rhost-expansion} applies within either reservoir.
All path endpoints lie outside its corresponding reservoir.
Thus endpoint edges remain independent until their designated queries.
Keep its parameters
\[
 \beta=\frac{2(\ln\ln n)^2}{\ln n},\qquad
 K_0=\left\lfloor\log_2\frac{\ln n}{6\ln\ln\ln n}\right\rfloor,
 \qquad K=\left\lfloor\log_2\frac{n}{\sqrt{\ln n}}\right\rfloor,
 \qquad L_0=2^{K+1}-1.
\]
For integral branching, put
\[
 a=\left\lfloor(5\beta)^{-1}\right\rfloor,
 \qquad d=\left\lfloor\log_a(n/200)\right\rfloor.
\]
Then $d=(1+o(1))\ln n/\ln\ln n$.
The change from real to integral branching has no effect on this estimate.

Replace the short-length parameters by
\[
 t=2,\qquad b=\lceil\sqrt{2\ln n}\rceil.
\]
The short base cycle has length $2b+1$.
Attach internally disjoint paths of lengths
\[
 d+2+j\quad\text{and}\quad d+2+jb
 \qquad(0\le j<b)
\]
to its $2b$ designated edges, as in the cited construction.
Choosing one path from each group produces every length in
\[
 [T+1,T+b^2],\qquad T=2(d+b+1).
\]
Indeed, each integer $0\le k<b^2$ has the form $k=j_0+bj_1$.
The resulting cycle has length $T+1+k$.

Here are the estimates needed to transfer the source's existence proof.

\emph{Initial cycles.}
The new length $2b+1$ is $o(\ln n/\ln\ln\ln n)$.
It is therefore below the longest initial cycle used in Lemma~4.1.
That lemma still supplies the required vertex-disjoint initial cycles.

\emph{Vertex budgets.}
The short gadget uses at most
\[
 O(bd+b^3)=O((\ln n)^{3/2})
\]
vertices.
This is $o(L_0)$.
The binary gadget, connectors, and short gadget together use at most $2L_0$ vertices.
Thus Lemma~\ref{lem:rhost-expansion} applies at every step.
Each remaining expander contains $n-o(n)$ vertices.

\emph{Path embeddings.}
Each additional stem has length at most $b(b-1)=O(\ln n)$.
The rooted trees in Lemma~4.5 still have fewer than $n/50$ vertices.
Their leaf sets have at least $\beta n/40$ vertices because
\[
 a^d>\frac{n}{200a}\ge\frac{\beta n}{40}.
\]
Process the short paths in their order around the base cycle.
For the new endpoint, expose at most $h=\lceil\beta n/100\rceil$ candidate neighbours in the remaining expander.
These edges have not appeared in any earlier exposure.
The first-neighbour failure probability is at most
\[
 \exp(-p_3h)\le(\ln n)^{-1}.
\]
The other endpoint has previously exposed at most $h$ edges.
At least $\beta n/40-h\ge\beta n/80$ leaves therefore have unexposed endpoint edges.
The remaining-leaf failure probability is at most
\[
 \exp(-p_3\beta n/80)=(\ln n)^{-5/4}.
\]
There are $2b=O(\sqrt{\ln n})$ short paths.
The sum of their failure probabilities tends to zero.
The binary paths and the final connector retain their original estimates.
The two-tree embeddings also remain below the size limit.
For $a\ge3$, their two complete trees use at most $3n/200$ vertices together.
The connecting stem has $O(L_0)=o(n)$ vertices.
Their maximum degree is at most $a+1\le\lfloor1/(4\beta)\rfloor$.

\emph{Length coverage.}
The binary gadget starts at
\[
 \ell_*=(K_0+1)(d+3)=(1+o(1))\log_2 n.
\]
There are enough base edges for all designated binary detours and the reserved connector edge.
Indeed, with $x=\ln\ln n$,
\[
 \ln a=x-2\ln x-\ln10+o(1),\qquad
 K_0+1>\frac{x-\ln6-\ln\ln x}{\ln2}.
\]
These inequalities give $\ell_*-K\to\infty$.
Since $b^2\ge2\ln n$, the short interval reaches $\ell_*$ for all sufficiently large $n$.
The remaining interval-covering shortcuts from Lemma~4.8 are unchanged.
Only the lengths $3,\ldots,T$ need individual short chords.
Here $T=O(\ln n/\ln\ln n)$.
Each missing-chord probability is at most $(1-p_5)^n\le(\ln n)^{-2}$.
Their sum tends to zero, as required by the argument of Lemma~4.9.

\emph{Hamilton completion.}
The remaining vertex set still has size at least $n-2L_0$.
Claim~4.7 therefore applies without a parameter change.
The number of long interval-covering chords is at most
\[
 m=\lfloor n2^{-K}\rfloor=O(\sqrt{\ln n}).
\]
The final graph has a Hamilton cycle and at most
\[
 K+bt+m+T+O(1)
 =K+2d+4b+m+O(1)
 \le\log_2 n+2d+O(\sqrt{\ln n})
\]
additional edges.
Since $\sqrt{\ln n}=o(\ln n/\ln\ln n)$, this proves the claimed upper bound.
The conclusion for $p\ge p_*$ follows by monotonicity.
\end{proof}

\paragraph{Further notation corrections in the checked source.}
These corrections do not require a new probabilistic input.
First, the source writes $b^t\ge\ln n>\ell_*$, although $\ell_*\sim\log_2 n$.
The choice $b^2\ge2\ln n$ gives the required overlap explicitly.
Second, Lemma~4.4 states one extra edge in the binary replacement path.
For a desired path length $\ell_i=2^i+1$, use a connecting stem of length $\ell_i-2d-2$.
Its positivity follows from $i>K_0$.
This gives the required gain $2^i$ when one cycle edge is replaced.
The short-chord argument above covers every length through $T$.
This also supplies the extra $t$ lengths omitted from the printed range in Lemma~4.9.

The exposure budget in Lemma~4.5 also needs correction.
The source first permits $n/\ln\ln n$ exposed edges at an endpoint.
This exceeds the guaranteed $\beta n/40$ leaves, so its later fresh-leaf subtraction does not follow.
The budget $h=\lceil\beta n/100\rceil$ repairs that step and supports the new union bound.

\paragraph{What remains open in this note.}
The lower and upper constants do not match.
The upper construction still needs $p\ge p_*$.
This fixed exposure scheme does not reach the full pancyclicity threshold range.
The later adaptive construction addresses that separate question.
The lower proof uses the outside-part length count, without the stronger forest probability estimate.


% Phase-two research fragment. Uses the theorem environments and packages
% of research-progress.tex. The first-moment kernel estimate is restated here.
\section{A coefficient-one random-host lower bound}
\label{sec:rhost-phase-two}

\paragraph{Status.}
This section improves the lower coefficient from $2/3$ to $1$.
It uses a weighted pseudoforest argument and the same random-host kernel count.
The proof is not formalised in Lean.
We have not established publication priority.

All logarithms are natural, except $\log_2$.
A pseudoforest is a graph whose components contain at most one cycle each.
All graphs in this section are finite and simple.
Isolated vertices may be retained in every spanning subgraph.

\begin{lemma}[Cycles and pseudoforests]
\label{lem:rhost-pseudoforest-cycles}
Let $W$ be a graph, let $P\subseteq W$ be a spanning pseudoforest,
and let $\mathcal C$ be a family of $k$ distinct cycles of $W$.
Then
\[
 \sum_{C\in\mathcal C}|E(C)\setminus E(P)|
 \ge 2\bigl(k-\beta(W)\bigr).
\]
\end{lemma}

\begin{proof}
Let $u$ be the number of unicyclic components of $P$.
Let $a$ count omitted edges whose endpoints lie in the same component of $P$.
Let $b$ count the remaining omitted edges.
If $c_P$ and $c_W$ are the component counts, then $b\ge c_P-c_W$.
Consequently,
\[
 \beta(W)=u+a+b-c_P+c_W\ge u+a.
\]

Write $n_j$ for the number of selected cycles using exactly $j$ omitted edges.
A unicyclic component has only one cycle, so $n_0\le u$.
A tree has at most one simple path between two fixed vertices.
A unicyclic graph has at most two.
Thus each of the $a$ internal omitted edges belongs to at most two cycles counted by $n_1$.
An omitted edge between different components belongs to none.
It follows that
\[
 2n_0+n_1\le2u+2a\le2\beta(W).
\]
Therefore
\[
 \sum_{C\in\mathcal C}|E(C)\setminus E(P)|
 =\sum_{j\ge0}jn_j
 \ge2k-2n_0-n_1
 \ge2(k-\beta(W)).
\]
\end{proof}

\begin{lemma}[Balanced pseudoforest cover]
\label{lem:rhost-balanced-cover}
Let $R\ge2$ be an integer.
Suppose every nonempty edge set $A\subseteq E(W)$ satisfies
\[
 (R-1)|A|\le R|V(A)|,
\]
where $V(A)$ is the set of endpoints of edges in $A$.
Then there are $R$ spanning pseudoforests $P_1,\ldots,P_R$ of $W$
such that each edge of $W$ is absent from exactly one $P_i$.
\end{lemma}

\begin{proof}
Form a bipartite graph.
Its left side has $R-1$ copies of every edge of $W$.
Its right side has $R$ copies of every vertex.
Each edge-copy is adjacent to all copies of both endpoints.
For a set $X$ of left vertices, let $A$ be the original edges represented in $X$.
Then
\[
 |X|\le(R-1)|A|\le R|V(A)|=|N(X)|.
\]
Hall's theorem gives a matching that covers every left vertex.

Forget the copy labels after recording the matched incidences.
This gives a bipartite multigraph with sides $E(W)$ and $V(W)$.
Its edge-vertices have degree $R-1$, and its original vertices have degree at most $R$.
Colour its edges properly with $R$ colours.
For completeness, such a colouring follows again from Hall's theorem.
Add dummy vertices and edges to make both sides equally large and every degree equal to $R$.
A regular bipartite multigraph has a perfect matching by Hall's theorem.
Remove one perfect matching at a time to obtain the $R$ colours.

For colour $i$, retain every original edge incident with an incidence of that colour.
Each retained edge has an assigned endpoint, and no two retained edges have the same assigned endpoint.
Every component therefore has at most as many edges as vertices.
The retained graph is a pseudoforest.
The $R-1$ incidences at each edge-vertex have distinct colours.
Thus each original edge is absent from exactly one colour.
\end{proof}

\begin{proposition}[Weighted cycle-length bound]
\label{prop:rhost-weighted-cycles}
Under the hypotheses of Lemma~\ref{lem:rhost-balanced-cover}, every family
$\mathcal C$ of distinct cycles satisfies
\[
 \boxed{\quad
 \beta(W)\ge\sum_{C\in\mathcal C}
 \left(1-\frac{|E(C)|}{2R}\right).
 \quad}
\]

In particular, one may sum over all cycles of length less than $2R$.
\end{proposition}

\begin{proof}
Apply Lemma~\ref{lem:rhost-pseudoforest-cycles} to every $P_i$ from
Lemma~\ref{lem:rhost-balanced-cover}.
Every edge is omitted exactly once, so
\[
 \sum_{i=1}^{R}\sum_{C\in\mathcal C}|E(C)\setminus E(P_i)|
 =\sum_{C\in\mathcal C}|E(C)|.
\]
The same sum is at least $2R(|\mathcal C|-\beta(W))$.
Rearrange this inequality.
\end{proof}

This deterministic inequality allows dependencies between the selected cycles.
It therefore avoids the theta obstruction to their complete independence.

\begin{lemma}[Local sparsity of the random host]
\label{lem:rhost-phase-two-sparsity}
Let $G\sim G(n,D/n)$, with $D\ge2$.
Put
\[
 R=\left\lfloor\frac{\ln n-40\ln\ln n}{\ln D}\right\rfloor,
 \qquad S=\lceil8(\ln n)^2\rceil.
\]
Assume $R\ge3$.
With probability $1-O((\ln n)^{-17})$, every subgraph $W$ with at most $S$ edges
satisfies the hypothesis of Lemma~\ref{lem:rhost-balanced-cover}.
\end{lemma}

\begin{proof}
The kernel enumeration in Lemma~\ref{lem:rhost-kernel} applies with this $S$ and $R$.
It bounds the probability of a connected subgraph $F$ satisfying
\[
 \delta(F)\ge2,\qquad q=|E(F)|-|V(F)|\ge1,
 \qquad |E(F)|\le\min\{S,Rq\}
\]
by
\[
 S\sum_{q=1}^{S}
 \left(\frac{(2S)^{10}D^R}{n}\right)^q
 =O((\ln n)^{-17}).
\]
Indeed, $D^R\le n/(\ln n)^{40}$ and $(2S)^{10}\le(\ln n)^{21}$ for large $n$.
Assume no such $F$ occurs.

Consider a connected component of an edge set $A\subseteq E(W)$.
Write $e$ and $v$ for its edge and vertex counts.
If $e\le v$, then $(R-1)e\le Rv$ immediately.
If $q=e-v\ge1$, repeatedly remove vertices of degree one.
The remaining nonempty core is connected, has minimum degree at least two, and retains the same $q$.
If the original component had $e\le Rq$, then its core would be forbidden.
Thus $e>Rq$, which implies $(R-1)e<Rv$.
Sum these inequalities over the components of $A$.
\end{proof}

\begin{theorem}[Coefficient-one lower bound]
\label{thm:rhost-coefficient-one}
Let $G\sim G(n,D/n)$, and let $R$ be as in Lemma~\ref{lem:rhost-phase-two-sparsity}.
Assume $D\ge2$ and $R\ge3$.
With probability $1-O((\ln n)^{-17})$, every spanning pancyclic subgraph $H$ satisfies
\[
 \boxed{\quad
 |E(H)|-n
 \ge\log_2(n-2)+R-3-\log_2(2R^2-R-3).
 \quad}
\]
In particular, if $D\sim\ln n$, then with high probability
\[
 \operatorname{Pex}(G)
 \ge\log_2 n+
 \frac{\ln n}{\ln D}
 -2\log_2\!\left(\frac{\ln n}{\ln D}\right)-O(1).
\]
Consequently,
\[
 \operatorname{Pex}(G)
 \ge\log_2 n+(1-o(1))\frac{\ln n}{\ln\ln n}.
\]
\end{theorem}

\begin{proof}
Choose in $H$ one cycle of every length $3,\ldots,2R-1$.
Let $W$ be their union.
Its number of edges is at most
\[
 M=\sum_{j=3}^{2R-1}j=2R^2-R-3\le S.
\]
On the event in Lemma~\ref{lem:rhost-phase-two-sparsity},
Proposition~\ref{prop:rhost-weighted-cycles} gives
\[
 \beta(W)\ge(2R-3)-\frac{2R^2-R-3}{2R}
 =R-\frac52+\frac{3}{2R}.
\]
Since $\beta(W)$ is an integer, $\beta(W)\ge R-2$.

Put $r=\beta(H)$.
The outside-part length count yields
\[
 n-2\le |E(W)|2^{r-\beta(W)}\le M2^{r-R+2}.
\]
The graph $H$ is connected, so $|E(H)|-n=r-1$.
Taking binary logarithms proves the finite bound.
If $D\sim\ln n$, then $R=\ln n/\ln D-O(1)$.
Also, $\log_2(2R^2-R-3)=2\log_2(\ln n/\ln D)+O(1)$.
\end{proof}

At the explicit probability $p_*$ from Theorem~\ref{thm:rhost-scale},
this improves the proved interval to
\[
 \log_2n+(1-o(1))\frac{\ln n}{\ln\ln n}
 \le\operatorname{Pex}(G(n,p_*))
 \le\log_2n+(2+o(1))\frac{\ln n}{\ln\ln n}.
\]
The upper bound and its probability range remain unchanged.

\paragraph{Scope of the improvement.}
The weighted inequality is deterministic and has no independence hypothesis on the cycles.
The random input only ensures the local edge-density condition.
The outside-part count then adds the rank cost of the short lengths to $\log_2n$.
This proof does not determine the sharp correction constant or extend the upper bound below $p_*$.

\subsection{Replacing individual short cycles by theta blocks}
\label{subsec:rhost-theta-upper}

\begin{lemma}[A theta block with prescribed lengths]
\label{lem:rhost-theta-block}
Let $s,x,y$ be integers at least three, with $x<y$.
Suppose
\[
 s+x+y\equiv0\pmod2,\qquad y-x<s<x+y.
\]
Set
\[
 \alpha=\frac{s+x-y}{2},\qquad
 \gamma=\frac{s+y-x}{2},\qquad
 z=\frac{x+y-s}{2}.
\]
An $s$-cycle with an additional internally disjoint path of length $z$
between two vertices at cyclic distance $\alpha$ contains cycles of lengths $s,x,y$.
Its rank is two.
It has a spanning path whose endpoints can be specified before the added path is embedded.
\end{lemma}

\begin{proof}
The three internally disjoint paths between the branch vertices have lengths
$\alpha,\gamma,z$.
Their pairwise sums are $s,x,y$.
Write the base cycle as $v_0v_1\cdots v_{s-1}v_0$, with branch vertices
$u=v_0$ and $v=v_\alpha$.
Delete the base edges $v_0v_1$ and $v_\alpha v_{\alpha+1}$.
Retain the full added path between $u$ and $v$.
The result is a spanning path with endpoints $v_1$ and $v_{\alpha+1}$.
This also applies when $\alpha=1$.
\end{proof}

\begin{lemma}[Packing the short length interval]
\label{lem:rhost-theta-interval}
Let $d,b$ be positive integers with $b=o(d)$ along the sequence under consideration.
Put
\[
 T=2(d+b+1),\qquad H=d+5,\qquad
 J=\left\lfloor\frac{T-H+1}{8}\right\rfloor.
\]
For all sufficiently large $d$, one can cover every length $3,\ldots,T$
with disjoint abstract blocks of total rank
\[
 T-2-4J\le\frac32d+b+5.
\]
Each block is either a cycle or a theta graph.
Every theta block has a base cycle of length at most $4J+2$,
and its additional path has length at least $d+4$ and at most $O(d)$.
\end{lemma}

\begin{proof}
For $0\le j<J$, write $h=H+8j$ and use the following four triples $(s,x,y)$:
\[
 \begin{array}{c|c|c|c}
 s&x&y&z=(x+y-s)/2\\ \hline
 4j+3&h+4&h+5&H+6j+3\\
 4j+4&h&h+2&H+6j-1\\
 4j+5&h+6&h+7&H+6j+4\\
 4j+6&h+1&h+3&H+6j-1
 \end{array}
\]
Each row satisfies Lemma~\ref{lem:rhost-theta-block}.
These blocks cover exactly the two disjoint sets
\[
 \{3,\ldots,4J+2\}\quad\text{and}\quad\{H,\ldots,H+8J-1\}.
\]
Their disjointness follows from $b=o(d)$.
Use one cycle for every remaining length in $\{3,\ldots,T\}$.
There are $4J$ theta blocks and $T-2-12J$ single-cycle blocks.
Their total rank is $8J+(T-2-12J)=T-2-4J$.
The claimed estimate follows from the definition of $J$.
\end{proof}

\begin{theorem}[An upper coefficient of three halves]
\label{thm:rhost-three-halves}
For the same $p_*$ as in Theorem~\ref{thm:rhost-scale}, and every $p\ge p_*$,
with high probability
\[
 \operatorname{Pex}(G(n,p))
 \le\log_2n+\left(\frac32+o(1)\right)\frac{\ln n}{\ln\ln n}.
\]
\end{theorem}

\begin{proof}
Keep $p_1,\ldots,p_5,\beta,a,d,K_0,K,L_0$ from the audited construction above.
Keep the two-digit gadget with $b=\lceil\sqrt{2\ln n}\rceil$.
Its shortest represented length is $T+1$, where $T=2(d+b+1)$.
We replace the individual short chords by the blocks in Lemma~\ref{lem:rhost-theta-interval}.
Write $A=4J$ for the number of theta blocks and $B=T-2-8J$ for the number of all short blocks.
Their total rank is $A+B=T-2-4J$.

\emph{Initial components.}
Split the $G_1$ exposure into four independent subexposures with equal edge probability
$q_1=1-(1-p_1)^{1/4}$.
Then $nq_1\to\infty$.
Use the first subexposure for the original binary base cycles $C_0,\ldots,C_{K_0}$.
Use the second for one base cycle of each short block.
Use the third for the two-digit base cycle $C_{\rm short}$.
The cycle-length input used in the original construction applies to each subexposure.

All requested lengths are at most the largest original binary base-cycle length, for sufficiently large $n$.
The original short-cycle separation estimate therefore makes all distinct selected cycles vertex-disjoint.
Cycles selected from different subexposures are also distinct with high probability.
Indeed, each pairwise intersection has edge probability $q_1^2$ and $nq_1^2=o(1)$.
Its expected total number of cycles is $o(1)$.
Finally, choose an edge $f$ in the fourth subexposure outside all selected vertices.
This is possible with high probability, conditionally on the earlier selections.

For each theta base cycle, fix its branch vertices and its spanning-path endpoints from
Lemma~\ref{lem:rhost-theta-block}.
For each single-cycle block, delete a designated edge to specify its spanning-path endpoints.
All these endpoints are now fixed, before $G_2$ is exposed.

\emph{The main cycle and the chain of short blocks.}
Build the binary base cycle $C_*$ through the designated edges of
$C_0,\ldots,C_{K_0}$ and through the extra edge $f$.
Use $K_0+2$ internally disjoint $G_2$ paths, each of length $d+2$.
Thus
\[
 |C_*|=(K_0+2)(d+3)=(1+o(1))\log_2n.
\]
Reserve $f$ as the connector edge.
The exact parameter estimates still give $|C_*|-K\to\infty$.
There are therefore enough other edges for the remaining binary detours.

Join one endpoint of $f$ to one endpoint of the designated edge $e_{\rm short}$
through all $B$ short blocks, in a fixed order.
Use $B+1$ additional $G_2$ paths of length $d+2$ between the fixed spanning-path endpoints.
Their interiors avoid every initial component and every earlier connecting path.
The total number of $G_2$ requests is $O(d)$.
Their total vertex cost is $O(d^2)$.

Here the corrected exposure argument remains valid for this larger number of requests.
For a first endpoint query, inspect $h_0=\lceil\beta n/100\rceil$ unused candidate edges.
Its failure probability is at most $(\ln n)^{-1}$.
The leaf query has at least $\beta n/80$ unused candidates and failure probability at most $(\ln n)^{-5/4}$.
Only the selected endpoint of $f$ appears in two requests.
In its first request, use it for the bounded first-endpoint query.
In its second request, use it for the leaf query.
All other endpoints appear once.
The total failure probability is
\[
 O\!\left(d\bigl((\ln n)^{-1}+(\ln n)^{-5/4}\bigr)\right)=o(1).
\]
The expansion and rooted-tree inputs are those proved above.

\emph{The additional paths in $G_3$.}
Embed the original binary detours and the two-digit detours with the audited procedure.
For every theta block, embed its additional path of length $z$ from the table.
This requires a rooted depth-$d$ tree with a stem of length $z-d-2\ge2$.
The new path uses $O(d)$ vertices, and all such paths use $O(d^2)$ vertices together.
Their endpoints lie in distinct base cycles and were fixed before $G_3$ was exposed.
The first-endpoint and leaf failure bounds above apply.
There are $O(d)$ such requests, so their total failure probability also tends to zero.
Expose only reservoir edges before querying these endpoint edges.

The total number of used vertices remains at most
\[
 L_0+O(d^2+b^3)=L_0+o(L_0)<2L_0.
\]
Thus the corrected robust expansion statement applies at every stage.

\emph{Hamilton completion and the edge count.}
Replace each short block by its fixed spanning path.
Their connecting chain joins one endpoint of $f$ to one endpoint of $e_{\rm short}$.
In the main binary unit, delete $f$ and replace every binary shortcut edge by its detour.
In the two-digit unit, delete $e_{\rm short}$ and replace every designated edge by its detour.
These replacements and the short-block chain form one spanning path of the exposed construction.
Its two endpoints are the unused endpoints of $f$ and $e_{\rm short}$.
Apply the unchanged $G_4$ Hamilton-path completion to these endpoints and the remaining vertices.

The resulting graph contains a Hamilton cycle.
Its edges outside that Hamilton cycle number
\[
 (K+1)+2+2b+(A+B)
 =K+2b+(T-2-4J)+3.
\]
Here a single-cycle block contributes one edge, and a theta block contributes two.
All short lengths $3,\ldots,T$ are present by Lemma~\ref{lem:rhost-theta-interval}.
The two-digit unit supplies $[T+1,T+b^2]$, which reaches $|C_*|$.
The binary unit supplies $[|C_*|,|C_*|+L_0]$ and $[n-L_0,n]$.
Its vertices remain on one contiguous interval of the Hamilton cycle.
The original long interval-covering argument in $G_5$ therefore applies without change.
It uses at most $m=\lfloor n2^{-K}\rfloor=O(\sqrt{\ln n})$ further chords.

The final excess is at most
\[
 K+2b+(T-2-4J)+m+3
 \le K+\frac32d+3b+m+8
 \le\log_2n+\frac32d+O(\sqrt{\ln n}).
\]
Since $d=(1+o(1))\ln n/\ln\ln n$, this proves the claim at $p_*$.
Monotonicity gives the result for $p\ge p_*$.
\end{proof}

Together, Theorems~\ref{thm:rhost-coefficient-one} and~\ref{thm:rhost-three-halves}
give, with high probability at $p=p_*$,
\[
 \boxed{\quad
 \log_2n+(1-o(1))\frac{\ln n}{\ln\ln n}
 \le\operatorname{Pex}(G(n,p_*))
 \le\log_2n+\left(\frac32+o(1)\right)\frac{\ln n}{\ln\ln n}.
 \quad}
\]

\subsection{Embedding whole theta blocks}
\label{subsec:rhost-whole-theta}

\begin{lemma}[Overlap estimate for a subdivided theta graph]
\label{lem:rhost-theta-moments}
Let $Q$ be a simple theta graph with $v$ vertices and $e=v+1$ edges.
Write $\ell_1,\ell_2,\ell_3$ for its three cycle lengths.
Let $J\sim G(n,D/n)$, with $D\ge54$ and $v^2/n=o(1)$.
Let $X$ count the copies of $Q$ in $J$.
Let $Y$ count ordered pairs of copies with a common vertex, including identical pairs.
Put
\[
 \mu=\mathbb E X=\frac{(n)_v}{|\operatorname{Aut}(Q)|}(D/n)^e,
 \qquad
 \varepsilon_Q=\frac{v^2}{n}
 +\sum_{i=1}^{3}\ell_iD^{-\ell_i}+nD^{-e}.
\]
There is an absolute constant $C$ such that
\[
 \operatorname{Var}(X)\le C\mu^2\varepsilon_Q,
 \qquad \mathbb E Y\le C\mu^2\varepsilon_Q.
\]
\end{lemma}

\begin{proof}
Use labelled embeddings first.
For two embeddings, their common vertices and common edges give a partially identified subgraph of $Q$.
An overlap with $u$ vertices and $f$ edges has normalized weight
\[
 (1+O(v^2/n))n^{-u}(n/D)^f.
\]
This follows by choosing the unshared vertices of the second embedding outside the first.
We sum over all possible overlap subgraphs and their injective maps into $Q$.
Allowing extra maps only increases the bound.

A connected tree overlap with $k$ edges has normalized weight $n^{-1}D^{-k}$.
There are at most $v9^k$ rooted tree subgraphs in the first copy.
A canonical depth-first traversal gives this bound, since the maximum degree is three.
There are at most $v3^k$ injective maps of each tree into the second copy.
The total contribution of one tree component is therefore at most
\[
 \frac{v^2}{n}\sum_{k\ge0}(27/D)^k\le\frac{2v^2}{n}.
\]
Allowing an unordered collection of tree components gives an exponential factor
$\exp(2v^2/n)$.
The nonempty forest overlaps contribute $O(v^2/n)$.

A proper connected overlap containing a cycle is unicyclic.
Its cycle is one of the three cycles of $Q$.
Outside that cycle, it contains only initial segments of the remaining theta path, from either endpoint.
For a cycle of length $\ell$, its image has at most three choices and at most $2\ell$ orientations and shifts.
Once this image is fixed, each additional initial segment has at most one possible image.
Summing their lengths gives a factor at most $(1-D^{-1})^{-2}$.
Thus these connected overlaps contribute at most
\[
 \frac{6}{(1-D^{-1})^2}\sum_{i=1}^{3}\ell_iD^{-\ell_i}.
\]
At most one unicyclic component can occur, since every two cycles of $Q$ meet.
Other components are trees and supply the same exponential factor as before.

An overlap with two independent cycles contains every edge of $Q$.
It contributes at most $12nD^{-e}$, because $|\operatorname{Aut}(Q)|\le12$.
These cases prove the bound for intersecting pairs of labelled embeddings.
Passing to unlabelled copies divides both pair counts and squared means by the same factor.

The contribution of vertex-disjoint pairs to $\mathbb E X^2$ is at most $\mu^2$.
Indeed, their vertex-placement factor is $(n-v)_v/(n)_v\le1$.
The remaining contribution is $\mathbb E Y$.
This also proves the variance bound.
\end{proof}

\begin{lemma}[Simultaneous disjoint theta embeddings]
\label{lem:rhost-theta-list}
Let $d\to\infty$, with $d=O(\ln n)$, and let $D=\Theta(d)$.
Assume $nD^{-d}$ tends to zero faster than every fixed negative power of $d$.
Consider $A=\Theta(d)$ simple theta graphs $Q_1,\ldots,Q_A$.
Suppose their numbers of edges lie between $d$ and $Cd$, for a fixed constant $C$.
Suppose their shortest cycle lengths are exactly $3,\ldots,A+2$,
and every other cycle has length at least $A+3$.
Then, with high probability, $G(n,D/n)$ contains vertex-disjoint copies of all $Q_i$.
They can also avoid any prescribed set of $O(d^2)$ vertices, even if that set is chosen after the graph is observed.
\end{lemma}

\begin{proof}
Choose $M=C_0d^2$, with $C_0$ larger than the total required vertex budget and the prescribed forbidden-set bound.
We show that every $Q_i$ has more than $M$ pairwise vertex-disjoint copies.
This permits a greedy choice of one copy of each type.

For one type, use $X,Y,\mu,\varepsilon_Q$ from Lemma~\ref{lem:rhost-theta-moments}.
Chebyshev's inequality bounds $\mathbb P(X<\mu/2)$ by $4C\varepsilon_Q$.
Markov's inequality bounds $\mathbb P(Y>\mu^2/(8M))$ by $8CM\varepsilon_Q$.
On the complementary event, form the graph whose vertices are the copies of $Q$.
Join two copies when they share a vertex.
It has an independent set of size at least $X^2/Y\ge2M$.
This bound follows by taking a random vertex order and retaining vertices earlier than all their neighbours.
Cauchy--Schwarz gives an expected retained size at least $X^2/Y$.

The sum of the error parameters over all types is
\[
 O\!\left(\frac{d^3}{n}+D^{-3}+d\,nD^{-d}\right).
\]
Here the distinct short cycle lengths give
$\sum_{s\ge3}sD^{-s}=O(D^{-3})$.
The longer cycle lengths give a smaller contribution.
After multiplication by $M=O(d^2)$, this sum tends to zero.
Indeed, $d^5/n=o(1)$ and $d^2D^{-3}=O(1/d)$.
The hypothesis handles the final term.

Each forbidden vertex can meet at most one member of a vertex-disjoint family of copies.
At every greedy step, fewer than $M$ vertices have been excluded.
An available copy therefore remains.
\end{proof}

\begin{lemma}[A denser packing of the required lengths]
\label{lem:rhost-whole-theta-interval}
Put $T=2(d+b+1)$, where $b\to\infty$ and $b=o(d)$.
Set
\[
 J=\left\lfloor\frac{T-2}{12}\right\rfloor,\qquad A=4J.
\]
There are $A$ theta blocks covering exactly the lengths $3,\ldots,3A+2$.
Their shortest cycle lengths are $3,\ldots,A+2$.
Their other cycle lengths are at least $A+3$.
Each block has between $d$ and $O(d)$ edges.
With one single-cycle block for each remaining length through $T$,
the total rank is at most
\[
 \frac43d+\frac43b+4.
\]
\end{lemma}

\begin{proof}
For $0\le j<J$, put $q=J-j-1$ and $h=A+3+8j$.
Use the four triples
\[
 (4q+3,h+4,h+5),\quad (4q+4,h,h+2),\quad
 (4q+5,h+6,h+7),\quad (4q+6,h+1,h+3).
\]
Lemma~\ref{lem:rhost-theta-block} gives the required theta graphs.
The low triples run through $3,\ldots,A+2$ in reverse blocks.
The high triples run through every length $A+3,\ldots,3A+2$ once.
The four edge counts in block $j$ are
\[
 6J+6j+7,\quad6J+6j+4,\quad6J+6j+10,\quad6J+6j+6.
\]
Their minimum is at least $6J+4\ge d+b-1\ge d$.
Their maximum is $O(d)$.
There are at most eleven remaining lengths through $T$.
The total rank is
\[
 2A+(T-2-3A)=T-2-A
 \le\frac23(T-2)+4=\frac43d+\frac43b+4.
\]
\end{proof}

\begin{theorem}[An upper coefficient of four thirds]
\label{thm:rhost-four-thirds}
For the same $p_*$ as in Theorem~\ref{thm:rhost-scale}, and every $p\ge p_*$,
with high probability
\[
 \operatorname{Pex}(G(n,p))
 \le\log_2n+\left(\frac43+o(1)\right)\frac{\ln n}{\ln\ln n}.
\]
\end{theorem}

\begin{proof}
Split $G_2$ into two independent subexposures of probability
$q_2=1-\sqrt{1-p_2}$.
Their union has the original probability $p_2$.
Replace the expansion parameter by
\[
 \beta=\frac{4(\ln\ln n)^2}{\ln n},\qquad
 a=\left\lfloor(5\beta)^{-1}\right\rfloor,
 \qquad d=\left\lfloor\log_a(n/200)\right\rfloor.
\]
The remaining parameters retain their previous values.
The depth satisfies $d=(1+o(1))\ln n/\ln\ln n$, and the second $G_2$ subexposure obeys
\[
 q_2\beta n\ge100\ln\ln n.
\]
Thus the corrected expansion and endpoint-query bounds retain their stated constants.
They also hold in $G_3$, whose edge probability is larger.

Prepare the binary base cycles, the two-digit base cycle, the reserve edge $f$,
and the at most eleven single-cycle blocks in $G_1$, as above.
Use the first $G_2$ subexposure to embed the complete theta list from
Lemma~\ref{lem:rhost-whole-theta-interval}, avoiding these initial vertices.
This is justified by Lemma~\ref{lem:rhost-theta-list}.
Indeed, $D=nq_2=\Theta(d)$, and
\[
 \ln D-\ln a\ge\ln(500\ln\ln n).
\]
Consequently,
\[
 nD^{-d}
 \le\exp\!\left(-\Omega\!\left(
 \frac{\ln n\,\ln\ln\ln n}{\ln\ln n}\right)\right).
\]
This is smaller than every fixed negative power of $d$.

All short blocks now exist before the second $G_2$ subexposure.
Choose their spanning-path endpoints and join them into the same connector chain as before.
Construct the binary base cycle through $f$ in this second subexposure as well.
There are $O(d)$ endpoint requests, with total failure probability $o(1)$.
Embed the binary and two-digit detours in $G_3$.
No additional theta path is needed there.
The vertex budget is still $L_0+O(d^2+b^3)<2L_0$.

Use the same $G_4$ Hamilton completion and $G_5$ long interval-covering chords.
The short blocks cover $3,\ldots,T$.
Their total rank is at most $4d/3+4b/3+4$.
Thus the excess is at most
\[
 K+2b+\left(\frac43d+\frac43b+4\right)+m+3
 \le\log_2n+\frac43d+O(\sqrt{\ln n}).
\]
The five original exposure probabilities have not increased.
This proves the claim at $p_*$, and monotonicity gives the remaining range.
\end{proof}

Combining the coefficient-one lower bound with the whole-theta construction gives
\[
 \boxed{\quad
 \log_2n+(1-o(1))\frac{\ln n}{\ln\ln n}
 \le\operatorname{Pex}(G(n,p_*))
 \le\log_2n+\left(\frac43+o(1)\right)\frac{\ln n}{\ln\ln n}
 \quad}
\]
with high probability.
This construction leaves the sharp coefficient and the full-threshold question unresolved.
Both are treated further in the later sections.

\section{Integer bands from traceable subdivisions of $K_4$}
\label{sec:rhost-k4-bands}

A graph is traceable if it has a Hamilton path.
Every subdivision below retains one edge of $K_4$ without subdivision.
This is necessary for the stated spanning-path construction.
A subdivision with an internal vertex on all six arms has no Hamilton path.
Indeed, such a path would delete three edges while reducing all four branch degrees.
Each deleted edge would then meet at most one branch vertex.

\begin{lemma}\label{lem:rhost-k4-bands}
For each integer $N\ge1$, there are $2(N+1)$ traceable subdivisions of $K_4$ with these properties:
\begin{enumerate}
\item Every cycle has length at least $2N+6$.
\item Every bicyclic core has at least $7N+13$ edges.
\item Each whole subdivision has at least $14N+15$ edges and at most $15N+14$ vertices.
\item Their cycle lengths include every integer in $[2N+6,14N+10]$ except possibly $4N+8$.
\end{enumerate}
Consequently, adding $2N+4$ individual cycles covers all lengths $3,\ldots,14N+10$.
The total cycle rank of these blocks is $8N+10$.
Their total number of vertices is at most $(61N^2+141N+94)/2$.
\end{lemma}

\begin{proof}
Label the vertices of $K_4$ by $a,b,c,d$.
List its six arms in the order
\[
 (ab,ac,ad,bc,bd,cd).
\]
For each integer $t\in\{0,1,\ldots,N\}$, define two subdivisions by the following arm lengths:
\begin{align*}
 A_t&=(8N-2t+3,\ 1,\ N-t+3,\ 2N+2t+3,\ N+t+3,\ 3N-t+3),\\
 B_t&=(1,\ N-t+2,\ N-t+3,\ 3N-t+3,\ 3N+3t+3,\ 6N+3).
\end{align*}
All arm lengths are positive integers.
The arm $ac$ of $A_t$ and the arm $ab$ of $B_t$ each have length one.
Every other arm has length at least two.

The six selected cycle lengths in the following table each run through an arithmetic progression with step two.
The interval gives the two endpoint values, irrespective of direction.
\[
\begin{array}{c|cc|cc}
 \text{band}&\text{cycle in }A_t&\text{length interval}&
              \text{cycle in }B_t&\text{length interval}\\\hline
 0&acd&[2N+7,4N+7]&abc&[2N+6,4N+6]\\
 1&acbd&[4N+10,6N+10]&abd&[4N+7,6N+7]\\
 2&bcd&[6N+9,8N+9]&acd&[6N+8,8N+8]\\
 3&abd&[8N+9,10N+9]&abcd&[8N+10,10N+10]\\
 4&abdc&[10N+10,12N+10]&abdc&[10N+9,12N+9]\\
 5&abcd&[12N+12,14N+12]&bcd&[12N+9,14N+9]
\end{array}
\]
The unlisted cycles have lengths $10N+7$ in $A_t$ and $8N+11$ in $B_t$.
These formulas show that every cycle has length at least $2N+6$.
Each band supplies the two opposite parities.
The endpoint intervals overlap across all interfaces except the single possible gap $4N+8$.
Their union therefore contains $[2N+6,14N+10]\setminus\{4N+8\}$.

The total arm lengths are
\[
 E(A_t)=15N-t+16,\qquad E(B_t)=14N+15.
\]
A bicyclic core in a subdivision of $K_4$ consists of five complete arms.
For $A_t$, the longest arm is $ab$, so every such core has at least $7N+t+13$ edges.
For $B_t$, the longest arm is $cd$, tied by $bd$ when $t=N$.
Every such core therefore has at least $8N+12\ge7N+13$ edges.
The whole subdivisions have cycle rank three and $E-2$ vertices.
This proves the edge and vertex bounds.

It remains to verify traceability.
In $A_t$, omit the edge $ac$.
Delete the edge at $b$ on the arm $ab$, and the edge at $d$ on the arm $cd$.
The remaining graph is the spanning path whose branch order is
\[
 a,d,b,c.
\]
Its endpoints are the internal vertices adjacent to $b$ on $ab$ and to $d$ on $cd$.
In $B_t$, omit $ab$.
Delete the edge at $c$ on $ac$, and the edge at $d$ on $bd$.
The remaining graph is a spanning path with branch order
\[
 a,d,c,b.
\]
Its endpoints are the internal vertices adjacent to $c$ on $ac$ and to $d$ on $bd$.
Every arm used to specify these endpoints has an internal vertex.

Add one individual cycle of each length $3,\ldots,2N+5$, and one of length $4N+8$.
This uses $2N+4$ cycles and fills the prefix and the possible gap.
The cycle-rank sum is
\[
 3\cdot2(N+1)+(2N+4)=8N+10.
\]
The $K_4$ subdivisions together have
\[
 \sum_{t=0}^{N}\bigl(15N-t+14+14N+13\bigr)
 =\frac{57N^2+111N+54}{2}
\]
vertices.
The individual cycles add $2N^2+15N+20$ vertices.
Adding these expressions gives the claimed total.
\end{proof}

For a target interval $3,\ldots,T$, set $N=\lceil T/14\rceil$.
The lemma gives total block rank at most $4T/7+18$.
All blocks together use $O(T^2)$ vertices.
If $d=7N$, the whole-edge threshold is at least $2d$, every bicyclic core has at least $d$ edges,
and every cycle has at least $2d/7$ edges.
These inequalities are the local density inputs for packing the blocks into a random host.
The lemma itself does not assert a random-host embedding theorem.

The exact formulas and the Hamilton paths are checked by
\nolinkurl{check_k4_random_host_blocks.py}.

% Further random-host refinement. Include after random-hosts-phase2.tex
% and k4-random-host-blocks.tex.
\section{A random-host upper coefficient of eight sevenths}
\label{sec:rhost-eight-sevenths}

\paragraph{Status.}
This section uses the explicit traceable $K_4$ blocks in Lemma~\ref{lem:rhost-k4-bands}.
It adds a uniform overlap estimate for their random embeddings.
This construction leaves the sharp coefficient and the full-threshold question unresolved.
Later sections improve the coefficient and address the threshold.

\begin{lemma}[Overlap estimate for a subdivided $K_4$]
\label{lem:rhost-k4-moments}
Let $Q$ be a subdivision of $K_4$, with $v$ vertices and $e=v+2$ edges.
Let $g$ be its shortest cycle length.
Let $b$ be the smallest number of edges in a connected subgraph of minimum degree at least two and rank two.
Let $J\sim G(n,D/n)$, with $D\ge54$ and $v^2/n=o(1)$.
Let $X$ count copies of $Q$, and let $Y$ count ordered intersecting pairs, including identical pairs.
Put $\mu=\mathbb E X$ and
\[
 \varepsilon_Q=\frac{v^2}{n}
 +(v+1)^{32}\left(D^{-g}+nD^{-b}+n^2D^{-e}\right).
\]
For an absolute constant $C$,
\[
 \operatorname{Var}(X)\le C\mu^2\varepsilon_Q,
 \qquad \mathbb E Y\le C\mu^2\varepsilon_Q.
\]
\end{lemma}

\begin{proof}
The normalized weight of an overlap with $u$ vertices and $f$ edges is
$(1+O(v^2/n))n^{-u}(n/D)^f$.
As in Lemma~\ref{lem:rhost-theta-moments}, all forest overlaps contribute $O(v^2/n)$.
Their disjoint components also supply at most the factor $\exp(2v^2/n)$ when a cyclic component is present.

A subdivision of $K_4$ has no two vertex-disjoint cycles.
Thus at most one overlap component contains a cycle.
Strip its degree-one vertices to obtain its nonempty core.
Each original subdivided edge is either wholly present in this core or absent from it.
The core is therefore one of the following:
a subdivided cycle, a subdivision of $K_4$ with one edge deleted, or all of $Q$.
Its excess $f-u$ is respectively zero, one, or two.
Its edge count is at least $g$, $b$, or $e$.

We give a polynomial bound for the number of connected cyclic overlap descriptions.
On each of the six original edge paths, such a component uses at most an initial segment from each endpoint.
Thus at most twelve cut positions specify its edge set, after the core type is chosen.
There are $O((v+1)^{12})$ possible source subgraphs.
Each has at most twelve vertices of degree one and four vertices of degree three.
Choose the images of these at most sixteen vertices in the second copy.
After suppressing degree-two vertices, at most twenty-four path segments remain.
Between two specified vertices in $Q$, there are at most $2^{\beta(Q)}=8$ simple paths.
A loop segment has at most fourteen oriented cycle images, since $Q$ has seven cycles.
Hence the segment images have only an absolute constant number of choices.
A component that is itself a cycle has at most $14v$ possible images and obeys the same bound.
There are therefore $O((v+1)^{28})$ connected cyclic overlap descriptions.

Their total normalized contribution is at most
\[
 O((v+1)^{28})\left(D^{-g}+nD^{-b}+n^2D^{-e}\right).
\]
The forest components and falling-factorial correction give the stated estimate with exponent $32$.
The argument applies first to labelled embeddings and then to unlabelled copies.
Vertex-disjoint pairs contribute at most $\mu^2$, which also proves the variance estimate.
\end{proof}

\begin{corollary}[A list of robust $K_4$ embeddings]
\label{cor:rhost-k4-list}
Let $d=(1+o(1))\ln n/\ln\ln n$ and $D=\Theta(d)$.
Assume $nD^{-d}$ tends to zero faster than every fixed negative power of $d$.
Consider $O(d)$ subdivisions of $K_4$.
Suppose each has $O(d)$ edges, shortest cycle length at least $d/4$,
every rank-two core has at least $d$ edges, and its full edge count is at least $2d$.
With high probability, all these graphs embed vertex-disjointly in $G(n,D/n)$.
They can avoid any prescribed set of $O(d^2)$ vertices, even if that set is chosen after observing the graph.
\end{corollary}

\begin{proof}
Apply Lemma~\ref{lem:rhost-k4-moments} to each graph in the list.
The sum of its error parameters, multiplied by $O(d^2)$, tends to zero.
The forest term contributes $O(d^5/n)$.
The cycle term tends to zero because $D^{-d/4}=n^{-1/4+o(1)}$.
The other terms tend to zero by the hypothesis on $nD^{-d}$.
The Chebyshev, Markov, and independent-set argument in Lemma~\ref{lem:rhost-theta-list}
therefore supplies more than $M$ vertex-disjoint copies of each type, for any fixed $M=O(d^2)$.
Choose one copy of each type greedily, avoiding the prescribed vertices and earlier choices.
\end{proof}

\begin{theorem}[An upper coefficient of eight sevenths]
\label{thm:rhost-eight-sevenths}
For the same $p_*$ as in Theorem~\ref{thm:rhost-scale}, and every $p\ge p_*$,
with high probability
\[
 \operatorname{Pex}(G(n,p))
 \le\log_2n+\left(\frac87+o(1)\right)\frac{\ln n}{\ln\ln n}.
\]
\end{theorem}

\begin{proof}
Use the exposure scheme and parameters of Theorem~\ref{thm:rhost-four-thirds}.
In particular, split $G_2$ into two independent subexposures of probability
$q_2=1-\sqrt{1-p_2}$, and use
\[
 \beta=\frac{4(\ln\ln n)^2}{\ln n},\qquad
 a=\lfloor(5\beta)^{-1}\rfloor,
 \qquad d=\lfloor\log_a(n/200)\rfloor.
\]
Let $b=\lceil\sqrt{2\ln n}\rceil$, $T=2(d+b+1)$, and $N=\lceil T/14\rceil$.
Replace the theta list by the two families of $K_4$ blocks from Lemma~\ref{lem:rhost-k4-bands}.
They use $2(N+1)$ blocks of rank three.
Add one single-cycle block for each length $3,\ldots,2N+5$, and one for length $4N+8$.
The resulting blocks cover every length through $14N+10\ge T$.
Their total rank is exactly
\[
 6(N+1)+(2N+4)=8N+10
 \le\frac87(d+b+1)+18.
\]

Prepare the single-cycle blocks, the binary base cycles, the two-digit base cycle,
and the reserve edge in $G_1$, using the separated subexposures described above.
Their total vertex count is $O(d^2)$.
Use the first $G_2$ subexposure for the complete $K_4$ blocks.
Corollary~\ref{cor:rhost-k4-list} applies.
Indeed, the metric lemma gives shortest cycles $2N+6>d/4$,
rank-two cores with at least $7N+13\ge d$ edges, and whole blocks with at least $14N+15\ge2d$ edges.
Also, $D=nq_2=\Theta(d)$ and $nD^{-d}$ has the required superpolynomial decay.

Each block has a specified unsubdivided arm and therefore has a spanning path.
Choose those paths after the embeddings are found.
Use the second $G_2$ subexposure to join the short blocks into the connector chain.
The previously audited endpoint-query and vertex-budget estimates remain unchanged.
Complete the binary and two-digit gadgets in $G_3$ and the Hamilton cycle in $G_4$.
Use the original $G_5$ long interval-covering chords.

The blocks cover every length through $T$.
The two-digit and binary gadgets cover the remaining intervals as before.
The final excess is at most
\[
 K+2b+8N+m+O(1)
 \le\log_2n+\frac87d+O(\sqrt{\ln n}).
\]
This gives the asserted asymptotic bound at $p_*$.
Monotonicity handles $p\ge p_*$.
\end{proof}

Together with Theorem~\ref{thm:rhost-coefficient-one}, this gives
\[
 \boxed{\quad
 \log_2n+(1-o(1))\frac{\ln n}{\ln\ln n}
 \le\operatorname{Pex}(G(n,p_*))
 \le\log_2n+\left(\frac87+o(1)\right)\frac{\ln n}{\ln\ln n}
 \quad}
\]
with high probability.

\section{An exact theta partition of the middle interval}
\label{sec:rhost-theta-partition}

\begin{lemma}\label{lem:rhost-theta-middle}
Let $d=9q$, where $q\ge1$ is an integer.
There are exactly $2d/9-1$ traceable theta graphs whose cycle lengths partition
\[
 \{d/3+2,d/3+3,\ldots,d-2\}.
\]
Each theta graph has exactly $d$ edges.
Every branch has at least two edges, and every cycle has length at least $d/3+2$.
The total cycle rank is $4d/9-2$, and the total vertex count is $(2d/9-1)(d-1)$.
\end{lemma}

\begin{proof}
Put
\[
 m=2q-1,\qquad h=q-1,\qquad A=3q+2.
\]
Thus $m=2h+1$ is odd.
For $0\le i<m$, define
\[
 \pi(i)=(i+h)\bmod m,\qquad \sigma(i)=3h-i-\pi(i).
\]
The map $\pi$ is a permutation of $\{0,\ldots,m-1\}$.
For $0\le i\le h$, we have $\sigma(i)=2h-2i$, which runs through all even values.
For $h+1\le i\le2h$, we have $\sigma(i)=4h+1-2i$, which runs through all odd values.
Hence $\sigma$ is also a permutation of the same set.

For each $i$, form the triple
\[
 L_{i,1}=A+i,\qquad
 L_{i,2}=A+m+\pi(i),\qquad
 L_{i,3}=A+2m+\sigma(i).
\]
The triples partition the consecutive interval from $A=d/3+2$ to $A+3m-1=d-2$.
Each triple sums to
\[
 3A+3m+i+\pi(i)+\sigma(i)
 =3A+3m+3h=18q=2d.
\]
Take a theta graph with branch lengths $d-L_{i,1}$, $d-L_{i,2}$, and $d-L_{i,3}$.
These are integers at least two and sum to $d$.
Its three cycle lengths are exactly the entries of the triple.
This also verifies every triangle inequality and the required even total without a separate parity adjustment.

A theta graph with all branches of length at least two has a Hamilton path.
Delete one edge at the first branch vertex on one branch, and one edge at the second branch vertex on another branch.
The remaining graph is a spanning path.
Each block has cycle rank two and $d-1$ vertices.
There are $m=2d/9-1$ blocks, giving the stated totals.
\end{proof}

For an arbitrary positive threshold $d_0$, take $d=9\lceil d_0/9\rceil$.
Then $d_0\le d<d_0+9$.
The covered lengths include every integer in $[\lceil d_0/3\rceil+5,\lfloor d_0\rfloor-2]$.
The number of theta blocks is $2d_0/9+O(1)$.
Their edge counts are at least $d_0$, so the rounding preserves the embedding density condition.

The permutation and length identities are checked by
\nolinkurl{check_theta_random_host_partition.py}.

% Independently audited by root and four_regular; finite checks are recorded below.
% Include after random-hosts-phase3.tex and theta-random-host-partition.tex.
\section{A wheel construction and an upper coefficient of ten ninths}
\label{sec:rhost-ten-ninths}

This section adds one external matching theorem to the previous construction.
All metric distributions and random-host embedding estimates are proved below.
The matching input is Pippenger's near-perfect matching theorem.
We use its version allowing $o(v)$ exceptional vertices, as stated after Theorem~1.1 in
Alon and Yuster, \emph{On a hypergraph matching problem},
\url{https://web.math.princeton.edu/~nalon/PDFS/ayhyper2.pdf}.

\begin{lemma}[Wheel length packing]\label{lem:rhost-wheel-lengths}
For integers $h\to\infty$, there are $(1+o(1))h/12$ traceable subdivisions of the four-rim wheel.
Their cycle lengths cover all but $o(h)$ integers in $[h+1,2h]$.
Each subdivision has $O(h)$ edges, shortest cycle length at least $2\lceil\sqrt h\rceil$,
and at least $(r-1)h$ edges in every connected core of cycle rank $r\in\{2,3,4\}$.
\end{lemma}

\begin{proof}
Put $s=\lceil\sqrt h\rceil$.
Choose $U,L,M$ independently and uniformly from the integers $s,\ldots,2s$.
Put $t=h+U$ and
\[
 c_j=4^{-j}\binom{2j}{j}\qquad(j\ge0).
\]
Choose four nonnegative integers $X_0,\ldots,X_3$ of sum $t$, with probability
\begin{equation}
 \Pr(X_0=x_0,\ldots,X_3=x_3)
   =\frac{c_{x_0}c_{x_1}c_{x_2}c_{x_3}}{t+1}.
 \label{eq:rhost-wheel-distribution}
\end{equation}
This is a probability distribution because $\sum_{j\ge0}c_jz^j=(1-z)^{-1/2}$.
For every pair $i,j$, the sum $X_i+X_j$ is uniform on $\{0,\ldots,t\}$.
Indeed, each convolution of two sequences $(c_j)$ is identically one.

Give the four spokes lengths $w_i=h+s-X_i$.
Give the cyclic rim arms $01,12,23,30$ lengths $1,L,1,M$.
Discard the outcome if a spoke length is nonpositive.
The probability of this event is $o(1)$.
To see this, fix a spoke and put $r=t-X_i$.
The generating function for the other three coordinates gives
\[
 \Pr(t-X_i=r)=\frac{c_{t-r}(2r+1)c_r}{t+1}.
\]
Since $c_j\le C/\sqrt{j+1}$, summing over $r\le s$ gives $O((s/h)^{3/2})$.
A union bound handles all four spokes.

There are twelve non-rim cycles, two for each pair of spokes.
Their lengths have the form
\[
 Z_{ij,a}=2h+2s-X_i-X_j+a.
\]
The two possible offsets for each pair are
\[
\begin{array}{c|c}
 ij&\text{offsets}\\\hline
 01,23&1,\ L+M+1\\
 12&L,\ M+2\\
 30&M,\ L+2\\
 02,13&L+1,\ M+1.
\end{array}
\]
All these lengths lie in
$J=[h+1,2h+6s+1]\cap\mathbb Z$.
For every $z\in I=[h+5s+1,2h]\cap\mathbb Z$ and each specified cycle label,
\begin{equation}
 \Pr(Z_{ij,a}=z)=\mathbb E\frac1{h+U+1}
                  =\frac{1+O(s/h)}h.
 \label{eq:rhost-wheel-marginal}
\end{equation}
For every integer $z$, the same probability is at most $1/h$.

We next bound the joint probabilities of two distinct cycle labels.
For the same spoke pair, the difference of the offsets is one of
$\pm(L+M)$, $\pm(L-M)$, or $\pm(L-M+2)$.
For fixed target lengths, this relation has probability $O(1/s)$.
Conditioning on $U,L,M$ then costs at most $1/h$ for the required pair sum.
For complementary spoke pairs, the sum of the two lengths is
\[
 3h+4s-U+a+a'.
\]
Conditioning on $L,M$ fixes at most one value of $U$.
The remaining pair-sum condition again costs at most $1/h$.
Both cases therefore give $O(1/(hs))$.

Suppose that the spoke pairs share one index.
Condition on $X_i+X_j=q$.
The split within this pair and the split within its complement are independent.
A split of an integer $q$ has masses $c_a c_{q-a}$, whose maximum is at most $C/\sqrt{q+1}$.
The shared-index pair sum is a convolution of splits of $q$ and $t-q$.
Its largest mass is at most
\[
 \min\left\{\frac C{\sqrt{q+1}},\frac C{\sqrt{t-q+1}}\right\}
 \le\frac C{\sqrt h}.
\]
The first pair sum has mass $1/(t+1)$.
Thus, uniformly in two target lengths and two distinct labels,
\begin{equation}
 \Pr(Z=z,Z'=z')=O(h^{-3/2}).
 \label{eq:rhost-wheel-codegree}
\end{equation}
Here we used $s=\lceil\sqrt h\rceil$.
Summing over $z=z'\in J$ shows that repeated cycle lengths occur with probability $O(h^{-1/2})$.
Discard these outcomes too.

For each twelve-element subset $A\subset J$, let $p_A$ be the probability that the retained outcome has length set $A$.
Then
\[
 \sum_Ap_A=1-o(1),\qquad
 \sum_{A\ni z}p_A\le\frac{12}h,\qquad
 \max_{z\ne z'}\sum_{A\supset\{z,z'\}}p_A=O(h^{-3/2}).
\]
Before the discarded outcomes, every vertex of $I$ has labelled incidence weight $(12+o(1))/h$.
Discarding outcomes removes total incidence weight $o(1)$.
Consequently, all but $o(h)$ vertices of $J$ have degree weight $(12+o(1))/h$.

This weighted system reduces to a simple hypergraph.
Set $\lambda=h(\ln h)^2$ and include each distinct set $A$ independently with probability $\lambda p_A$.
These probabilities are below one for large $h$, by the codegree bound.
Chernoff bounds show that the maximum degree is at most $(12+o(1))(\ln h)^2$.
All but $o(h)$ vertices have degree $(12+o(1))(\ln h)^2$.
The maximum codegree is $O(\ln h)$.
For completeness, each degree is a sum of independent indicators with mean at most $12(\ln h)^2$.
Use an error $(\ln h)^{-1/4}$ and a union bound over $O(h)$ vertices.
Each codegree has mean $O(h^{-1/2}(\ln h)^2)$.
The upper-tail estimate at $C\ln h$ permits a union bound over $O(h^2)$ pairs.
Pippenger's theorem, with its exceptional-vertex extension, gives a matching covering $|J|-o(h)$ vertices.
It has $h/12+o(h)$ edges.
Choose one retained wheel metric for each matching edge.
This proves the length coverage.

We verify the graph conditions.
The total spoke length is $3h+4s-U$, and the total edge count is at least $3h+4s+2$.
Every spoke has length at most $h+s$.
Every pair of spokes has total length at least $h$.
A non-rim cycle therefore has length at least $h+1$, while the rim has length $L+M+2\ge2s$.

A connected core of rank two contains either two spokes and the full rim, or at least three spokes.
Its edge count is at least $h$.
A core of rank three contains either three spokes and the full rim, or all four spokes.
Its edge count is at least $2h$.
A core of rank four is the whole wheel and has at least $3h$ edges.
These classifications follow by writing the rank as $k+r-v$,
where $k$ counts spokes and $r,v$ count rim edges and rim vertices.

Finally, delete the unit rim edges $01,23$.
Delete the hub-end edge of spokes $0,1$.
The remaining graph is a spanning path, with branch order $1,2,\text{hub},3,0$.
If a deleted spoke has length one, its rim vertex is an endpoint.
This proves traceability.
\end{proof}

\begin{lemma}[Wheel embeddings]\label{lem:rhost-wheel-embeddings}
Let $d\sim\ln n/\ln\ln n$, $h=d+O(1)$, and $D=\Theta(d)$.
Suppose $nD^{-d}$ tends to zero faster than every fixed negative power of $d$.
Any list of $O(d)$ wheel subdivisions supplied by Lemma~\ref{lem:rhost-wheel-lengths}
embeds vertex-disjointly in $G(n,D/n)$ with high probability.
The embeddings can avoid any set of $O(d^2)$ vertices chosen after observing the graph.
\end{lemma}

\begin{proof}
The overlap proof of Lemma~\ref{lem:rhost-k4-moments} extends to this fixed wheel kernel.
Here are the required details.
A wheel has no two vertex-disjoint cycles, so an overlap has at most one cyclic component.
After removing leaves, its core consists of complete kernel arms.
There are only constantly many core types.
For a core of rank $r$, its normalized overlap weight is at most $n^{r-1}D^{-e_r}$.
The metric lemma gives
\[
 e_1\ge2\lceil\sqrt h\rceil,\quad e_2\ge h,\quad e_3\ge2h,\quad e_4\ge3h.
\]
On eight arms, at most sixteen cut positions specify a connected overlap.
There are at most sixteen leaves and five other vertices of degree different from two.
Choose their images and then the images of the suppressed path segments.
A graph of cycle rank four has at most sixteen simple paths between prescribed endpoints.
It has at most fifteen simple cycles, so loop segments also have only constantly many choices.
Thus the number of cyclic overlap descriptions is at most $C(v+1)^{48}$.
Forest overlaps contribute $O(v^2/n)$, as before; maximum degree four only changes the geometric-series constant.
For large $n$, $D$ exceeds that constant.

Both the relative variance and the normalized intersecting-pair expectation are at most
\[
 C\left[\frac{v^2}n+(v+1)^{48}
 \left(D^{-2\lceil\sqrt h\rceil}+nD^{-h}+n^2D^{-2h}+n^3D^{-3h}\right)\right].
\]
This tends to zero faster than every fixed negative power of $d$.
The Chebyshev, Markov, and independent-set argument from Lemma~\ref{lem:rhost-theta-list}
therefore gives more than $C_0d^2$ disjoint copies of each metric, for every fixed $C_0$.
A greedy choice proves the assertion, including the adaptive forbidden set.
\end{proof}

\begin{theorem}[An upper coefficient of ten ninths]\label{thm:rhost-ten-ninths}
For the same $p_*$ as in Theorem~\ref{thm:rhost-scale} and every $p\ge p_*$,
with high probability
\[
 \operatorname{Pex}(G(n,p))
 \le\log_2n+\left(\frac{10}9+o(1)\right)\frac{\ln n}{\ln\ln n}.
\]
\end{theorem}

\begin{proof}
Use the exposure scheme of Theorem~\ref{thm:rhost-eight-sevenths}.
Put $h=9\lceil d/9\rceil$.
Use single cycles through length $h/3+1$.
Use the constant-sum theta partition in Lemma~\ref{lem:rhost-theta-middle} for $[h/3+2,h-2]$.
Use Lemma~\ref{lem:rhost-wheel-lengths} for $[h+1,2h]$.
Add one single-cycle block for every still-uncovered length through $T=2(d+b+1)$.
There are $o(d)$ such extra lengths, since $T-2h=O(b)=o(d)$.
The rank sum of the short blocks is at most
\[
 \frac h3+\frac{4h}9+4\frac h{12}+o(d)
 =\left(\frac{10}9+o(1)\right)d.
\]
All these choices depend only on $n$, before exposing the random host.

Prepare the single cycles and the previous base cycles in $G_1$.
Embed the theta and wheel blocks in the first half of $G_2$.
The theta moment estimate applies because every theta cycle has length at least $h/3+2$ and its edge count is $h\ge d$.
Lemma~\ref{lem:rhost-wheel-embeddings} supplies the wheel blocks.
Each type has more than $C_0d^2$ disjoint copies, so both lists can be selected greedily together.
Every block has a spanning path and the total number of vertices is $O(d^2)$.
Use the second half of $G_2$ for the connector chain.
The remaining exposures and interval coverage are unchanged.
The final excess is
\[
 K+2b+m+\left(\frac{10}9+o(1)\right)d
 =\log_2n+\left(\frac{10}9+o(1)\right)\frac{\ln n}{\ln\ln n}.
\]
\end{proof}

The finite checker \texttt{check\_random\_host\_wheels.py} verifies 3962 wheel metrics, all 42 kernel core types, and 120 explicit spanning paths.
It also verifies 96 pair-marginal identities with exact rational arithmetic.
These checks do not certify the asymptotic matching or embedding arguments.

\section{A weighted matching consequence of hypergraph edge colouring}
\label{sec:rhost-weighted-matching}

We use the Pippenger--Spencer edge-colouring theorem:
for fixed uniformity, an almost regular hypergraph with relative codegree $o(1)$ has chromatic index $(1+o(1))$ times its degree.
The primary paper is
N.~Pippenger and J.~Spencer, \emph{Asymptotic behavior of the chromatic index for hypergraphs},
Journal of Combinatorial Theory, Series~A \textbf{51} (1989), 24--42,
\url{https://doi.org/10.1016/0097-3165(89)90074-5}.
A precise statement also appears as Lemma~2.1 of Alon and Yuster,
\url{https://web.math.princeton.edu/~nalon/PDFS/ayhyper2.pdf}.

\begin{lemma}[Maximum degree and a weighted objective]\label{lem:rhost-max-degree-colouring}
Fix $r\ge3$.
Let $H$ be a simple $r$-uniform hypergraph with maximum degree at most $D\to\infty$ and maximum codegree $o(D)$.
Give each edge a nonnegative value $v_e$.
Then $H$ has a matching $M$ such that
\[
 \sum_{e\in M}v_e\ge\frac{1-o(1)}D\sum_{e\in E(H)}v_e.
\]
No lower bound on vertex degrees is required.
\end{lemma}

\begin{proof}
We reduce to the regular form of the edge-colouring theorem.
Take $L$ disjoint copies of $H$, where $L$ is a sufficiently large multiple of $r$.
For each original vertex $u$, group its $L$ copies together.
Independently choose $D-\deg_H(u)$ uniformly random partitions of that group into $r$-sets.
Add all sets in those partitions as edges.
Every vertex now has degree $D$, counted with multiplicity.

For a fixed pair within one group, its probability of sharing a part in one partition is $(r-1)/(L-1)$.
The probability that it shares a part in at least three partitions is $O(D^3/L^3)$.
A union bound over all groups and pairs gives $O(|V(H)|D^3/L)$.
The expected number of repeated added edges is $O(|V(H)|D^2L^{2-r})$.
Choose $L$ so large that both bounds tend to zero.
There is therefore a choice with no repeated added edges and with every new pair codegree at most two.

An original edge lies within one copy and uses distinct vertex positions.
An added edge lies within one vertex-position group and uses distinct copies.
No pair is contained in both an original and an added edge.
The resulting simple hypergraph is $D$-regular and has codegree at most $\max\{\Delta_2(H),2\}=o(D)$.
By Pippenger--Spencer, its edges split into $(1+o(1))D$ matchings.
Give added edges value zero and copied original edges their original values.
Some colour class has total value at least
\[
 \frac{L}{(1+o(1))D}\sum_{e\in E(H)}v_e.
\]
Within that colour class, some copy of $H$ carries at least the average of this value over the $L$ copies.
Its original edges form the required matching.
\end{proof}

\begin{lemma}[Weighted length packing]\label{lem:rhost-weighted-matching}
Fix integers $r\ge3$ and $C\ge1$, and constants $a,V,c>0$.
For $h\to\infty$, let $\mathcal A_h$ be a family of distinct $r$-element sets on at most $h^C$ vertices.
Give each set a nonnegative weight $p_A$ and a value $v_A\in[0,V]$.
Suppose
\[
 \sum_Ap_A=O(1),\qquad \sum_Ap_Av_A\ge c,
\]
\[
 \max_z\sum_{A\ni z}p_A\le\frac{a+o(1)}h,
 \qquad
 \max_{z\ne z'}\sum_{A\supset\{z,z'\}}p_A=O(h^{-3/2}).
\]
Then some pairwise disjoint subfamily $\mathcal M_h$ satisfies
\[
 \sum_{A\in\mathcal M_h}v_A
 \ge\left(\frac h a-o(h)\right)\sum_Ap_Av_A.
\]
\end{lemma}

\begin{proof}
Put $\lambda=h(\ln h)^2$.
Each $p_A$ is bounded by a pair weight, so $\lambda p_A<1$ for large $h$.
Include every distinct set independently with probability $\lambda p_A$.
This gives a simple $r$-uniform hypergraph $H$.

Every vertex degree has mean at most $(a+o(1))(\ln h)^2$.
Chernoff bounds, followed by a union bound over polynomially many vertices, give
\[
 \Delta(H)\le(a+o(1))(\ln h)^2
\]
with probability tending to one.
For example, an additional relative allowance $(\ln h)^{-1/4}$ suffices.
Every codegree has mean $O(h^{-1/2}(\ln h)^2)$.
A Chernoff upper-tail bound at $K\ln h$, followed by a union bound over polynomially many pairs, gives
$\Delta_2(H)=O(\ln h)$.

Let $W=\sum_{A\in E(H)}v_A$ and $S=\sum_Ap_Av_A$.
Then $\mathbb EW=\lambda S$ and
\[
 \operatorname{Var}(W)\le V^2\lambda\sum_Ap_A=O(\lambda).
\]
Since $S\ge c$, Chebyshev's inequality gives $W=(1-o(1))\lambda S$ with probability tending to one.
Fix a realization satisfying all these conclusions.
If needed, increase the degree parameter to an integer
$D=(a+o(1))(\ln h)^2\ge\Delta(H)$ in the regularisation proof of Lemma~\ref{lem:rhost-max-degree-colouring}.
That proof remains valid with this upper degree bound.
Its matching has value at least
\[
 \frac{(1-o(1))\lambda S}{(a+o(1))(\ln h)^2}
 =\left(\frac h a-o(h)\right)S.
\]
\end{proof}

\paragraph{Variable set sizes and private vertices.}
The lemma also applies when every real set has size between two and a fixed integer $r$.
Aggregate equal real sets before padding, retaining their weight and a fixed objective value for that set.
If a real set $A$ has fewer than $r$ vertices, add $r-|A|$ new private vertices used only by that set.
Every individual weight $p_A$ is at most a weighted codegree of two real vertices of $A$.
Thus each private vertex has weighted degree $O(h^{-3/2})$.
Every new pair codegree has the same bound, including pairs of private vertices from one padded set.
If the real ground set has $O(h)$ vertices, there are polynomially many possible real sets and private vertices.
The hypotheses therefore remain valid for the padded $r$-uniform family.
A matching of padded sets restricts to disjoint real sets and preserves its objective value.

In particular, wheel outcomes may retain $k\in\{4,\ldots,12\}$ required non-rim lengths and one marked rim length.
Their real sets have sizes $k+1\in\{5,\ldots,13\}$.
After padding to uniformity thirteen, their values $k-3\in\{1,\ldots,9\}$ count the saved block rank.

% Audited research proof; stronger bounds appear in the following refinement.
% Include after random-host-weighted-matching.tex and random-hosts-phase4.tex.
\section{Using the wheel rims as additional short cycles}
\label{sec:rhost-seventy-nine}

The previous wheel construction did not use its rim cycle to fill a distinct required length.
We now use a small proportion of longer rims for this purpose.
The resulting coefficient is $79/72<10/9$.

\begin{lemma}[A weighted wheel packing]\label{lem:rhost-wheel-profit}
Let $h\to\infty$ through multiples of nine.
There is a list of $O(h)$ traceable wheel subdivisions with the following properties.
Their non-rim cycle lengths form pairwise disjoint twelve-element subsets of $[h+1,2h]$.
Some wheels have a marked rim length in $[3,h/3+1]$.
The marked rim lengths are pairwise distinct.
If $a$ wheels have no marked rim and $b$ wheels have a marked rim, then
\[
 8a+9b\ge\left(\frac{49}{72}-o(1)\right)h.
\]
Every wheel has $O(h)$ edges, shortest cycle length at least $2\lceil\sqrt h\rceil$,
and at least $(r-1)h$ edges in every connected core of rank $r\in\{2,3,4\}$.
\end{lemma}

\begin{proof}
Put $s=\lceil\sqrt h\rceil$.
Use the spoke distribution from Lemma~\ref{lem:rhost-wheel-lengths}:
$U$ is uniform on $\{s,\ldots,2s\}$, the four deficits have sum $h+U$ and distribution
\eqref{eq:rhost-wheel-distribution}, and spoke $i$ has length $h+s-X_i$.
Independently choose one of two types.

With probability $2/3$, use the previous small rim $(1,L,1,M)$, where $L,M$ are independent and uniform on $\{s,\ldots,2s\}$.
Ignore its rim length and give the resulting twelve-element length set value eight.
With probability $1/3$, choose
\[
 R\ \text{uniform on }\{6s,\ldots,\lfloor h/36\rfloor\}.
\]
Choose $L$ uniformly from the integer interval with center $\lfloor(R-2)/2\rfloor$ and radius $\lfloor s/2\rfloor$.
Put $M=R-2-L$ and use the rim $(1,L,1,M)$.
Mark its rim length $R$ and give the thirteen-element length set value nine.
For sufficiently large $h$, these rim arms are positive.
The marked rim length is in $[3,h/3+1]$.

Discard outcomes with a nonpositive spoke, repeated non-rim cycle lengths, or a non-rim cycle outside $[h+1,2h]$.
For the small-rim type, the discarded probability is $o(1)$.
We estimate it for the large-rim type.
The first two causes still have probability $o(1)$.
The spoke calculation is unchanged.
For repeats, the codegree calculation below applies.

For a fixed spoke pair with rim-arc lengths $a,a'$, a cycle exceeds $2h$ only if
\[
 X_i+X_j<2s+\max\{a,a'\}.
\]
The pair deficit is uniform on $\{0,\ldots,h+U\}$.
The two pairs $01,23$ have largest rim arc $R-1$.
The other four pairs have largest rim arc at most $R/2+s/2+2$.
A union bound over these six pairs gives
\[
 \Pr(\text{some non-rim cycle exceeds }2h\mid R)
 \le\frac{4R+14s+6}{h}.
\]
Every non-rim cycle has length at least $h+1$.
Since $\mathbb E R=h/72+O(s)$, the retained probability for the large-rim type is at least $17/18-o(1)$.
Thus the expected retained value is at least
\begin{equation}
 8\cdot\frac23+9\cdot\frac13\left(\frac{17}{18}-o(1)\right)
 =\frac{49}{6}-o(1).
 \label{eq:rhost-wheel-value}
\end{equation}

For every upper-interval length, its weighted degree is at most $12/h$.
This follows from the uniform pair-deficit marginal, separately for each of the twelve labels.
For a marked rim length, the weighted degree is at most
\[
 \frac{1/3}{\lfloor h/36\rfloor-6s+1}
 =\frac{12+o(1)}h.
\]
All weighted codegrees are $O(h^{-3/2})$.
For two upper-interval lengths, the shared-index argument from Lemma~\ref{lem:rhost-wheel-lengths} is unchanged.
For complementary spoke pairs, the independent variable $U$ gives the same bound.
For two arcs of one spoke pair, their difference is smoothed either by $R$, or by the interval of $\Theta(s)$ choices for $L$.
For a rim length and an upper-interval length, first fixing $R$ costs $O(1/h)$.
Fixing the pair deficit then costs at most $1/h$.
Two distinct marked rim lengths never occur in the same outcome.
These bounds also show that repeated upper-interval labels have probability $o(1)$.

Apply Lemma~\ref{lem:rhost-weighted-matching} to these weighted sets, with values eight and nine.
The maximum weighted degree is $(12+o(1))/h$ and the weighted total value is at least $49/6-o(1)$.
It gives a disjoint collection with total value at least $(49/72-o(1))h$.
Choose one retained wheel metric for each selected set.

All chosen metrics satisfy the claimed density and traceability properties.
Indeed, their total spoke length is $3h+4s-U$, every spoke is at most $h+s$,
and every two spokes have total length at least $h$.
Both rim types have length at least $2s$ and at most $O(h)$.
The core classification and spanning-path deletion from Lemma~\ref{lem:rhost-wheel-lengths} apply verbatim.
\end{proof}

\begin{theorem}[An upper coefficient of $79/72$]\label{thm:rhost-seventy-nine}
For the same $p_*$ as in Theorem~\ref{thm:rhost-scale}, and every $p\ge p_*$,
with high probability
\[
 \operatorname{Pex}(G(n,p))
 \le\log_2n+\left(\frac{79}{72}+o(1)\right)\frac{\ln n}{\ln\ln n}.
\]
\end{theorem}

\begin{proof}
Use the exposure scheme of Theorem~\ref{thm:rhost-ten-ninths}, with $h=9\lceil d/9\rceil$.
Start with single-cycle blocks of lengths $3,\ldots,h/3+1$,
the theta partition for $[h/3+2,h-2]$, and single-cycle blocks of lengths $h-1,\ldots,2h$.
Their total rank is $16h/9+O(1)$.

Take the wheel list from Lemma~\ref{lem:rhost-wheel-profit}.
For each wheel, replace its twelve single-cycle blocks in the upper interval by that wheel.
If its rim is marked, also remove the single-cycle block of that rim length.
The disjointness of the chosen length sets makes all these replacements compatible.
A wheel has rank four, so an unmarked wheel saves eight and a marked wheel saves nine.
The resulting short-block rank is at most
\[
 \frac{16h}9-\left(\frac{49}{72}-o(1)\right)h
 =\left(\frac{79}{72}+o(1)\right)h.
\]
Add $o(d)$ single-cycle blocks for the remaining required lengths through $T=2(d+b+1)$.

The wheel overlap proof in Lemma~\ref{lem:rhost-wheel-embeddings} uses only the stated core thresholds,
the lower bound on the shortest cycle, and the $O(h)$ total edge bound.
It therefore applies to both rim types in the present list.
The theta embedding, connector chain, remaining exposures, and final excess count are unchanged.
Since $h=d+O(1)$ and $b,m=o(d)$, the asserted bound follows.
\end{proof}

% Audited research proof; stronger bounds appear in the following refinement.
% Include after random-host-weighted-matching.tex and random-hosts-phase5.tex.
\section{Keeping wheels with partially useful length sets}
\label{sec:rhost-one-fifty-one}

A wheel remains useful when some of its cycle lengths exceed the target interval.
Retain its useful lengths and charge its full rank.
This avoids the main loss in the preceding construction.

\begin{lemma}[Partial wheel length sets]\label{lem:rhost-partial-wheels}
Let $h\to\infty$ through multiples of nine.
There is a list of traceable wheel subdivisions with marked cycle-length sets $A$ such that:
\begin{enumerate}
\item each $A$ contains its rim length in $[3,h/3+1]$ and at least four non-rim lengths in $[h+1,2h]$;
\item the marked sets are pairwise disjoint;
\item $\sum_A(|A|-4)\ge(35/48-o(1))h$.
\end{enumerate}
Every wheel satisfies the same core thresholds and edge bounds as Lemma~\ref{lem:rhost-wheel-profit}.
\end{lemma}

\begin{proof}
Put $s=\lceil\sqrt h\rceil$ and use the previous spoke distribution.
Choose $R$ uniformly from $\{6s,\ldots,\lfloor h/12\rfloor\}$.
Choose $L$ uniformly from the integer interval with center $\lfloor(R-2)/2\rfloor$ and radius $\lfloor s/2\rfloor$.
Put $M=R-2-L$ and use the rim $(1,L,1,M)$.
Its marked rim length is $R$.
Retain only those non-rim cycle lengths that lie in $[h+1,2h]$.
Let $k$ be their number, counted initially by cycle labels.
The prospective value of this outcome is $k-3$.

Each spoke-pair deficit is uniform on $\{0,\ldots,h+U\}$.
For a rim arc of length $a$, its cycle exceeds $2h$ with probability at most $(2s+a)/h$.
The two arcs for each of the six pairs sum to $R$.
All non-rim cycles have length at least $h+1$.
Consequently,
\[
 \mathbb E(12-k)\le\frac{24s+6\mathbb E R}{h}
 =\frac14+o(1),
\]
and
\[
 \mathbb E(k-3)\ge\frac{35}{4}-o(1).
\]
Discard outcomes with a nonpositive spoke or repeated non-rim cycle lengths.
Their total probability is $o(1)$, by the calculations in the preceding section.
The value is bounded, so this changes its expectation by $o(1)$.
Also discard outcomes with $k\le3$.
Their values are nonpositive, so this cannot lower the retained expected value.
Each retained marked set has size $k+1\in\{5,\ldots,13\}$ and value $|A|-4\in\{1,\ldots,9\}$.

Every upper-interval vertex has weighted degree at most $12/h$.
Every rim vertex has weighted degree at most
\[
 \frac1{\lfloor h/12\rfloor-6s+1}=\frac{12+o(1)}h.
\]
The weighted codegree bound $O(h^{-3/2})$ is unchanged.
Restricting the marked set to its useful cycle lengths only decreases these degrees and codegrees.
For a rim and a non-rim length, fixing $R$ has probability $O(1/h)$ and fixing the pair deficit costs at most $1/h$.
Apply the variable-size version of Lemma~\ref{lem:rhost-weighted-matching}, with maximum value nine.
It supplies disjoint marked sets of total value at least
\[
 \frac{h}{12+o(1)}\left(\frac{35}{4}-o(1)\right)
 =\left(\frac{35}{48}-o(1)\right)h.
\]
Choose one metric witnessing each selected set.
The core bounds, total edge count, and spanning path follow exactly as in Lemma~\ref{lem:rhost-wheel-profit}.
\end{proof}

\begin{theorem}[An upper coefficient of $151/144$]\label{thm:rhost-one-fifty-one}
For the same $p_*$ as in Theorem~\ref{thm:rhost-scale}, and every $p\ge p_*$,
with high probability
\[
 \operatorname{Pex}(G(n,p))
 \le\log_2n+\left(\frac{151}{144}+o(1)\right)\frac{\ln n}{\ln\ln n}.
\]
\end{theorem}

\begin{proof}
Use the short-block baseline of Theorem~\ref{thm:rhost-seventy-nine}, of rank $16h/9+O(1)$.
For every marked set in Lemma~\ref{lem:rhost-partial-wheels}, replace its $|A|$ single-cycle blocks by the corresponding wheel.
This saves $|A|-4$ in cycle rank.
The disjointness of the marked sets makes these replacements compatible.
The resulting rank is at most
\[
 \frac{16h}{9}-\left(\frac{35}{48}-o(1)\right)h
 =\left(\frac{151}{144}+o(1)\right)h.
\]
The remaining $o(d)$ single-cycle blocks, simultaneous random embeddings, connector chain, and later exposures are unchanged.
All discarded cycle lengths belong to an embedded wheel; they need not belong to its marked set.
Every required length still has either its retained baseline block or a marked replacement.
The previous excess calculation proves the assertion.
\end{proof}

% Independently audited by root and four_regular; proof status is recorded in the shared handoff.
% Include after random-hosts-phase6.tex.
\section{Centering the wheel length intervals}
\label{sec:rhost-centered-wheels}

The useful-length construction improves further when each spoke is shortened by approximately one quarter of the rim length.
This centers the long cycle intervals around their required range.
Some outcomes then have invalid spokes; their probability is bounded explicitly below.

\begin{theorem}[An upper coefficient of $199/192$]\label{thm:rhost-centered-wheels}
For the same $p_*$ as in Theorem~\ref{thm:rhost-scale}, and every $p\ge p_*$,
with high probability
\[
 \operatorname{Pex}(G(n,p))
 \le\log_2n+\left(\frac{199}{192}+o(1)\right)\frac{\ln n}{\ln\ln n}.
\]
\end{theorem}

\begin{proof}
Use the distribution from Lemma~\ref{lem:rhost-partial-wheels}, with one change:
put
\[
 B=h-\lfloor R/4\rfloor+s,\qquad w_i=B-X_i.
\]
The total deficit remains $t=h+U$.
Before discarding invalid outcomes, the total edge length is
\[
 4B-t+R=3h+4s-U+R-4\lfloor R/4\rfloor\ge3h+2s.
\]

For a fixed rim arc of length $a$, the corresponding cycle length is uniform on
\[
 [2B-t+a,2B+a]\cap\mathbb Z.
\]
The probability that this length lies outside $[h+1,2h]$ is at most
\[
 \frac{|a-R/2|+O(s)}h.
\]
The constant in this estimate is absolute, uniformly over all choices.
For each of the two unit-rim pairs, the two arcs are $1,R-1$.
Their two absolute deviations sum to $R-2$.
For the other four pairs, both arcs differ from $R/2$ by $O(s)$.
Thus, if $k$ counts useful non-rim labels,
\[
 \mathbb E(12-k)\le\frac{2\mathbb E R+O(s)}h=\frac1{12}+o(1),
 \qquad \mathbb E(k-3)\ge\frac{107}{12}-o(1).
\]

We bound the probability $p_{\rm bad}$ that a spoke is nonpositive.
The marginal mass formula is
\[
 \Pr(X_i=t-r)=\frac{c_{t-r}(2r+1)c_r}{t+1}.
\]
Stirling's formula gives $c_j\sim(\pi j)^{-1/2}$.
Riemann sums, together with $c_j\le C/\sqrt{j+1}$ at the endpoints, show that $X_i/h$
converges in distribution to a variable with density
\[
 f(y)=\frac2\pi y^{-1/2}(1-y)^{1/2},\qquad 0<y<1.
\]
Also, $R/h$ converges to an independent uniform variable on $[0,1/12]$.
The independence follows because $R$ is independent of $U$ and the deficit distribution.
The limiting event $X_i/h\ge1-R/(4h)$ has a boundary of probability zero.
A union bound therefore gives
\[
 \limsup p_{\rm bad}
 \le 12\int_0^{1/12}\frac8\pi
       \int_0^{x/4}\frac{\sqrt u}{\sqrt{1-u}}\,du\,dx.
\]
For $0\le x\le1/12$,
\[
 \frac8\pi\int_0^{x/4}\frac{\sqrt u}{\sqrt{1-u}}\,du
 \le\frac{2x^{3/2}}{3\pi\sqrt{1-x/4}}
 \le\frac9{40}x^{3/2}.
\]
For the last inequality, use $\pi\ge3$ and $\sqrt{47/48}>80/81$.
Consequently,
\[
 p_{\rm bad}\le\frac9{100}\left(\frac1{12}\right)^{3/2}+o(1)
             =\frac3{800\sqrt3}+o(1).
\]

Discard outcomes with a nonpositive spoke or repeated non-rim cycle lengths.
The latter event still has probability $o(1)$.
Discarding a bad-spoke outcome loses at most nine in value.
Discard outcomes with $k\le3$, whose values are nonpositive.
The retained expected value is at least
\[
 \frac{107}{12}-\frac{27}{800\sqrt3}-o(1).
\]
The maximum weighted degree remains $(12+o(1))/h$.
Indeed, the pair-deficit interval still has $h+U+1$ equally likely positions, and the rim distribution is unchanged.
The weighted codegree remains $O(h^{-3/2})$.
The shared-index split proof is unchanged.
For complementary pairs, condition on $R,L$ and use the independent variable $U$.
For two arcs of one pair, use the same $R,L$ smoothing as before.
For a rim and a long length, fixing $R$ and then a pair sum costs $O(h^{-2})$.
Thus Lemma~\ref{lem:rhost-weighted-matching} supplies disjoint marked sets of total value at least
\[
 \left(\frac{107}{144}-\frac9{3200\sqrt3}-o(1)\right)h.
\]

Every retained metric has the required local density.
Its total length is at least $3h$ and its largest arm is at most $h+s$.
A three-spoke rank-three core therefore has length at least $2h$.
A four-spoke rank-three core omits at most one rim arm, of length at most $R/2+O(s)$, and also exceeds $2h$.
For a two-spoke rank-two core, include the full rim and use
\[
 2B-t+R\ge h+R/2.
\]
A rank-two core with at least three spokes has length at least $2h-R+s>h$.
Every non-rim cycle has length at least $h-R/2\gg\sqrt h$, while the rim has length at least $6s$.
The edge count is $O(h)$ and the same four-edge deletion gives a spanning path.
The wheel embedding estimate therefore applies.

Replace baseline single cycles by these marked wheel sets, as in Theorem~\ref{thm:rhost-one-fifty-one}.
The short-block rank is at most
\[
 \left(\frac{149}{144}+\frac9{3200\sqrt3}+o(1)\right)h.
\]
Since $\sqrt3>5/3$ and $27/16000<1/576$, the constant is less than $199/192$.
The previous exposure scheme and excess calculation finish the proof.
\end{proof}

The wheel checker also verifies 3000 centered metrics against all 42 kernel core types.
These finite checks support the metric formulas; the distributional limit and random matching arguments are proved above.

\section{Inputs for the threshold construction}
\label{sec:rhost-threshold-tools}

All graphs are simple.
A path's length is its number of edges.
The neighbourhood $N_J(A)$ is external: it contains vertices outside $A$ adjacent to $A$.
The cycle rank of a connected graph is $|E|-|V|+1$.
We write $\ln$ for the natural logarithm and $\log_2$ for the binary logarithm.
An event holds with high probability if its probability tends to one as $n\to\infty$.

\begin{lemma}[Uniform robust expansion]\label{lem:rhost-threshold-expansion}
Suppose $D\to\infty$, $D\le(\ln n)^B$ for a fixed $B$, and $\beta=20\ln D/D$.
Let $R\subseteq[n]$ have size $n-o(n)$ and expose $J\sim G(R,D/n)$.
For any prescribed $r_n=o(n)$, with high probability every $U\subseteq R$ of size at least $n-r_n$
contains a set $U'$ with
\[
 |U'|>|U|-\beta n,\qquad
 |N_{J[U']}(A)|\ge\frac{|A|}{3\beta}
 \quad(0<|A|\le\beta n).
\]
The assertion also holds conditionally on a random $R$ chosen independently of $J$.
\end{lemma}

\begin{proof}
Put $s=\lceil\beta n\rceil$.
The probability that two disjoint $s$-sets in $R$ have no joining edge is at most
\[
 \binom ns^2e^{-Ds^2/n}
 \le\exp\{s[2\ln(en/s)-Ds/n]\}=o(1).
\]
Indeed, $Ds/n\sim20\ln D$ whereas $\ln(en/s)\sim\ln D$.
Assume no such pair exists.
Starting from $U$, repeatedly remove a nonempty set of at most $\beta n$ vertices that violates the desired expansion.
If the total removed set $B$ first reaches size $\beta n$, then $\beta n\le|B|<2\beta n$.
Its external neighbourhood in the original $J[U]$ has size less than $|B|/(3\beta)<2n/3$.
Thus $U\setminus(B\cup N_{J[U]}(B))$ has at least $n/3-r_n-2\beta n$ vertices.
For large $n$, this exceeds $s$.
Taking $s$ vertices there and $s$ vertices in $B$ gives the forbidden pair.
Hence the deletion process removes fewer than $\beta n$ vertices and proves the assertion.
The probability estimate is uniform in $R$.
\end{proof}

The rooted tree theorem below is the form of Haxell's embedding result used in the construction.
It is stated as Theorem~2.2 by Alon and Krivelevich,
\href{https://doi.org/10.1137/23M1598969}{\emph{Sparse Pancyclic Subgraphs of Random Graphs}} (2025).
That statement follows from Haxell, \emph{Tree embeddings}, Journal of Graph Theory 36 (2001), 121--130,
and Balogh, Csaba, Pei, and Samotij,
\href{https://doi.org/10.37236/278}{\emph{Large bounded degree trees in expanding graphs}} (2010).

\begin{theorem}[Rooted tree embedding, external]\label{thm:rhost-rooted-tree}
Let $N,\Delta,k$ be positive integers.
Suppose a graph $J$ satisfies
\[
 |N_J(A)|\ge\Delta|A|+1\quad(0<|A|\le k),
\]
\[
 |N_J(A)|\ge\Delta|A|+N\quad(k<|A|\le2k).
\]
Then every tree with at most $N$ vertices and maximum degree at most $\Delta$
embeds in $J$ with any specified root mapped to any specified vertex of $J$.
\end{theorem}

\begin{lemma}[Trees with many leaves at one depth]\label{lem:rhost-threshold-trees}
Use $\beta$ from Lemma~\ref{lem:rhost-threshold-expansion} and set
\[
 a=\lfloor(5\beta)^{-1}\rfloor,\quad M=\lfloor n/1000\rfloor,
 \quad d=\lceil\log_a M\rceil.
\]
Every expander $J[U']$ supplied by that lemma contains, rooted at any prescribed vertex,
a tree with exactly $M$ leaves at depth $d$, fewer than $3M$ vertices, and maximum degree at most $a+1$.
It also contains a tree formed by joining the roots of two such trees with any positive-length $o(n)$ stem.
\end{lemma}

\begin{proof}
Take $k=\lfloor\beta n/2\rfloor$, $\Delta=\lfloor1/(4\beta)\rfloor$, and $N=\lfloor n/50\rfloor$.
For $0<|A|\le k$, the expansion surplus over $\Delta|A|$ is at least $|A|/(12\beta)>1$.
For $k<|A|\le2k$, it exceeds $n/24>N$.
Thus Theorem~\ref{thm:rhost-rooted-tree} applies.
For sufficiently large $n$, we have $a\ge3$ and $a+1\le\Delta$.

Build complete $a$-ary levels through depth $d-1$.
Give the $a^{d-1}$ vertices at that depth between one and $a$ children each, with exactly $M$ children in total.
This is possible because $a^{d-1}\le M\le a^d$.
All leaves have depth $d$.
The earlier levels have fewer than $2M$ vertices, so the whole tree has fewer than $3M$ vertices.
Two such trees, joined at their roots by a positive-length $o(n)$ stem, have fewer than $6M+o(n)<N$ vertices.
Their maximum degree is at most $a+1$.
Apply the rooted embedding theorem to each required tree.
\end{proof}

\begin{lemma}[A short list of disjoint cycles]\label{lem:rhost-cycle-list}
Fix constants $A,B>0$.
Suppose $D\to\infty$ and $D\le(\ln n)^B$.
In $G(n,D/n)$, with high probability, one can select vertex-disjoint cycles from any fixed list of $O(\ln n)$ lengths,
all between $3$ and $A\ln n$, provided each length has bounded multiplicity.
\end{lemma}

\begin{proof}
Let $X_\ell$ count cycles of length $\ell=O(\ln n)$ and let $Y_\ell$ count ordered intersecting pairs, including identical copies.
Write $\mu_\ell=\mathbb EX_\ell$.
We first establish
\[
 \operatorname{Var}X_\ell,\ \mathbb EY_\ell
 \le C\mu_\ell^2\left(\frac{\ell^2}n+\ell D^{-\ell}\right).
\]
Every proper overlap is a union of paths and isolated vertices.
An overlap with $u$ vertices and $f$ edges has normalized weight $(1+O(\ell^2/n))n^{-u}(n/D)^f$.
An isolated vertex has at most $\ell^2$ descriptions.
A path with $j\ge1$ edges has at most $4\ell^2$ descriptions, each of weight $n^{-1}D^{-j}$.
For several components, division by their order factorial bounds the total by the exponential of the one-component sum.
Thus all proper nonempty overlaps contribute $O(\ell^2/n)$.
The full overlap contributes $(1+O(\ell^2/n))2\ell D^{-\ell}$.
Vertex-disjoint pairs contribute at most $\mu_\ell^2$, proving the variance bound as well.

Put $q=\lfloor\ln n/(10\ln D)\rfloor$.
Chebyshev's inequality, summed over $3\le\ell\le q$, shows that every required short length has any fixed number of copies.
Indeed, $\sum_{\ell\ge3}\ell D^{-\ell}=O(D^{-3})$, and $\mu_\ell\to\infty$ uniformly.
With high probability, all distinct cycles of lengths at most $q$ are vertex-disjoint.
The same proper-overlap enumeration for two possibly different lengths bounds the expected number of intersecting pairs by
\[
 n^{-1}(\ln n)^{O(1)}D^{2q}=o(1).
\]
Choose all prescribed cycles in this range first.

For $q<\ell\le A\ln n$, put $\varepsilon_\ell=C(\ell^2/n+\ell D^{-\ell})\le n^{-1/10+o(1)}$.
Let $M_0=\lceil(\ln n)^3\rceil$.
Chebyshev's inequality gives $X_\ell\ge\mu_\ell/2$ except with probability $O(\varepsilon_\ell)$.
Markov's inequality gives $Y_\ell\le\mu_\ell^2/(8M_0)$ except with probability $O(M_0\varepsilon_\ell)$.
In the graph whose vertices are copies and whose edges join intersecting copies, the elementary independent-set bound is $X_\ell^2/Y_\ell$.
Thus at least $2M_0$ copies are pairwise vertex-disjoint.
A union bound gives this simultaneously for all required lengths.
The prescribed list uses only $O((\ln n)^2)$ vertices in total.
Greedily choose the remaining cycles: each previously used vertex excludes at most one member of any disjoint-copy family.
\end{proof}

The Hamilton-path input below is Lemma~2.1 in the cited Alon--Krivelevich paper.
It follows from the Hamiltonicity threshold and rotation-extension arguments; see Section~4 of
Krivelevich, \href{https://arxiv.org/abs/1507.00205}{\emph{Long paths and Hamiltonicity in random graphs}} (2016).

\begin{lemma}[Hamilton completion, external input]\label{lem:rhost-threshold-hamilton}
If $Nq-\ln N-\ln\ln N\to\infty$, then $G(N,q)$ has, with high probability,
a set $S$ of at least $N/4$ vertices such that every $s\in S$ has a set $T_s$ of at least $N/4$ vertices,
and each $t\in T_s$ is joined to $s$ by a Hamilton path.
Consequently, two external vertices with fresh independent $q$-edges to this graph can be joined through all $N$ vertices with high probability.
\end{lemma}

\begin{proof}[Proof of the consequence]
Expose the induced graph first.
The first external vertex has a neighbour $s\in S$ except with probability at most $e^{-qN/4}$.
Condition on that neighbour.
The second external vertex has a neighbour in $T_s$ with the same failure bound.
These two edges and the stated Hamilton path give the required path.
\end{proof}

\begin{lemma}[Long interval chords]\label{lem:rhost-threshold-windows}
Let a Hamilton cycle on $n$ vertices contain a distinguished arc of length
$2s-1+\ell-1$, where $s\gg n^{0.9}$ and $\ell=o(n^{0.9})$.
Suppose replacing paths inside this arc permits every length reduction from $0$ through $2s-1$.
Let $m=\lfloor n/s\rfloor\ge3$.
In an independent $G(n,D/n)$ with $D\to\infty$, with high probability $m-2$ additional edges
supply cycle intervals that, together with
\[
 [\ell,\ell+2s-1]\quad\text{and}\quad[n-2s+1,n],
\]
cover every integer from $\ell$ through $n$.
\end{lemma}

\begin{proof}
Write $W=n^{0.9}$ and let the distinguished arc have length $b_0=2s+\ell-2$.
For each $3\le i\le m$, there are at least $W^2/100$ vertex pairs whose joining arc contains the distinguished arc
and has length $a_i\in[is-W,is+W]$.
To see this, extend the two ends of the distinguished arc by distances $u,v\ge0$.
Put $t_i=\min\{is-b_0,n-b_0-W/2\}$.
Vary each distance independently in a window of width $W/4$ centered at $t_i/2$.
The lower bound $is-b_0\ge s-o(s)\gg W$ keeps both distances positive.
The upper adjustment leaves a complementary arc of length at least $W/4$.
All these pairs are distinct non-cycle edges and satisfy the required length window.

The probability that some window has no random edge is at most
\[
 m\exp(-Dn^{0.8}/100)=o(1).
\]
Choose one edge from each window.
It closes a cycle on the arc of length $a_i$, and the allowed reductions give
$[a_i+1-(2s-1),a_i+1]$.
Consecutive intervals overlap because $s\gg W$.
The first overlaps $[\ell,\ell+2s-1]$, and the last overlaps $[n-2s+1,n]$ because $ms>n-s$.
This proves the assertion.
\end{proof}

Finally, every connected pancyclic graph with $n+k$ edges has cycle rank $k+1$.
Its distinct cycles give distinct nonzero vectors in its binary cycle space.
Hence $n-2\le2^{k+1}-1$, so $k\ge\log_2(n-1)-1$.

% Independently audited by root, four_regular, and path_stability.
\section{Adaptive attachments at the full pancyclicity threshold}
\label{sec:rhost-full-threshold}

The endpoint failure probability need not be small enough for a direct union bound.
We instead select the attachment edges after exposing each connector tree.
Long base cycles supply spare candidate edges.
For the short cycles, the number of candidates grows with their length.

\begin{theorem}[The leading term at the full threshold]\label{thm:rhost-full-threshold}
Suppose
\[
 np-\ln n-\ln\ln n\longrightarrow\infty.
\]
Then with high probability
\[
 \operatorname{Pex}(G(n,p))=(1+o(1))\log_2 n.
\]
\end{theorem}

\begin{proof}
Put
\[
 w=\min\{np-\ln n-\ln\ln n,\ln\ln n\},\qquad D=w/100000.
\]
Expose five independent graphs with
\[
 p_1=p_2=p_3=p_5=D/n,\qquad
 p_4=(\ln n+\ln\ln n+w/2)/n.
\]
Their union is dominated by $G(n,p)$.
We have $D\to\infty$, without any specified rate.

Use
\[
 \beta=20\ln D/D,\quad a=\lfloor(5\beta)^{-1}\rfloor,\quad
 M=\lfloor n/1000\rfloor,\quad d=\lceil\log_a M\rceil.
\]
Thus $d\sim\ln n/\ln D=o(\ln n)$.
Lemmas~\ref{lem:rhost-threshold-expansion} and~\ref{lem:rhost-threshold-trees} supply the robust expanders and equal-depth leaf trees.
In particular, every permitted reservoir contains an induced expander of size $n-o(n)$.
It contains the required rooted trees and two-tree structures.
Each leaf tree has exactly $M$ leaves at depth $d$ and fewer than $3M$ vertices.
Write $H=\lfloor n/4000\rfloor$ and $\delta=\exp(-D/5000)=o(1)$.
A fresh endpoint has a neighbour in any specified $H$-vertex subset with failure at most $\delta$.
The same bound holds for a specified leaf set after at most $2H$ earlier endpoint queries.

Set
\[
 K_0=\lceil\log_2(4d)\rceil,\quad
 \ell_*=(K_0+2)(d+3),\quad b=\lceil\sqrt{2\ell_*}\rceil,
\]
\[
 c=8b+5,\qquad T=c+4d+3=4d+8b+8,
\]
\[
 K=\left\lfloor\log_2\frac{nw}{100\ln n}\right\rfloor,
 \qquad L_0=2^{K+1}-1,\qquad m=\lfloor n2^{-K}\rfloor.
\]
These parameters are fixed before exposure.
They satisfy
\[
 T,b,d,m=o(\ln n),\quad \ell_*=o(L_0),\quad L_0\gg n^{0.9},
 \quad 2L_0\le\frac{nw}{25\ln n}.
\]
Also, $\ell_*/K\to\infty$, since $D\le\ln\ln n$ and $\ln d\sim\ln\ln n$.

\emph{Initial cycles.}
In $G_1$, select vertex-disjoint bit cycles $C_i$ of lengths $2^i+2$, for $0\le i\le K_0$.
Select a digit base cycle $C_{\rm short}$ of length $c$.
Select one further cycle $S_j$ of each length $3\le j\le T$.
Lemma~\ref{lem:rhost-cycle-list} applies: all lengths are $O(\ln n)$, and their multiplicities are at most three.
The selected cycles use $O(T^2+d)$ vertices.
Also select an edge $f$ outside their vertices.
Such an edge exists with high probability: $G_1$ has $\Omega(nD)$ edges and maximum degree at most $(\ln n)^2$,
while the selected vertices number only $O((\ln n)^2)$.
Let $H_1$ contain all these cycles and $f$.

\emph{An adaptive chain of cycles.}
We describe a connector procedure in $G_2$.
Expose first only edges with both endpoints outside $H_1$.
The universal robust-expansion event then determines every required reservoir expander.
Edges from vertices of $H_1$ to the reservoir remain independent until queried.

Start at a specified endpoint outside the reservoir.
Query at most $H$ fresh candidate neighbours in the current expander.
If a neighbour $z$ exists, embed a leaf tree rooted at $z$.
This first step fails with probability at most $\delta$.
Suppose the next object is an untouched cycle $Q$ of length $j$.
Take a matching of $\lfloor j/2\rfloor$ candidate edges in $Q$.
For a candidate edge $xy$, first query the edges from $x$ to the previous leaf set.
If one exists, it gives an incoming path of length $d+2$ from the preceding endpoint to $x$.
Temporarily exclude this path's internal vertices from the next reservoir.
Then query at most $H$ fresh edges from $y$ to the resulting expander.
If this succeeds, retain the incoming path, select $xy$, and embed the next leaf tree at the chosen neighbour of $y$.
If either query fails, discard that candidate edge and retain no part of its tentative incoming path.
The previous leaf tree remains available for the next candidate.

Each candidate has two endpoints unused in any previous candidate.
The reservoir was exposed before all these endpoint queries.
Even after conditioning on earlier failures and on the tentative incoming path, the edges from the new endpoint $y$ remain fresh.
Thus each conditional candidate failure probability is at most $2\delta$.
A cycle of length $j$ fails with probability at most $(2\delta)^{\lfloor j/2\rfloor}$.
For any list with bounded multiplicity of cycle lengths, their total failure probability is at most
\[
 O\!\left(\sum_{j\ge3}(2\delta)^{\lfloor j/2\rfloor}\right)=o(1).
\]
At the end, one fixed target endpoint attaches to the final leaf set with failure at most $\delta$.
Only the chosen paths, not the unused tree vertices or failed queries, enter the constructed subgraph.

Apply this procedure first to the bit cycles, starting at one endpoint of $f$ and ending at the other.
The selected edge of $C_i$ is denoted $e_i$.
Together with $f$ and the $K_0+2$ connectors, these edges form a cycle $C_*$ of length $\ell_*$.
Each other arc of $C_i$ replaces $e_i$ with gain exactly $2^i$.

Fix an edge $e_{\rm short}$ of $C_{\rm short}$.
Apply the procedure again to $S_3,\ldots,S_T$.
Start at the endpoint of $f$ used for the first bounded query, and end at one endpoint of $e_{\rm short}$.
The starting endpoint has only $H$ earlier queries, so another $H$ fresh candidate edges remain available.
This second chain is disjoint from the first construction except at its designated endpoint.
It contains all single-cycle blocks needed for the short lengths.

\emph{Adaptive binary and digit detours.}
Expose the $G_3$ reservoir edges outside the entire graph constructed so far.
All its base-cycle endpoint edges remain fresh.
For each $i>K_0$, embed two depth-$d$ leaf trees whose roots are joined by a stem of length
\[
 2^i+1-2d-2=2^i-2d-1>0.
\]
The total tree size is below $n/50$.
Choose an unused candidate edge $xy$ of $C_*$, and query its two endpoints against the two leaf sets.
If both attachments succeed, the resulting detour has length $2^i+1$.
Its replacement gain over $xy$ is $2^i$.
If an attachment fails, abandon that candidate edge and try another with two fresh endpoints.

Preselect $2(K-K_0)$ pairwise vertex-disjoint candidate edges of $C_*$,
avoiding $f$ and all endpoints of $e_0,\ldots,e_{K_0}$.
There are enough because $\ell_*/K\to\infty$.
Every conditional failure probability is at most $2\delta$.
The expected number of failures in this entire candidate pool is at most $4\delta(K-K_0)$.
If the pool is exhausted before completing all detours, at least $K-K_0$ candidates have failed.
Markov's inequality shows that this event has probability at most $4\delta=o(1)$.
After each success, retain only its detour path and re-form the next reservoir.

The digit base cycle has at least $4b$ pairwise vertex-disjoint candidate edges avoiding the endpoints of $e_{\rm short}$.
Use them to obtain $2b$ successful detours, in two groups of $b$.
The desired detour lengths are
\[
 2d+3+j,\qquad 2d+3+jb\qquad(0\le j<b).
\]
The two-tree stems have lengths $1+j$ and $1+jb$, so they are positive.
The same expected-failures argument succeeds with probability $1-o(1)$.
The designated edges need not be consecutive on the base cycle.
Choosing one detour from each group gives every cycle length in
\[
 [T+1,T+b^2].
\]
Indeed, the length is $c-2+(2d+3+j_0)+(2d+3+bj_1)=T+1+j_0+bj_1$.
This interval overlaps $[\ell_*,\ell_*+L_0]$, since $T=o(\ell_*)$ and $b^2\ge2\ell_*$.

\emph{Hamilton completion and final coverage.}
Delete $f$, $e_{\rm short}$, and the selected edge of each $S_j$ when choosing a spanning path through the constructed blocks.
Use every bit detour, every digit detour, and the long arc of each $S_j$ in this path.
The separate short-cycle chain joins one end of the binary path to one end of the digit path.
Thus all constructed vertices lie on one spanning path.
The binary path is a contiguous segment of it.
All deleted edges remain present as chords in the constructed subgraph.

The total used vertex count is at most $2L_0$.
The binary detours contribute $L_0+O(\ell_*)$ vertices.
The short-cycle chain and digit gadget contribute only a polynomial in $\ln n$.
For the remaining vertex set $X$,
\[
 |X|p_4-\ln|X|-\ln\ln|X|
 \ge(1/2-1/25-o(1))w\longrightarrow\infty.
\]
Lemma~\ref{lem:rhost-threshold-hamilton} therefore joins the two remaining endpoints through all vertices of $X$.
This is the completion input from Alon--Krivelevich Lemma~2.1 and Claim~4.7.
Its endpoint failure probability is $o(1)$.
This gives a Hamilton cycle.

The bit cycles give every length in $[\ell_*,\ell_*+L_0]$.
Replacing bit detours on the Hamilton cycle gives every length in $[n-L_0,n]$.
The individual cycles $S_j$ and the digit gadget cover all lengths below the first interval.
Use $G_5$ only for the long interval-covering chords.
Lemma~\ref{lem:rhost-threshold-windows} applies because the binary segment is contiguous, has length $L_0+\ell_*-1$,
and $\ell_*=o(n^{0.9})=o(L_0)$.
Each required window has at least $n^{1.8}/100$ candidate edges.
The failure probability is at most $m\exp(-Dn^{0.8}/100)=o(1)$.
Their overlapping length intervals cover the remaining gap.

The constructed graph has excess at most
\[
 K+2b+T+m+O(1)=\log_2 n+o(\ln n).
\]
The short-cycle chain contributes $T-2$ to cycle rank.
Failed queries contribute no retained edges.
The universal cycle-count lower bound gives the asserted equality to leading order.
\end{proof}

\begin{corollary}[A quantitative full-threshold error]\label{cor:rhost-full-threshold-error}
Write $\omega=np-\ln n-\ln\ln n\to\infty$ and $w=\min\{\omega,\ln\ln n\}$.
Then with high probability
\[
 \operatorname{Pex}(G(n,p))
 \le\log_2n+(4+o(1))\frac{\ln n}{\ln w}.
\]
\end{corollary}

\begin{proof}
The preceding construction has $\ln a=(1+o(1))\ln w$, so $d\sim\ln n/\ln w$.
Moreover, $b/d=O(\sqrt{K_0/d})=o(1)$ and $m/d=O(\ln w/w)=o(1)$.
Therefore $T=4d+o(d)$, and the displayed excess bound gives the result.
\end{proof}

\section{A limitation of cycle, theta, and $K_4$ block constructions}
\label{sec:rhost-block-barrier}

This section concerns a construction method, not a lower bound for arbitrary pancyclic subgraphs.
The method supplies its short cycle lengths inside separate cycle, theta, and subdivided-$K_4$ blocks.
Its cost is the sum of the cycle ranks of those blocks.
The following obstruction shows that this method cannot attain coefficient one under the stated density constraints.

\subsection{An elementary obstruction to the proposed balanced mixture}

\begin{lemma}\label{lem:rhost-excess-length}
Let $d>0$.
For a subdivision of $K_4$ with total arm length $E\ge2d$, suppose every arm has length at most $E-d$.
Then
\[
 \sum_C(|C|-d)_+\le4E-6d.
\]
For a theta graph of total length $E\ge d$,
\[
 \sum_C(|C|-d)_+\le2E-2d.
\]
Here $x_+=\max\{x,0\}$, and each sum is over all cycles of the stated graph.
\end{lemma}

\begin{proof}
Each arm of $K_4$ belongs to four of its seven cycles.
It suffices to prove $\sum_C\min\{|C|,d\}\ge6d$.
The left side is a concave function of the six nonnegative arm lengths.
For fixed $E\ge2d$, minimise it on the polytope
\[
 \sum_iw_i=E,\qquad 0\le w_i\le E-d.
\]
Its vertices have two nonzero coordinates, equal to $E-d$ and $d$.
When $E=2d$, both coordinates equal $d$.
Any two distinct edges of $K_4$ occur together in exactly two cycles.
They each occur in four cycles, so their union meets exactly six cycles.
At each polytope vertex, the sum of the truncated cycle lengths is therefore $6d$.
Concavity proves the bound throughout the polytope.

For a theta graph, every arm belongs to two cycles.
On the simplex of three nonnegative arm lengths summing to $E\ge d$,
concavity gives $\sum_C\min\{|C|,d\}\ge2d$.
Indeed, at every simplex vertex, the two nonzero cycle lengths both equal $E$.
Subtract these truncated sums from the full cycle-length sums.
\end{proof}

Consider the proposed mixture of $N+O(1)$ subdivided-$K_4$ blocks and $N+O(1)$ theta blocks.
Let $d=5N$, and suppose their arm lengths are $O(N)$.
Require the $K_4$ blocks to have total length at least $10N$ and every rank-two core at least $5N$.
Require each theta block to have total length at least $5N$.
Suppose their cycles cover all lengths $3,\ldots,10N$.
There are only $10N+O(1)$ cycle occurrences, so all but $O(1)$ must supply distinct required lengths.
Consequently, their total cycle-length sum is $50N^2+O(N)$.
Lemma~\ref{lem:rhost-excess-length} bounds their total excess above $5N$ by
\[
 50N^2-6(5N)N-2(5N)N+O(N)=10N^2+O(N).
\]
But the required lengths have excess sum
\[
 \sum_{\ell=5N+1}^{10N}(\ell-5N)
 =\frac{25}{2}N^2+O(N),
\]
a contradiction for large $N$.
This excludes arbitrary choices and parameterisations of these blocks.
It is not restricted to affine families or to a particular parity assignment.

\subsection{A quantitative barrier for arbitrary mixtures}

Define a nonnegative weight function on $[0,\infty)$ by
\[
 w(x)=
 \begin{cases}
 1,&0\le x\le1/3,\\
 4/3-x,&1/3\le x\le1,\\
 1/2-x/6,&1\le x\le2,\\
 0,&x>2.
 \end{cases}
\]
Its two formulas agree at $1/3$ and $1$.
At $2$, use the third formula.

\begin{lemma}[Finite rational certificate]\label{lem:rhost-block-weight}
Let $Q$ be a subdivision of $K_4$ with total length at least $2d$ and every rank-two core of length at least $d$.
Then
\[
 \sum_{C\subseteq Q}w(|C|/d)\le3.
\]
For a theta graph $Q$ of total length at least $d$, the corresponding sum is at most two.
For a single-cycle block, it is at most one.
\end{lemma}

\begin{proof}
The single-cycle assertion follows from $0\le w\le1$.
The other assertions have finite rational certificates, described explicitly below.
Scale the arm lengths by $d$ and write them as a nonnegative vector $x$.
For $K_4$, the metric inequalities are
\[
 \sum_{i=1}^6x_i\ge2,\qquad
 \sum_{i\ne j}x_i\ge1\quad(1\le j\le6).
\]
For a theta graph, the metric inequality is $x_1+x_2+x_3\ge1$.

Assign each cycle length to one of the four intervals in the definition of $w$.
There are $4^7=16384$ assignments for $K_4$, and $4^3=64$ for a theta graph.
Each assignment gives rational linear inequalities $Ax\le b$ and a linear objective $c\cdot x+c_0$.
The objective is the sum of cycle weights on that region.
For every assignment, the supplied certificate gives a rational vector $y$ satisfying
\[
 y\ge0,\qquad A^{\mathsf T}y\ge c,\qquad
 b\cdot y+c_0\le r,
\]
where $r=3$ for $K_4$ and $r=2$ for a theta graph.
Since $x\ge0$, these inequalities imply
\[
 c\cdot x+c_0\le y^{\mathsf T}Ax+c_0\le b\cdot y+c_0\le r.
\]
Every permitted metric lies in a listed region.
At the discontinuity $x=2$, use the region with the nonzero boundary weight.
Thus the certificates prove the claims for all real nonnegative arm lengths.
\end{proof}

The file \nolinkurl{random-host-rank3-barrier-certificates.json} contains the rational multipliers.
The standard-library verifier \nolinkurl{check_random_host_rank3_barrier.py} checks every inequality exactly.
Its optional generation mode uses a linear-programming solver only to find candidate multipliers.
Every candidate is converted to rational numbers and checked exactly before it is accepted.
The proof does not use floating-point solver success as a certificate.

\begin{proposition}[A barrier of $37/36$]\label{prop:rhost-block-barrier}
Suppose a collection of single-cycle, theta, and subdivided-$K_4$ blocks supplies all lengths $3,\ldots,\lfloor2d\rfloor$ internally.
Suppose each theta block has total length at least $d$.
Suppose each $K_4$ block has total length at least $2d$ and each of its rank-two cores has length at least $d$.
Then the sum of the block cycle ranks is at least
\[
 \frac{37}{36}d-O(1).
\]
\end{proposition}

\begin{proof}
The weight function is nonnegative.
Select one block cycle of each required length and apply Lemma~\ref{lem:rhost-block-weight} to all blocks.
Their rank sum is at least
\[
 \sum_{\ell=3}^{\lfloor2d\rfloor}w(\ell/d)
 =d\int_0^2w(x)\,dx+O(1).
\]
The function has bounded variation, so the Riemann-sum error is bounded independently of $d$.
Finally,
\[
 \int_0^2w(x)\,dx
 =\frac13+\frac49+\frac14=\frac{37}{36}.
\]
\end{proof}

The obstruction permits arbitrary numbers, metrics, and parameterisations of the three block types.
It also permits duplicate and unused cycle lengths.
It does not cover constructions using larger kernels or short cycles that traverse several blocks.
Thus it supports the search for wheel blocks without excluding coefficient one for the original random-host problem.


\section{The exact four-regular theorem}

The full proof is in the companion manuscript \nolinkurl{four-regular-theorem.pdf}, with source and certificates beside it.
The following statements summarize its conclusions.
Parallel edges are distinct, and two parallel edges form a cycle.
For order two, a doubled cycle means four parallel edges joining two vertices.

\begin{theorem}
Every finite loopless four-regular multigraph on $n$ vertices has at most $2^n+n$ cycles.
Equality holds precisely for doubled cycles and $K_5$.
Moreover, such a graph has more than $2^n$ cycles if and only if it is a doubled cycle, $K_5$, or $K_4$ with two opposite edges doubled.
\end{theorem}

The last exceptional graph has $19$ cycles: two digons, eight triangles, and nine Hamilton cycles.
Complementation first reduces the problem to unions of two Hamilton cycles.
Doubled-edge expansion and gluing across two doubled pairs preserve the bound.
A parity count then reduces any remaining case to at most two matching doubled pairs.
The published degree-sequence path theorem, used after one or two vertex deletions, puts every such graph of order at least eleven strictly below $2^n$.
An induction verifies the analytic estimate at all larger orders from seven exact rational base inequalities.

The remaining finite check consists of $925$ matching instances on at most ten vertices, generated from $3234$ underlying simple graphs.
Every instance was counted independently by vertex-simple path enumeration and binary cycle-space enumeration.
Only $K_5$ attains $2^n+n$ in this reduced class.
The independent comparison was also run on all $5012$ reduced instances through order eleven.

\begin{corollary}
Every connected loopless Eulerian multigraph of rank $r\ge1$ has at most
\[
 2^{r-1}+r-1
\]
cycles.
For $r\ge3$, equality holds precisely for subdivisions of doubled cycles and of $K_5$.
The graphs with more than $2^{r-1}$ cycles are precisely the subdivisions of the three four-regular families in the preceding theorem.
At rank one the maximum is one, and at rank two it is two.
\end{corollary}

The companion paper gives the full reduction, equality analysis, source citations, and reproduction instructions.
The small-core certificate is \nolinkurl{four-regular-theorem-certificate.json}.
It can be regenerated with \nolinkurl{verify_four_regular_theorem.py} after building the nauty generator and the supplied C++ checker.
The theorem uses computer assistance for this finite check; it is not a formal Lean theorem.

% Research fragment. The general edge-half problem remains open here.
\section{Partial results for the cubic edge-half problem}
\label{sec:cubic-edge-frequency}

Write $c_e(G)$ for the number of cycles containing an edge $e$.
The edge-half problem asks for an edge with $c_e(G)\le c(G)/2$.
The two usual exceptions among connected simple cubic graphs are $K_4$ and $K_{3,3}$.
The results below settle several structural cases, but do not settle the general problem.
No priority claim is made for these elementary reductions.

\begin{proposition}\label{prop:cubic-small-structures}
Let $G$ be a connected simple cubic graph other than $K_4$ or $K_{3,3}$.
If $G$ has a bridge, a two-edge cut, a triangle, or a four-cycle, then some edge satisfies
\[
 c_e(G)\le\frac12c(G).
\]
\end{proposition}

\begin{proof}
A bridge lies in no cycle.
We may therefore suppose that $G$ is bridgeless.

First suppose that $\delta(S)$ is a two-edge cut.
Both induced shores $N_1,N_2$ are connected.
Otherwise one shore component has a cut of size at most one.
On each shore, the two cut endpoints are distinct.
If they coincided, the third edge at that endpoint would be a bridge.
The terminals have degree two within their shore.

Let $p_i$ count terminal paths in $N_i$, and put $c_i=c(N_i)$.
Every cycle crossing the cut uses a terminal path on each shore, so
\[
 c(G)=c_1+c_2+p_1p_2.
\]
Relabel the shores so that $c_1\le c_2$.
At one terminal of $N_1$, let $f,g$ be its two internal edges.
The terminal paths partition according to their first edge.
Choose $f$ to occur in at most $p_1/2$ terminal paths.
At most $c_1$ cycles wholly inside $N_1$ contain $f$.
Consequently,
\[
 c_f(G)\le c_1+\frac12p_1p_2\le\frac12c(G).
\]

Next suppose that $W$ is a triangle.
Put $c_0=c(G-V(W))$.
Let $S$ count paths in $G-E(W)$ joining distinct vertices of $W$.
Their internal vertices lie outside $W$, since every vertex of $W$ has outside degree one.
Every cycle meeting both $W$ and its exterior is one such path plus one of the two corresponding triangle arcs.
Thus
\[
 c(G)=c_0+1+2S,\qquad c_e(G)=1+S\quad(e\in E(W)).
\]
If $n=|V(G)|\ge6$, the exterior has $n-3$ vertices and $3n/2-6$ edges.
It therefore contains a cycle, since its edge count is at least its vertex count.
Hence $c_0\ge1$, and every triangle edge satisfies the required inequality.
The only simple cubic graph of order four is $K_4$.

Finally suppose that $W$ is a four-cycle.
A chord gives a triangle, so assume that $W$ is induced.
We claim the exact identity
\begin{equation}\label{eq:cubic-square-frequency}
 \sum_{e\in E(W)}c_e(G)=2c(G)+2-2c(G-V(W)).
\end{equation}
To prove it, group cycles by their edge sets outside $E(W)$.
A cycle with two attachment vertices has one outside path and two possible rim completions.
The two completed cycles use four rim edges in total.
A cycle with four attachment vertices has two disjoint outside paths.
Every connected completion uses a perfect matching of the four-cycle, and therefore uses exactly two rim edges.
The cycle $W$ itself uses four rim edges.
A cycle avoiding $E(W)$ avoids $V(W)$, since each rim vertex has outside degree one.
These cases prove~\eqref{eq:cubic-square-frequency}.

If $n\ge8$, the exterior has $n-4$ vertices and $3n/2-8$ edges.
Its edge count is at least its vertex count, so it contains a cycle.
Equation~\eqref{eq:cubic-square-frequency} gives
\[
 \min_{e\in E(W)}c_e(G)
 \le\frac14\sum_{e\in E(W)}c_e(G)\le\frac12c(G).
\]
For $n=6$, the complement of $G$ is a two-regular simple graph.
It is either a six-cycle or two triangles.
The corresponding cubic graphs are the triangular prism and $K_{3,3}$.
The former has a triangle and the latter is excluded.
\end{proof}

Thus any unresolved connected simple cubic example is three-edge-connected and has girth at least five.
The square identity also gives the quantitative bound
\[
 \min_{e\in E(W)}c_e(G)
 \le\frac12c(G)-\frac12\bigl(c(G-V(W))-1\bigr).
\]

\subsection{A weighted consequence of the four-regular classification}

\begin{proposition}\label{prop:cubic-matching-weight}
Let $G$ be a connected simple cubic graph on $n\ge6$ vertices.
For every perfect matching $M$ of $G$,
\begin{equation}\label{eq:cubic-matching-weight}
 \sum_{C\text{ a cycle of }G}2^{|E(C)\cap M|}\le2^n-\frac n2.
\end{equation}
If $G$ is bridgeless, then
\begin{equation}\label{eq:cubic-length-weight}
 \sum_{C\text{ a cycle of }G}2^{|E(C)|/3}\le2^n-\frac n2.
\end{equation}
\end{proposition}

\begin{proof}
Double every edge of $M$, obtaining a loopless four-regular multigraph $Q$.
Its cycles of length two are exactly the $n/2$ doubled matching edges.
Each longer cycle projects to one cycle $C$ of $G$ and has $2^{|E(C)\cap M|}$ lifts.
Hence
\[
 c(Q)=\frac n2+\sum_C2^{|E(C)\cap M|}.
\]
The four-regular classification gives $c(Q)\le2^n$.
Indeed, its underlying graph is cubic, so $Q$ is neither a doubled cycle nor $K_5$.
The remaining exceptional graph has order four.
This proves~\eqref{eq:cubic-matching-weight}.

Suppose now that $G$ is bridgeless.
Every odd cut has odd size and has at least three edges.
The vector assigning weight $1/3$ to every edge satisfies the perfect matching polytope constraints.
Edmonds' theorem therefore gives a probability distribution on perfect matchings with edge marginals $1/3$.
For a fixed cycle $C$, convexity gives
\[
 \mathbb E_M 2^{|E(C)\cap M|}
 \ge2^{\mathbb E_M|E(C)\cap M|}=2^{|E(C)|/3}.
\]
Average~\eqref{eq:cubic-matching-weight} and sum this inequality over the cycles.
This proves~\eqref{eq:cubic-length-weight}.
\end{proof}

The external input is J.~Edmonds, \emph{Maximum matching and a polyhedron with $0,1$-vertices},
Journal of Research of the National Bureau of Standards, Section~B \textbf{69B} (1965), 125--130,
\url{https://doi.org/10.6028/jres.069b.013}.
The original paper is available at
\url{https://nvlpubs.nist.gov/nistpubs/jres/69B/jresv69Bn1-2p125_A1b.pdf}.
The weighted inequalities do not by themselves establish the edge-half property.

\subsection{Exact formulas across a three-edge cut}

Let a nontrivial three-edge cut split a three-edge-connected simple cubic graph into connected shores $N_1,N_2$.
The three cut endpoints on each shore are distinct.
Otherwise an internal edge at a repeated endpoint, together with the remaining cut edge, gives a two-edge cut.
Label corresponding terminals by $1,2,3$.
Put $c_i=c(N_i)$, and let $p_{jk}$ and $q_{jk}$ count terminal paths in the two shores.
Then
\begin{equation}\label{eq:cubic-three-cut}
 S=\sum_{1\le j<k\le3}p_{jk}q_{jk},\qquad
 c(G)=c_1+c_2+S.
\end{equation}
The cut edge at terminal $j$ belongs to
\[
 \sum_{k\ne j}p_{jk}q_{jk}
\]
cycles, with subscripts interpreted as unordered pairs.
Hence some cut edge satisfies the edge-half bound whenever
\[
 S\le3(c_1+c_2).
\]
For an internal edge $e$ of $N_1$, let $p_{jk}(e)$ count terminal paths using $e$.
Its exact cycle frequency is
\begin{equation}\label{eq:cubic-three-cut-internal}
 c_e(G)=c_e(N_1)+\sum_{j<k}q_{jk}p_{jk}(e).
\end{equation}

One sufficient extension would be the following weighted three-terminal statement.
For every connected simple graph $N$ with three degree-two terminals and all other vertices cubic, and every $w_{jk}\ge0$, some edge $e$ would satisfy
\begin{equation}\label{eq:cubic-three-pole-candidate}
 2c_e(N)+2\sum_{j<k}w_{jk}p_{jk}(e)
 \le c(N)+1+\sum_{j<k}w_{jk}p_{jk}.
\end{equation}
This statement is not proved here.
If it holds, apply it with $w_{jk}=q_{jk}$.
The other shore contains a cycle, since it has $(3|V(N_2)|-3)/2$ edges and odd order at least three.
Thus $c_2\ge1$, and~\eqref{eq:cubic-three-cut-internal} would give an edge with frequency at most $c(G)/2$.

A simpler terminal-visitation bound is false.
Take the cubic graph with vertex set $\{0,\ldots,13\}$ and edge list
\[
\begin{split}
 &(0,5),(1,6),(2,7),(0,8),(1,8),(3,8),(0,9),\\
 &(2,9),(4,9),(2,10),(5,10),(6,10),(1,11),(4,11),\\
 &(7,11),(3,12),(4,12),(6,12),(3,13),(5,13),(7,13).
\end{split}
\]
It is three-edge-connected and has girth five.
After deleting vertex $2$, there are $50$ paths between terminals $7$ and $9$.
Exactly $26$ visit the third terminal $10$.
Thus even this restricted class does not guarantee that at most half the terminal paths visit the third terminal.
This does not disprove~\eqref{eq:cubic-three-pole-candidate}.

\subsection{A counterexample to the mean-length route}

A sufficient condition for the edge-half property would be mean cycle length at most $3n/4$.
This condition fails even at order eight.
Replace one vertex of $K_{3,3}$ by a triangle, attaching its three old incident edges to distinct triangle vertices.
The resulting simple cubic graph has cycle polynomial
\[
 z^3+3z^4+6z^5+6z^6+6z^7+6z^8.
\]
Indeed, $K_{3,3}$ has nine four-cycles and six six-cycles.
The replaced vertex belongs to six cycles of each length.
Each such cycle has two lifts, whose lengths increase by one and two.
The remaining three four-cycles persist, and the new triangle contributes one cycle.
The mean cycle length is therefore
\[
 \frac{171}{28}>6=\frac{3\cdot8}{4}.
\]
Nevertheless, every new triangle edge belongs to exactly $13$ of the $28$ cycles.
Thus the mean-length condition is stronger than the edge-half conclusion and cannot prove it in general.

\subsection{Independent finite checks}

The script \texttt{check\_cubic\_edge\_frequency.py} enumerates connected simple cubic graphs using nauty's \texttt{geng}.
It counts cycles as connected nonzero vectors of the binary cycle space.
For all $621$ such graphs through order fourteen, it verifies every triangle and induced-square identity above.
It also verifies~\eqref{eq:cubic-matching-weight} for all $8\,140$ perfect matchings in the graphs of order at least six.
The order-eight counterexample is checked separately, including its cycle polynomial and three triangle-edge frequencies.
The exact output is saved in \texttt{cubic-edge-frequency-results.json}.
These checks support the arguments but are not premises of the propositions.

\section{Other directions and the remaining gaps}\label{sec:remaining}

\subsection{An edge-frequency consequence}

The following implication was proposed in the supplied research list.
It follows from the source manuscript's formalized forest estimate.
It is recorded here with its denominator and graph assumptions.

\begin{proposition}\label{prop:research-edge}
Let $G$ be connected and simple, with sufficiently large rank $r$.
Suppose an edge $e$ lies on a cycle $W$ of length at most $\sqrt r$.
Then
\[
 c_e(G)\le\left(\frac14+\frac{10}{\log_2\log_2\log_2 r}\right)2^r.
\]
If also $c(G)\ge(1/2-\eta)2^r$ with $0\le\eta<1/2$, then
\[
 \frac{c_e(G)}{c(G)}
 \le\frac{1/4+10/(\log_2\log_2\log_2 r)}{1/2-\eta}.
\]
\end{proposition}

\begin{proof}
Restriction outside $W$ has fibers of size two, namely $\{x,x\mathbin\triangle E(W)\}$.
Exactly one member of each fiber contains $e$.
A cycle containing $e$ has a forest outside part.
Thus $c_e(G)\le p(G,W)2^{r-1}$.
Apply the source manuscript's forest theorem and divide by the stated lower bound on $c(G)$.
\end{proof}

A connected simple graph of minimum degree three has at most $2r-2$ vertices.
A breadth-first search gives a cycle of length $O(\log r)$.
Hence such a graph has an edge satisfying the proposition for large $r$.
For families with $c(G)/2^r\to1/2$, this gives an edge in at most $(1/2+o(1))c(G)$ cycles.
This is not the exact edge-half theorem for arbitrary graphs.

At a vertex $v$ of degree $d\ge4$, the exact local half bound is elementary:
\[
 \sum_{e\ni v}c_e(G)=2c_v(G)\le2c(G).
\]
Some incident edge therefore has $c_e(G)\le c(G)/2$.
A bridge lies in no cycle.
Thus a global counterexample with minimum degree at least three would have to be cubic and bridgeless.

\subsection{Why the general error-term questions remain}

The extended marked-set range in Section~\ref{sec:research-path} concerns a fixed positive error.
It does not supply polynomial decay in the error as $r$ grows.
The source proof's weighted cycle supply remains logarithmic in its available length ratio.
The moment-only example in Section~\ref{sec:inverse} prevents stronger conclusions from those moments alone.

The original chord construction gives
\[
 \frac{c(G)}{2^r}\le p(G^+,W)+2^{-t}.
\]
Its witness can have $t$ components.
The current forest estimate allows only $O(\log r)$ components.
Consequently, this direct use of the construction cannot make $2^{-t}$ exponentially small in $r$.
Fable's edge-deletion recurrence removes this separate error: for an edge $e$ of a short cycle,
\[
 c(G)=c(G-e)+c_e(G),\qquad c_e(G)\le p(G,W)2^{r-1}.
\]
It uses a single marked cycle at each rank and transfers the forest error directly to the cycle count.
The remaining obstacle is a sufficiently strong quantitative bound on that forest error or on the least edge frequency.

The fixed-edge-budget problem needs a different improvement as well.
Multiplying $2^r$ by a constant does not improve its exponential base when $r$ is proportional to the edge count.
No result here proves the proposed base $(2+2\sqrt2)^{1/5}$.
The degree-sequence paper supplying that target is Dvo\v r\'ak, Morrison, Noel, Norin, and Postle,
\emph{Bounding the number of cycles in a graph in terms of its degree sequence},
\url{https://arxiv.org/abs/1907.12091}.

\subsection{A boundary for matroid extensions}

Let $M=M^*(K_{r+1})$ be the cographic matroid of the complete graph.
Its binary cycle space is the graph's cut space and has dimension $r$.
Every nonempty cut $\delta(S)$ is a bond, since both induced sides are connected.
Thus every nonzero vector is a circuit, and $M$ has $2^r-1$ circuits.
The one-half theorem cannot extend to all binary matroids or even all regular matroids.
This counterexample was already identified in the supplied research list.

\subsection{Targets not resolved by this work}

General additive path stability is proved in Section~\ref{sec:general-additive-path-stability}: every terminal path has length at least $b-C_\varepsilon$.
It uses Fable's localization and layer argument and the subsequent short-cycle packing in bounded layers.
The exact four-regular bound and its equality cases are proved in the companion paper.
The leading random-host bound now holds throughout the full pancyclicity threshold range.
The sharp second-order term remains open, including the gap between coefficients $1$ and $199/192$ at the stated higher density and its behavior closer to the threshold.
Polynomial cycle-count error is proved for fixed proper minor-closed classes; the corresponding unrestricted bound and an exponential saving remain open.
No new result here proves monotonicity of minimum pancyclic excess or the proposed edge-pancyclic density bound.
The supplied list's distinction between cycle-space counts and obligations through every edge remains essential.


\section{Verification records}

The scripts beside this report use exact integer or rational arithmetic for finite graph identities and certificate inequalities.
They test the claimed finite statements; asymptotic probability and separator arguments are mathematical proofs.
\begin{itemize}
\item \nolinkurl{check_no_waste.py}: all marked cyclic edge subsets in the connected graph atlas through six vertices.
\item \nolinkurl{check_path_stability.py} and \nolinkurl{check_path_stability_phase2.py}: terminal series-parallel recurrences and explicit path-count obstructions.
\item \nolinkurl{check_path_minor_reduction.py}: all above-half terminal pairs and marked paths in the connected graph atlas through seven vertices.
\item \nolinkurl{check_link_stability.py}: exact small-cut moments and inverse inequalities.
\item \nolinkurl{check_random_host_encoding.py} and \nolinkurl{check_random_host_phase2.py}: explicit length encodings, pseudoforest inequalities, and theta-block constructions.
\item \nolinkurl{check_k4_random_host_blocks.py}: the integer bands, local density conditions, and spanning paths of the $K_4$ blocks.
\item \nolinkurl{check_random_host_wheels.py} and \nolinkurl{check_theta_random_host_partition.py}: wheel metrics, exact pair marginals, spanning paths, and the middle-interval theta partition.
\item \nolinkurl{check_random_host_rank3_barrier.py}: exact dual certificates for the limits of cycle, theta, and $K_4$ block constructions.
\item \nolinkurl{check_additive_layer_structure.py}: minimum-span cyclic windows and their greedy rank bound.
\item \nolinkurl{check_cubic_edge_frequency.py}: triangle and square identities and perfect-matching weighted cycle counts.
\item \nolinkurl{check_full_threshold_integration.py}: adaptive attachment positions, binary and digit lengths, Hamilton integration, and exact final excess.
\item Fable's \nolinkurl{fable_check_tools.py}: exact checks of the edge recursion, localization, product formula, cubic splitting, and layer boundaries.
\item \nolinkurl{verify_four_regular_theorem.py}: the complete finite part of the four-regular theorem, checked by two independent cycle counters.
\item \nolinkurl{four_regular_vertex_deletion_certificates.py}: exact analytic inequalities and induction constants for the four-regular theorem.
\end{itemize}

The accompanying JSON files contain completed results.
The larger four-regular experiments are retained as additional checks; they are not needed for the compact proof.
The source hashes and final PDF review status are recorded in \nolinkurl{verification-extended.json}.
\end{document}
